\documentclass[a4paper,12pt]{report}

% -------------------------------------------------------
% Pakete
% -------------------------------------------------------
\usepackage[T1]{fontenc}
\usepackage[utf8]{inputenc}
\usepackage[ngerman]{babel}
\usepackage{lmodern}
\usepackage{microtype}
\usepackage{geometry}
\geometry{left=3cm, right=2.5cm, top=3cm, bottom=3cm}

\usepackage{amsmath, amssymb, amsthm}
\usepackage{mathtools}
\usepackage{array}
\usepackage{booktabs}
\usepackage{longtable}
\usepackage{enumitem}
\usepackage{graphicx}
\usepackage{caption}
\usepackage{subcaption}
\usepackage{float}
\usepackage{tikz}
\usetikzlibrary{shapes.geometric, arrows.meta, positioning, calc, decorations.pathmorphing}

\usepackage{listings}
\usepackage{xcolor}

\lstdefinestyle{sqlstyle}{
  language=SQL,
  basicstyle=\ttfamily\small,
  keywordstyle=\color{blue!70!black}\bfseries,
  commentstyle=\color{gray}\itshape,
  stringstyle=\color{green!50!black},
  numbers=left,
  numberstyle=\tiny\color{gray},
  stepnumber=1,
  frame=single,
  rulecolor=\color{black!30},
  backgroundcolor=\color{gray!5},
  breaklines=true,
  showstringspaces=false,
  tabsize=2,
  captionpos=b,
  extendedchars=true,
  literate=%
    {Ö}{{\"O}}1 {Ä}{{\"A}}1 {Ü}{{\"U}}1 {ß}{{\ss}}1
    {ü}{{\"u}}1 {ä}{{\"a}}1 {ö}{{\"o}}1
}
\lstset{style=sqlstyle}

% Hyperref
\usepackage[hidelinks, colorlinks=false]{hyperref}
\usepackage{url}
\Urlmuskip=0mu plus 1mu

% Bibliographie
\usepackage[numbers,sort&compress]{natbib}

% Kopf-/Fusszeilen
\usepackage{fancyhdr}
\pagestyle{plain}
\renewcommand{\headrulewidth}{0pt}
\renewcommand{\footrulewidth}{0pt}

% Verhindere Überlauf über den rechten Rand
\sloppy
\setlength{\emergencystretch}{3em}

% Theoreme
\theoremstyle{definition}
\newtheorem{definition}{Definition}[chapter]
\newtheorem{theorem}{Satz}[chapter]
\newtheorem{lemma}[theorem]{Lemma}
\newtheorem{corollary}[theorem]{Korollar}
\newtheorem{proposition}[theorem]{Proposition}
\theoremstyle{remark}
\newtheorem{remark}{Bemerkung}[chapter]
\newtheorem{example}{Beispiel}[chapter]

\newcommand{\R}{\mathcal{R}}
\newcommand{\att}{\mathit{attr}}
\newcommand{\FD}{\mathit{FD}}
\newcommand{\MVD}{\mathit{MVD}}
\newcommand{\depth}{\mathit{depth}}
\newcommand{\breadth}{\mathit{breadth}}
\newcommand{\cost}{\mathit{cost}}
\newcommand{\NF}{\mathit{NF}}

\title{Relationale Datenbanksysteme \\ \large Eine formal fundierte Einführung mit
einem tiefenorientierten Entwurfsparadigma}
\author{Stephan Epp}
\date{\today}

% -------------------------------------------------------
\begin{document}

\begin{titlepage}
  \centering
  \vspace*{2cm}
  {\LARGE\bfseries Relationale Datenbanksysteme}\\[0.5em]
  {\Large Eine formal fundierte Einführung mit einem\\
  tiefenorientierten Entwurfsparadigma}
  \vspace{1.5cm}

  \rule{\textwidth}{0.4pt}\\[0.5em]
  \begin{tabular}{@{}r p{11cm}@{}}
    \textbf{Themengebiet:} & Datenbanktheorie, Schemaentwurf, SQL, PL/SQL, Anwendungsentwicklung \\[2pt]
    \textbf{Schwerpunkt:}  & Tiefenorientierter Entwurf -- Tiefenpfade vs.\ sternförmige Strukturen \\[2pt]
    \textbf{Autor:}        & Stephan Epp \\[2pt]
    \textbf{Datum:}        & \today \\
  \end{tabular}\\[0.5em]
  \rule{\textwidth}{0.4pt}
  \vfill
\end{titlepage}

\tableofcontents

\chapter*{Vorwort}
\addcontentsline{toc}{chapter}{Vorwort}

Das vorliegende Buch entwickelt eine vollständige, formal fundierte
Einführung in relationale Datenbanksysteme, deren Gliederung sich an dem
klassischen Lehrwerk \emph{Relationale Datenbanksysteme -- Eine praktische
Einführung} von Kleinschmidt und Rank orientiert. Im Unterschied zu rein
praxisorientierten Darstellungen wird in diesem Buch ein durchgängiges
formales Entwurfsparadigma -- der \emph{tiefenorientierte Datenbankentwurf}
-- in den Mittelpunkt gestellt: Es wird formal nachgewiesen, dass die
strukturelle Orientierung von Fremdschlüsselrelationen in die \emph{Tiefe}
(lineare bzw.\ baumartige Ketten) gegenüber der \emph{Breite}
(sternförmige Hub-Anordnungen) hinsichtlich Normalisierungsgrad,
funktionaler Abhängigkeitsstruktur, Anomalievermeidung,
Abfragekomplexität und Wartbarkeit überlegen ist. Die Argumentation
stützt sich auf die klassische relationale Algebra nach Codd, die
Normalformtheorie bis zur 4NF sowie graphentheoretische Resultate zu
Tiefen- und Breitensuche. Alle Kapitel des klassischen Kanons --
formale Modelle, SQL, Datenbankentwurf, Datenbankbetrieb, PL/SQL,
Programmiersprachenanbindung und WWW-Integration -- werden konsequent
unter diesem Paradigma neu beleuchtet und mit vollständigen formalen
Beweisen versehen.

\section*{Motivation und Entstehungshintergrund}

Die Idee zu diesem Buch entstand aus einer wiederkehrenden
Beobachtung: In zahlreichen produktiven Datenbankschemata -- gleich ob
in kleinen Anwendungen oder in gewachsenen Unternehmenssystemen --
findet sich immer wieder dasselbe strukturelle Muster: eine zentrale
\glqq Hub\grqq{}-Tabelle, an die eine Vielzahl weiterer Tabellen über
Fremdschlüssel angebunden ist, ohne dass die natürliche, oft tief
geschachtelte Abhängigkeitsstruktur der modellierten Domäne
abgebildet wird. Dieses Muster -- in diesem Buch als
\emph{Sternschema} formalisiert (Definition~\ref{def:stern}) -- wird
in der Praxis häufig aus pragmatischen Gründen gewählt: Es erscheint
auf den ersten Blick einfacher zu verstehen, da \glqq alles an einer
Stelle\grqq{} zu finden ist, und es vermeidet scheinbar die
\glqq Umständlichkeit\grqq{} mehrstufiger Verbundoperationen.

Die in diesem Buch entwickelte Theorie zeigt, dass dieser scheinbare
Vorteil bei genauerer Betrachtung trügerisch ist. Sobald die
modellierte Domäne -- wie praktisch immer -- transitive
Abhängigkeiten aufweist (ein \texttt{Auftrag} gehört zu einem
\texttt{Kunden}, der einer \texttt{Region} zugeordnet ist, die zu
einem \texttt{Land} gehört), erzwingt bereits die klassische
Normalisierungstheorie nach Codd eine \emph{tiefe}, kettenförmige
Zerlegung (Theorem~\ref{thm:transitivitaet}). Ein Sternschema, das
diese Tiefe künstlich in die Breite presst, verletzt zwangsläufig die
dritte Normalform und erkauft sich die scheinbare Einfachheit mit
Redundanz, Anomalieanfälligkeit und eingeschränkter Wartbarkeit.

Dieses Buch versteht sich daher nicht als bloße Wiederholung
etablierter Lehrbuchinhalte, sondern als ein \emph{durchgängiges
Argument}: Jedes Kapitel des klassischen Kanons -- von den formalen
Grundlagen über SQL, Datenbankentwurf und Datenbankbetrieb bis hin zu
PL/SQL, Hostsprachenanbindung und WWW-Integration -- wird daraufhin
untersucht, welche \emph{zusätzliche} Erkenntnis sich ergibt, wenn man
es konsequent durch die Brille des tiefenorientierten
Entwurfsparadigmas betrachtet. Die Antwort ist in jedem Fall
dieselbe: Die Tiefenstruktur des Schemas -- formal erfasst durch den
Fremdschlüsselgraphen $G_{\mathcal{S}}$ und die in
Definition~\ref{def:masse} eingeführten Maße $\depth$, $\breadth$ und
$\tau$ -- ist nicht nur ein Entwurfsdetail, sondern die zentrale
Determinante für nahezu jede praktisch relevante Eigenschaft des
Systems.

\section*{An wen sich dieses Buch richtet}

Dieses Buch wendet sich an Studierende der Wirtschaftsinformatik,
der Informatik und verwandter Studiengänge, die bereits eine erste
Einführung in relationale Datenbanken erhalten haben oder parallel
zu diesem Buch erhalten, ebenso wie an Praktikerinnen und Praktiker,
die Datenbankschemata entwerfen, warten oder refaktorisieren und ein
formales Vokabular suchen, um intuitiv empfundene
\glqq Code-Smells\grqq{} in Datenbankschemata präzise zu benennen und
zu begründen. Vorausgesetzt werden grundlegende Kenntnisse der
Mengenlehre und ein erster Kontakt mit SQL, wie er typischerweise in
einer einführenden Lehrveranstaltung vermittelt wird; die formalen
Beweise dieses Buches sind jedoch -- in Anlehnung an die in
\cite{Cormen2009} gepflegte Tradition -- so gestaltet, dass sie auch
ohne Vorkenntnisse in mathematischer Logik nachvollziehbar sind.

Da die Gliederung dieses Buches bewusst eng an das Lehrwerk von
Kleinschmidt und Rank~\cite{Kleinschmidt2002} angelehnt ist, eignet es
sich insbesondere als \emph{begleitende} oder \emph{vertiefende}
Lektüre zu Lehrveranstaltungen, die sich an jenem Werk orientieren:
Jedes Kapitel dieses Buches kann parallel zum entsprechenden Kapitel
des Vorbilds gelesen werden, wobei dieses Buch jeweils die zusätzliche,
formal-beweisende Perspektive des tiefenorientierten Entwurfs
beisteuert.

\section*{Wie dieses Buch zu lesen ist}

Die Kapitel~1 bis~8 folgen der klassischen Gliederung eines
einführenden Datenbanklehrbuchs und können -- mit Ausnahme der jeweils
am Ende eingestreuten Querverweise auf Kapitel~\ref{ch:tiefenentwurf}
-- weitgehend unabhängig voneinander gelesen werden, sofern die
Leserin oder der Leser bereits über grundlegende SQL-Kenntnisse
verfügt. Wer hingegen die zentrale formale Argumentation dieses
Buches \emph{zuerst} kennenlernen möchte, kann direkt mit
Kapitel~\ref{ch:tiefenentwurf} (\glqq Tiefenorientierter
Datenbankentwurf: Ein formales Gütekriterium\grqq{}) beginnen, das in
sich geschlossen die zentralen Definitionen, Sätze und Beweise dieses
Buches enthält, und anschließend in den übrigen Kapiteln verfolgen,
wie sich diese Theorie auf SQL (Kapitel~\ref{ch:sql}),
Datenbankentwurf (Kapitel~\ref{ch:entwurf}), Datenbankbetrieb
(Kapitel~\ref{ch:betrieb}), PL/SQL (Kapitel~\ref{ch:plsql}),
Hostsprachenanbindung (Kapitel~\ref{ch:hostsprachen}) und
WWW-Integration (Kapitel~\ref{ch:www}) auswirkt.

Sämtliche Definitionen, Sätze, Lemmata, Korollare, Propositionen und
Bemerkungen sind kapitelweise nummeriert und werden im gesamten Buch
konsistent referenziert; ein Symbolverzeichnis findet sich in
Anhang~A.2. Beweise sind durch das Symbol $\blacksquare$
abgeschlossen. SQL-Beispiele orientieren sich, sofern nicht anders
angegeben, an PostgreSQL, da dieses System sowohl die in
Kapitel~\ref{ch:sql} behandelten Standardkonstrukte als auch die in
Abschnitt~\ref{sec:distinct-on} behandelte Erweiterung
\texttt{DISTINCT ON} unterstützt.

\section*{Überblick über die Kapitel}

\textbf{Kapitel~\ref{ch:einleitung}} (\glqq Was sind
Datenbanksysteme?\grqq{}) führt die grundlegenden Begriffe ein --
DBMS, ACID-Eigenschaften, die Drei-Ebenen-Architektur nach
ANSI/SPARC -- und schließt mit einer Übersicht über Aufbau und
Leitidee des Buches.

\textbf{Kapitel~\ref{ch:modelle}} (\glqq Formale Modelle für
Datenbanken\grqq{}) behandelt das Relationenmodell und das
Entity-Relationship-Modell und führt mit
Definition~\ref{def:er-graph} den Fremdschlüsselgraphen
$G_{\mathcal{S}}$ ein, der im weiteren Verlauf des Buches als
zentrales formales Objekt dient.

\textbf{Kapitel~\ref{ch:sql}} (\glqq Datenbankoperationen mit
SQL\grqq{}) behandelt SELECT-Anfragen, Views und
Datenmodifikationen und enthält mit
Abschnitt~\ref{sec:distinct-on} eine ausführliche, formal untermauerte
Behandlung des PostgreSQL-spezifischen Sprachmittels
\texttt{DISTINCT ON} und seiner Beziehung zum
Gruppenrepräsentantenproblem.

\textbf{Kapitel~\ref{ch:entwurf}} (\glqq Datenbank-Entwurf\grqq{})
behandelt Anomalien, Normalformen und Dekomposition und enthält mit
Abschnitt~\ref{sec:minimalcover} einen vollständigen
Synthesealgorithmus zur Schemaminimierung samt Pseudocode,
Laufzeitanalyse und Korrektheitsbeweis.

\textbf{Kapitel~\ref{ch:betrieb}} (\glqq Datenbankbetrieb\grqq{})
behandelt Transaktionen, Mehrbenutzerbetrieb, Zugriffsrechte und
Systemtabellen und zeigt, wie sich die Tiefenpfadstruktur auf
Sperrgranularität und Rechtevergabe auswirkt.

\textbf{Kapitel~\ref{ch:plsql}} (\glqq PL/SQL\grqq{}) behandelt
prozedurale SQL-Erweiterungen, Cursor, Fehlerbehandlung, Prozeduren,
Funktionen und Trigger.

\textbf{Kapitel~\ref{ch:hostsprachen}} (\glqq Datenbankzugriff in
höheren Programmiersprachen\grqq{}) behandelt die
Impedanzdiskrepanz, Embedded SQL, ODBC/JDBC, die Perl-DBI-Schnittstelle
und objektrelationales Mapping und enthält mit
Theorem~\ref{thm:repository-zerlegung} einen formalen Nachweis, dass
ein Tiefenpfadschema eine eindeutige, lokal konsistente Zerlegung des
Anwendungscodes in Repository-Module induziert.

\textbf{Kapitel~\ref{ch:www}} (\glqq WWW-Integration von
Datenbanken\grqq{}) behandelt die Architektur webbasierter
Datenbankanwendungen, CGI-Skripte, Web-Formulare und
Mehrbenutzerbetrieb im Web-Kontext.

\textbf{Kapitel~\ref{ch:tiefenentwurf}} (\glqq Tiefenorientierter
Datenbankentwurf: Ein formales Gütekriterium\grqq{}) bildet das
formale Herzstück dieses Buches und enthält die vollständige,
in sich geschlossene Argumentation für das Prinzip
\glqq Tiefe vor Breite\grqq{}, einschließlich der
graphentheoretischen Fundierung mittels Tiefen- und Breitensuche und
eines durchgängigen Praxisbeispiels.

\textbf{Kapitel~\ref{ch:ausblick}} (\glqq Fazit und Ausblick\grqq{})
fasst die Ergebnisse zusammen und entwickelt ein umfangreiches
Forschungsprogramm, das von automatisiertem Schema-Tooling über
temporale und verteilte Datenbanken, Sicherheits- und
Datenschutzaspekte, formale Verifikation und maschinelles Lernen bis
hin zu hybriden Graph-relationalen Systemen reicht.

\section*{Danksagung}

Dieses Buch wäre ohne die in Abschnitt~1.4 genannte strukturelle
Anlehnung an das Lehrwerk von Kleinschmidt und
Rank~\cite{Kleinschmidt2002} nicht in dieser Form entstanden; ihm
gebührt der Dank für eine Gliederung, die sich seit nunmehr über zwei
Jahrzehnten in der Lehre bewährt hat und die sich -- wie dieses Buch
zeigt -- auch als Gerüst für eine deutlich formalere, beweisorientierte
Darstellung eignet. Für verbleibende Fehler und Unzulänglichkeiten
trägt allein der Autor die Verantwortung.

\bigskip
\noindent Bielefeld, \today \hfill Stephan Epp

\chapter{Was sind Datenbanksysteme?}
\label{ch:einleitung}

\section{Unterschiede von DBMS zu klassischen Dateisystemen}

Klassische dateiorientierte Anwendungssysteme verwalten Daten in
applikationsspezifischen Dateien, deren Struktur fest in der Anwendung
kodiert ist. Dies führt zu einer Reihe struktureller Probleme, die in
einem Datenbankmanagementsystem (DBMS) systematisch vermieden werden:

\begin{itemize}
  \item \textbf{Datenredundanz:} Dieselbe Information wird in mehreren
    Dateien dupliziert, was Speicherplatz verschwendet und
    Inkonsistenzen begünstigt.
  \item \textbf{Programm-Daten-Abhängigkeit:} Änderungen am
    Datenformat erfordern Anpassungen aller zugreifenden Programme.
  \item \textbf{Fehlende Mehrbenutzerfähigkeit:} Der gleichzeitige
    Zugriff mehrerer Anwendungen auf dieselben Daten ist nicht oder nur
    eingeschränkt möglich.
  \item \textbf{Fehlende Konsistenzsicherung:} Es existiert kein
    zentraler Mechanismus, der die Einhaltung von
    Integritätsbedingungen über Dateigrenzen hinweg garantiert.
\end{itemize}

Ein DBMS hingegen stellt eine zentrale Verwaltungskomponente bereit, die
Daten unabhängig von den zugreifenden Anwendungsprogrammen organisiert,
speichert und konsistent hält. Die Daten werden dabei in einem
gemeinsamen, formal beschriebenen Schema abgelegt, auf das mehrere
Anwendungen über eine deklarative Anfragesprache zugreifen.

\begin{definition}[Datenbankmanagementsystem]
\label{def:dbms}
Ein Datenbankmanagementsystem (DBMS) ist ein Softwaresystem, das die
Definition, Speicherung, Manipulation und Verwaltung einer Datenbank
$\mathcal{D}$ derart bereitstellt, dass
\begin{enumerate}
  \item die logische Struktur der Daten (das \emph{Schema}) von ihrer
    physischen Speicherung getrennt ist (Datenunabhängigkeit),
  \item mehrere Anwendungsprogramme nebenläufig und konsistent auf
    $\mathcal{D}$ zugreifen können,
  \item die Einhaltung definierter Integritätsbedingungen
    $\mathcal{I}$ für jeden Zustand von $\mathcal{D}$ garantiert wird,
  \item ein Wiederanlauf nach Fehlern in einen konsistenten Zustand
    möglich ist (Persistenz und Recovery).
\end{enumerate}
\end{definition}

\section{Forderungen an ein DBMS}

Aus Definition~\ref{def:dbms} lassen sich die klassischen Forderungen an
ein DBMS ableiten, die in der Literatur häufig unter dem Akronym
\emph{ACID} für Transaktionseigenschaften zusammengefasst werden
(siehe Kapitel~\ref{ch:betrieb}):

\begin{itemize}
  \item \textbf{Atomarität (Atomicity):} Eine Transaktion wird entweder
    vollständig oder gar nicht ausgeführt.
  \item \textbf{Konsistenz (Consistency):} Jede Transaktion überführt
    die Datenbank von einem konsistenten in einen konsistenten Zustand.
  \item \textbf{Isolation:} Nebenläufig ausgeführte Transaktionen
    dürfen sich nicht gegenseitig beeinflussen, als ob sie seriell
    ausgeführt würden.
  \item \textbf{Dauerhaftigkeit (Durability):} Die Effekte einer
    erfolgreich abgeschlossenen Transaktion bleiben auch im Fehlerfall
    erhalten.
\end{itemize}

Darüber hinaus fordert man von einem modernen DBMS Effizienz bei
Anfragen über große Datenmengen, Skalierbarkeit, ein feingranulares
Zugriffsrechtesystem sowie Erweiterbarkeit des Schemas ohne
Unterbrechung des laufenden Betriebs.

\section{Abstraktionsebenen im DBMS}

Die Architektur eines DBMS wird klassisch in drei Abstraktionsebenen
gegliedert (ANSI/SPARC-Architektur):

\begin{enumerate}
  \item \textbf{Externe Ebene:} Sichten (Views) auf die Daten, die für
    einzelne Anwendungen oder Benutzergruppen zugeschnitten sind.
  \item \textbf{Konzeptuelle Ebene:} Das logische Gesamtschema der
    Datenbank, unabhängig von physischer Speicherung und individuellen
    Sichten.
  \item \textbf{Interne Ebene:} Die physische Speicherorganisation,
    Indexstrukturen, Zugriffspfade.
\end{enumerate}

\begin{figure}[H]
\centering
\begin{tikzpicture}[
  layer/.style={draw, rectangle, rounded corners, minimum width=7cm,
    minimum height=1.2cm, fill=blue!5},
  node distance=0.6cm
]
\node[layer] (ext) {Externe Ebene (Views)};
\node[layer, below=of ext] (kon) {Konzeptuelle Ebene (Gesamtschema)};
\node[layer, below=of kon] (int) {Interne Ebene (Speicherstruktur, Indizes)};
\draw[<->] (ext) -- (kon);
\draw[<->] (kon) -- (int);
\end{tikzpicture}
\caption{Drei-Ebenen-Architektur eines DBMS (ANSI/SPARC)}
\label{fig:ansi-sparc}
\end{figure}

Diese strikte Trennung garantiert \emph{logische} und \emph{physische
Datenunabhängigkeit}: Änderungen an der internen Speicherorganisation
bleiben für die externe Ebene unsichtbar, und Änderungen am
konzeptuellen Schema -- sofern sie additiv erfolgen -- beeinträchtigen
bestehende externe Sichten nicht. Genau diese Eigenschaft wird in
Kapitel~\ref{ch:tiefenentwurf} zum formalen Ausgangspunkt für den
Nachweis genommen, dass tiefenorientierte Schemata eine überlegene
Schemaevolution ermöglichen.

\section{Aufbau und Leitidee dieses Buches}

Die Gliederung dieses Buches folgt bewusst der Struktur des klassischen
Lehrwerks von Kleinschmidt und Rank~\cite{Kleinschmidt2002}: Kapitel~2
behandelt die formalen Modelle (Relationenmodell,
Entity-Relationship-Modell), Kapitel~3 die Datenbankoperationen mit SQL,
Kapitel~4 den Datenbankentwurf einschließlich Normalisierung,
Kapitel~5 den Datenbankbetrieb (Transaktionen, Mehrbenutzerbetrieb,
Zugriffsrechte), Kapitel~6 PL/SQL, Kapitel~7 die Anbindung an höhere
Programmiersprachen und Kapitel~8 die WWW-Integration von Datenbanken.

Im Unterschied zum Vorbild wird jedoch in jedem dieser Kapitel ein
zusätzlicher, durchgängiger formaler Strang verfolgt: das
\emph{tiefenorientierte Entwurfsparadigma}. Dieses Paradigma -- in
Kapitel~\ref{ch:tiefenentwurf} vollständig und in sich geschlossen
hergeleitet -- besagt, dass relationale Schemata, deren
Fremdschlüsselstruktur als gerichteter, azyklischer \emph{Tiefenpfad}
(im Sinne eines DFS-Walds) organisiert ist, sternförmigen
(BFS-artigen) Schemata in allen klassischen Gütekriterien überlegen
sind. Die Kapitel~2 bis~8 verweisen an den jeweils relevanten Stellen
auf dieses Paradigma und illustrieren es konkret an SQL-Beispielen,
Normalformen, Transaktionsanalysen und WWW-Architekturen.

\chapter{Formale Modelle für Datenbanken}
\label{ch:modelle}

\section{Das Relationenmodell}
\label{sec:relmodell}

Das Relationenmodell, eingeführt von Codd~\cite{Codd1970}, bildet die
mathematische Grundlage relationaler Datenbanksysteme. Es basiert auf
dem Begriff der \emph{Relation} im Sinne der Mengenlehre.

\begin{definition}[Relation]
\label{def:relation}
Sei $D_1, D_2, \ldots, D_n$ eine Folge von \emph{Domänen} (Wertebereichen).
Eine Relation $\R$ über dem Schema $(A_1:D_1, \ldots, A_n:D_n)$ ist eine
endliche Teilmenge des kartesischen Produkts
\[
  \R \subseteq D_1 \times D_2 \times \cdots \times D_n .
\]
Die Elemente $A_1,\ldots,A_n$ heißen \emph{Attribute}, ein Element
$t \in \R$ heißt \emph{Tupel}.
\end{definition}

\begin{definition}[Relationenschema, Datenbankschema]
\label{def:relschema}
Ein Relationenschema ist ein Paar $R = (\att(R), \Sigma_R)$, bestehend
aus einer endlichen Menge von Attributen $\att(R) = \{A_1,\ldots,A_n\}$
und einer Menge von Integritätsbedingungen $\Sigma_R$ (insbesondere
funktionalen Abhängigkeiten, vgl.\ Abschnitt~\ref{sec:normalformen}). Ein
Datenbankschema $\mathcal{S} = \{R_1,\ldots,R_m\}$ ist eine endliche
Menge von Relationenschemata.
\end{definition}

\begin{definition}[Schlüssel]
\label{def:schluessel}
Eine Menge $K \subseteq \att(R)$ heißt \emph{Schlüssel} von $R$, wenn
\begin{enumerate}
  \item $K$ jedes Tupel von $R$ eindeutig identifiziert (Eindeutigkeit), und
  \item keine echte Teilmenge $K' \subsetneq K$ diese Eigenschaft besitzt
    (Minimalität).
\end{enumerate}
Eine Menge $A \subseteq \att(R')$ in einem anderen Schema $R'$ heißt
\emph{Fremdschlüssel} bezüglich $R$, wenn $A$ und der Primärschlüssel von
$R$ über derselben Domäne definiert sind und jeder in $A$ vorkommende
Wert entweder \texttt{NULL} ist oder als Primärschlüsselwert in $R$
existiert (referentielle Integrität).
\end{definition}

\section{Das Entity-Relationship-Modell}
\label{sec:ermodell}

Das von Chen eingeführte Entity-Relationship-Modell (ER-Modell) dient
der konzeptuellen Modellierung auf einer höheren Abstraktionsebene als
das Relationenmodell. Zentrale Konstrukte sind:

\begin{itemize}
  \item \textbf{Entitätstypen} repräsentieren Klassen von Objekten der
    Anwendungswelt (z.\,B. \textit{Kunde}, \textit{Auftrag}).
  \item \textbf{Attribute} beschreiben Eigenschaften von Entitäten oder
    Beziehungen.
  \item \textbf{Beziehungstypen} (Relationships) modellieren
    Assoziationen zwischen Entitätstypen, charakterisiert durch ihre
    \emph{Kardinalität} ($1{:}1$, $1{:}n$, $n{:}m$).
\end{itemize}

\begin{definition}[ER-Schema als gerichteter Graph]
\label{def:er-graph}
Ein ER-Schema lässt sich als gerichteter Graph $G_{ER} = (V,E)$ auffassen,
wobei $V$ die Menge der Entitätstypen ist und
$(E_i, E_j) \in E$ genau dann gilt, wenn ein Beziehungstyp zwischen $E_i$
und $E_j$ mit Kardinalität $1{:}n$ oder $n{:}1$ existiert, gerichtet von
der Seite mit Kardinalität $n$ (der „vielen“ Seite, die den
Fremdschlüssel trägt) zur Seite mit Kardinalität $1$.
\end{definition}

Definition~\ref{def:er-graph} ist der zentrale Bezugspunkt für das in
Kapitel~\ref{ch:tiefenentwurf} entwickelte tiefenorientierte
Entwurfsparadigma: Die Frage, ob ein Datenbankschema „in die Tiefe“ oder
„in die Breite“ orientiert ist, wird formal als eine Eigenschaft des
Graphen $G_{ER}$ -- nämlich seiner Tiefe $\depth(G_{ER})$ im Sinne von
Definition~\ref{def:masse} -- charakterisiert.

\section{Vom Entity-Relationship- zum Relationenmodell}
\label{sec:er-zu-rel}

Die Transformation eines ER-Schemas in ein Relationenschema folgt
festen Regeln:

\begin{enumerate}
  \item Jeder Entitätstyp wird zu einem Relationenschema mit den
    Attributen der Entität und einem Primärschlüssel.
  \item Ein Beziehungstyp mit Kardinalität $1{:}n$ wird durch Aufnahme
    des Primärschlüssels der „1“-Seite als Fremdschlüssel in das
    Relationenschema der „$n$“-Seite realisiert.
  \item Ein Beziehungstyp mit Kardinalität $n{:}m$ erfordert ein eigenes
    Relationenschema, dessen Schlüssel sich aus den Primärschlüsseln
    beider beteiligter Entitätstypen zusammensetzt.
\end{enumerate}

\begin{proposition}
\label{prop:er-rel-graph}
Sei $G_{ER}=(V,E)$ ein ER-Schema gemäß Definition~\ref{def:er-graph} und
$\mathcal{S}$ das nach obigen Regeln daraus abgeleitete relationale
Schema. Dann existiert für jede Kante $(R_i,R_j)\in E$ ein
Fremdschlüssel in $\att(R_i)$, der auf den Primärschlüssel von $R_j$
verweist, sofern die Beziehung nicht $n{:}m$ ist (in diesem Fall wird
die Kante in zwei Kanten über ein zusätzliches Verbindungsschema
zerlegt). Insbesondere bleibt die Graphstruktur von $G_{ER}$ -- bis auf
diese Zerlegung von $n{:}m$-Kanten -- im Fremdschlüsselgraphen
$G_{\mathcal{S}}$ des relationalen Schemas erhalten.
\end{proposition}

\begin{proof}
Für $1{:}n$-Beziehungen folgt die Behauptung unmittelbar aus
Transformationsregel~2: Der Primärschlüssel von $R_j$ (der „1“-Seite)
wird als Attribut in $\att(R_i)$ aufgenommen und erfüllt nach
Definition~\ref{def:schluessel} die Bedingungen eines Fremdschlüssels,
da seine Werte aus dem Wertebereich des Primärschlüssels von $R_j$
stammen und referentielle Integrität gefordert wird. Für $n{:}m$-Beziehungen
erzeugt Regel~3 ein neues Schema $R_{ij}$ mit Fremdschlüsseln auf
\emph{beide} Seiten $R_i$ und $R_j$; die ursprüngliche Kante
$(R_i,R_j)$ wird somit durch den Pfad $R_i \leftarrow R_{ij} \rightarrow R_j$
ersetzt, was strukturell einer Aufspaltung in zwei gerichtete Kanten mit
gemeinsamem Zielknoten $R_{ij}$ entspricht. In beiden Fällen bleibt die
Adjazenzstruktur -- modulo dieser lokalen Zerlegung -- erhalten, womit die
Behauptung gezeigt ist.
\end{proof}

Diese Proposition ist von zentraler Bedeutung, da sie es erlaubt, alle
graphentheoretischen Aussagen über $G_{ER}$ (insbesondere die in
Kapitel~\ref{ch:tiefenentwurf} bewiesenen Sätze über Tiefe und Breite)
direkt auf den Fremdschlüsselgraphen des relationalen Schemas zu
übertragen.

\section{Wichtige Operationen auf Relationen (Tabellen)}
\label{sec:relalg}

Die relationale Algebra definiert eine Menge von Operationen, die
Relationen auf neue Relationen abbilden und damit abgeschlossen
(„relational vollständig“) sind:

\begin{itemize}
  \item \textbf{Selektion} $\sigma_{\varphi}(\R)$: Auswahl der Tupel aus
    $\R$, die die Bedingung $\varphi$ erfüllen.
  \item \textbf{Projektion} $\pi_{A_1,\ldots,A_k}(\R)$: Beschränkung auf
    die angegebenen Attribute, mit anschließender Duplikateliminierung.
  \item \textbf{Vereinigung, Differenz, Durchschnitt:} mengentheoretische
    Operationen auf Relationen mit gleichem Schema.
  \item \textbf{Kartesisches Produkt} $\R \times \mathcal{S}$: Kombination
    aller Tupelpaare.
  \item \textbf{Verbund (Join)} $\R \Join_{\varphi} \mathcal{S}$: Selektion
    auf dem kartesischen Produkt gemäß einer Verbundbedingung $\varphi$,
    typischerweise einer Gleichheit zwischen Fremdschlüssel- und
    Primärschlüsselattributen (\emph{Equi-Join}).
  \item \textbf{Umbenennung} $\rho_{A \to B}(\R)$: Umbenennung von
    Attributen.
\end{itemize}

\begin{definition}[Natürlicher Verbund über Fremdschlüsselkante]
\label{def:fk-join}
Sei $(R_i,R_j)$ eine Kante im Fremdschlüsselgraphen $G_{\mathcal{S}}$
gemäß Proposition~\ref{prop:er-rel-graph}, mit Fremdschlüssel
$F \subseteq \att(R_i)$, der auf den Primärschlüssel $K \subseteq
\att(R_j)$ verweist. Der \emph{kanonische Verbund} entlang dieser Kante
ist
\[
  R_i \Join_{F=K} R_j .
\]
\end{definition}

Definition~\ref{def:fk-join} liefert die formale Grundlage für die in
Abschnitt~\ref{sec:abfragen} und Kapitel~\ref{ch:tiefenentwurf}
durchgeführte Analyse der Join-Komplexität entlang von Tiefenpfaden im
Vergleich zu Sternschemata.

\section{Zusatzforderungen für relationale DBMS}
\label{sec:codd-rules}

Codd formulierte zwölf Regeln, die ein DBMS erfüllen muss, um als
„vollständig relational“ zu gelten~\cite{Codd1970,Codd1974}. Für die
vorliegende Arbeit sind insbesondere die folgenden Forderungen
relevant:

\begin{itemize}
  \item \textbf{Informationsregel:} Sämtliche Information -- inklusive
    Metadaten -- wird ausschließlich in Form von Werten in Tabellen
    dargestellt.
  \item \textbf{Garantierter Zugriff:} Jeder Datenwert ist eindeutig
    über Tabellenname, Primärschlüsselwert und Attributname adressierbar.
  \item \textbf{Systematische Behandlung von Nullwerten:} \texttt{NULL}
    repräsentiert fehlende oder unanwendbare Information unabhängig vom
    Datentyp.
  \item \textbf{Strukturunabhängigkeit:} Die Anfragesprache darf nicht
    von der physischen Speicherstruktur abhängen.
  \item \textbf{Integritätsunabhängigkeit:} Integritätsbedingungen
    (insbesondere referentielle Integrität gemäß
    Definition~\ref{def:schluessel}) werden im Schema selbst definiert
    und vom DBMS automatisch durchgesetzt, unabhängig von
    Anwendungsprogrammen.
\end{itemize}

Die letztgenannte Forderung -- Integritätsunabhängigkeit -- ist von
besonderer Bedeutung für die in Kapitel~\ref{ch:tiefenentwurf}
diskutierte Wartbarkeit: Da Fremdschlüsselbeziehungen im Schema selbst
deklariert werden, lässt sich die Tiefe eines Schemas
$\depth(G_{\mathcal{S}})$ direkt aus dem Datenbankkatalog (vgl.\
Abschnitt~\ref{sec:systemtabellen}) ablesen, ohne den
Anwendungsquellcode analysieren zu müssen.

\chapter{Datenbankoperationen mit SQL}
\label{ch:sql}

\section{Anforderungen an einen Datenbank-Server}

Ein Datenbank-Server muss nebenläufige Verbindungen mehrerer Clients
verwalten, Anfragen optimieren und ausführen sowie Transaktionen gemäß
den ACID-Eigenschaften (vgl.\ Kapitel~\ref{ch:einleitung}) garantieren.
Die in den folgenden Abschnitten vorgestellten SQL-Konstrukte werden in
diesem Buch konsequent anhand des in Kapitel~\ref{ch:tiefenentwurf}
entwickelten Tiefenpfadschemas illustriert (vgl.\
Abschnitt~\ref{sec:beispiel}).

\section{SQL-Standards}
\label{sec:sql-standards}

SQL (Structured Query Language) wurde mit ANSI SQL-86 erstmals
standardisiert und seitdem mehrfach erweitert: SQL-92 (umfangreiche
Erweiterung der Datentypen und Integritätsbedingungen), SQL:1999
(rekursive Anfragen, objektrelationale Erweiterungen, Trigger), sowie
nachfolgende Standards (SQL:2003, SQL:2008, SQL:2011, SQL:2016) mit
Erweiterungen für XML, Window-Funktionen, temporale Tabellen und
JSON-Unterstützung. Die in diesem Buch verwendeten Konstrukte
orientieren sich am SQL:1999-Standard mit punktuellen Hinweisen auf
spätere Erweiterungen.

\section{Sprachstruktur von SQL}

SQL gliedert sich in mehrere Teilsprachen:

\begin{itemize}
  \item \textbf{DDL} (Data Definition Language): \texttt{CREATE},
    \texttt{ALTER}, \texttt{DROP} -- Definition von Schemaobjekten.
  \item \textbf{DML} (Data Manipulation Language): \texttt{SELECT},
    \texttt{INSERT}, \texttt{UPDATE}, \texttt{DELETE}.
  \item \textbf{DCL} (Data Control Language): \texttt{GRANT},
    \texttt{REVOKE} -- Verwaltung von Zugriffsrechten.
  \item \textbf{TCL} (Transaction Control Language): \texttt{COMMIT},
    \texttt{ROLLBACK}, \texttt{SAVEPOINT}.
\end{itemize}

\section{Abrufen von Informationen mit SELECT}
\label{sec:select}

Die allgemeine Form einer \texttt{SELECT}-Anweisung lautet:

\begin{lstlisting}
SELECT [DISTINCT] Ausdrucksliste
FROM   Tabellenliste
[WHERE  Bedingung]
[GROUP BY Attributliste]
[HAVING Gruppenbedingung]
[ORDER BY Attributliste]
\end{lstlisting}

Semantisch entspricht eine \texttt{SELECT}-Anweisung einer Komposition
relationaler Algebra-Operatoren gemäß Abschnitt~\ref{sec:relalg}:

\[
\pi_{\text{Ausdrucksliste}}\Big(
  \sigma_{\text{Bedingung}}\big(
    R_1 \times R_2 \times \cdots \times R_k
  \big)
\Big)
\]

mit anschließender optionaler Gruppierung, Aggregation und Sortierung.

\begin{example}[Verbund über eine Fremdschlüsselkante]
\label{ex:join-tiefe}
Sei $\textsc{Auftrag}$ ein Relationenschema mit Fremdschlüssel
\texttt{kunden\_id} auf $\textsc{Kunde}$ (vgl.\ Definition~\ref{def:fk-join}).
Die Anfrage
\begin{lstlisting}
SELECT a.auftrag_id, k.name, a.datum
FROM   Auftrag a
JOIN   Kunde k ON a.kunden_id = k.kunden_id
WHERE  a.datum >= DATE '2026-01-01'
\end{lstlisting}
realisiert den kanonischen Verbund $\textsc{Auftrag}
\Join_{\texttt{kunden\_id}=\texttt{kunden\_id}} \textsc{Kunde}$ mit
anschließender Selektion und Projektion. Im Tiefenpfadschema aus
Abschnitt~\ref{sec:beispiel} entspricht dieser Verbund genau
einer Kante des DFS-Walds.
\end{example}

Mehrstufige Verbunde entlang eines Tiefenpfads
$R_1 \rightarrow R_2 \rightarrow \cdots \rightarrow R_k$ werden durch
Verkettung mehrerer \texttt{JOIN}-Klauseln realisiert; die
Komplexitätsanalyse hierzu findet sich in
Abschnitt~\ref{sec:abfragen}.

\section{Das Sprachmittel \texttt{DISTINCT ON}: Gruppenrepräsentanten entlang des Tiefenpfads}
\label{sec:distinct-on}

Im Folgenden wird ein Sprachmittel vorgestellt, das zwar nicht Teil des
SQL-Standards ist, aber für die in diesem Buch entwickelte
Tiefenpfadarchitektur (Kapitel~\ref{ch:tiefenentwurf}) von besonderer
Bedeutung ist: die PostgreSQL-spezifische Erweiterung
\texttt{SELECT DISTINCT ON (\ldots)}. Dieses Sprachmittel löst auf
deklarative, mengenorientierte Weise eine Aufgabe, die in der Praxis
außerordentlich häufig auftritt: das Auffinden eines
\emph{Gruppenrepräsentanten} -- z.\,B.\ des jeweils neuesten Tupels pro
Fremdschlüsselwert.

\subsection{Motivation: Standard-\texttt{DISTINCT} versus Gruppenrepräsentanten}

Der SQL-Standard kennt nur das einfache Schlüsselwort
\texttt{DISTINCT}, das gemäß Abschnitt~\ref{sec:select} als
Duplikateliminierung nach der Projektion $\pi$ wirkt:

\begin{definition}[Standard-\texttt{DISTINCT}]
\label{def:distinct-standard}
Sei $\R$ eine Relation und $A_1,\ldots,A_k\in\att(\R)$. Die Anweisung
\texttt{SELECT DISTINCT $A_1,\ldots,A_k$ FROM $\R$} liefert das
Ergebnis
\[
  \pi_{A_1,\ldots,A_k}(\R)
\]
im Sinne der relationalen Algebra (Abschnitt~\ref{sec:relalg}), d.\,h.\
die Menge \emph{aller verschiedenen Wertekombinationen}
$(t[A_1],\ldots,t[A_k])$, die in $\R$ auftreten -- ohne Information
darüber, \emph{welches} Tupel von $\R$ diese Kombination geliefert
hat, falls mehrere Tupel dieselbe Kombination aufweisen.
\end{definition}

Häufig stellt sich jedoch eine andere, semantisch verwandte, aber
algebraisch verschiedene Aufgabe: Gegeben eine Gruppierung von $\R$
nach einer Attributmenge $G\subseteq\att(\R)$ -- etwa $G=\{$
\texttt{kunden\_id}$\}$ --, soll für jede Gruppe \emph{ein vollständiges
Tupel} zurückgegeben werden, nämlich dasjenige, das bezüglich einer
Ordnung $\prec$ auf $\att(\R)$ \emph{minimal} bzw.\ \emph{maximal}
ist. Diese Aufgabe wird im Folgenden formal als
\emph{Gruppenrepräsentantenproblem} bezeichnet.

\begin{definition}[Gruppenrepräsentantenproblem]
\label{def:gruppenrepraesentant}
Sei $\R$ eine Relation, $G\subseteq\att(\R)$ eine
\emph{Gruppierungsattributmenge} und $\prec$ eine totale
Ordnung auf den Tupeln von $\R$, induziert durch eine
Attributfolge $(B_1,\ldots,B_m)\subseteq\att(\R)$ mit zugehörigen
Sortierrichtungen (\texttt{ASC}/\texttt{DESC}). Für jeden in $\R$
auftretenden Wert $g\in\pi_G(\R)$ sei
\[
  \R_g := \sigma_{(A_1,\ldots,A_{|G|})=g}(\R)
\]
die Gruppe aller Tupel mit Gruppierungswert $g$. Das
Gruppenrepräsentantenproblem besteht darin, für jedes $g\in\pi_G(\R)$
genau das (bezüglich $\prec$ kleinste) Tupel
\[
  t_g := \min_{\prec} \R_g
\]
zu bestimmen und die Ergebnisrelation
$\{\,t_g \mid g\in\pi_G(\R)\,\}$ zurückzugeben.
\end{definition}

\begin{remark}
Definition~\ref{def:gruppenrepraesentant} ist \emph{nicht} durch
Standard-\texttt{DISTINCT} (Definition~\ref{def:distinct-standard})
ausdrückbar, da \texttt{DISTINCT} nur auf den
\emph{projizierten} Attributen operiert und keine Möglichkeit bietet,
„das übrige Tupel mitzuführen“, das zu einer bestimmten
Wertekombination gehört.
\end{remark}

\subsection{Das Standard-SQL-Vorgehen: Fensterfunktion mit \texttt{ROW\_NUMBER()}}
\label{sec:row-number}

Innerhalb des SQL-Standards (SQL:2003 ff., vgl.\
Abschnitt~\ref{sec:sql-standards}) wird das
Gruppenrepräsentantenproblem über eine \emph{Fensterfunktion}
(Window-Funktion) gelöst, typischerweise \texttt{ROW\_NUMBER()} in
Kombination mit einer \texttt{PARTITION BY}-Klausel:

\begin{lstlisting}
WITH RankedOrders AS (
  SELECT
    kunden_id, auftrag_id, datum, gesamtbetrag,
    ROW_NUMBER() OVER (
      PARTITION BY kunden_id
      ORDER BY datum DESC
    ) AS rn
  FROM Auftrag
)
SELECT kunden_id, auftrag_id, datum, gesamtbetrag
FROM RankedOrders
WHERE rn = 1;
\end{lstlisting}

Semantisch geschieht hier Folgendes: Die \texttt{PARTITION
BY}-Klausel zerlegt $\R=\textsc{Auftrag}$ in die Gruppen $\R_g$ aus
Definition~\ref{def:gruppenrepraesentant} mit
$G=\{\texttt{kunden\_id}\}$. Innerhalb jeder Gruppe nummeriert
\texttt{ROW\_NUMBER()} die Tupel gemäß der durch \texttt{ORDER BY
datum DESC} induzierten Ordnung $\prec$ durch, beginnend bei~$1$ für
das bezüglich $\prec$ kleinste (hier: das Tupel mit dem größten
\texttt{datum}, da \texttt{DESC}) Tupel. Der äußere Filter
\texttt{WHERE rn = 1} selektiert genau dieses minimale Tupel pro
Gruppe -- also exakt $t_g$ aus
Definition~\ref{def:gruppenrepraesentant}.

\begin{remark}[Konzeptionelle Zweistufigkeit]
\label{rem:zweistufig}
Das \texttt{ROW\_NUMBER()}-Vorgehen ist konzeptionell
\emph{zweistufig}: In einer ersten, materialisierten
Zwischenrelation (\texttt{RankedOrders}, realisiert über eine
Common Table Expression bzw.\ einen Subquery) wird \emph{jedem}
Tupel von $\R$ eine Rangnummer zugewiesen -- also $|\R|$ Tupel
erzeugt --, bevor in einer zweiten Stufe die Filterung
\texttt{rn = 1} auf $|\pi_G(\R)|$ Tupel reduziert. Dieser
Zwischenschritt ist algebraisch nicht notwendig, da das Endergebnis
nur $|\pi_G(\R)|\leq|\R|$ Tupel umfasst; er ist jedoch eine direkte
Konsequenz davon, dass Fensterfunktionen im SQL-Standard
\emph{zeilenweise} (pro Eingabetupel ein Ausgabewert) definiert sind
und nicht direkt als gruppierende Aggregation mit „Tupel-Rückgabe“
formuliert werden können.
\end{remark}

\subsection{\texttt{DISTINCT ON}: Direkte deklarative Lösung in PostgreSQL}
\label{sec:distinct-on-postgres}

PostgreSQL bietet mit \texttt{SELECT DISTINCT ON (\ldots)} eine
direkte, einstufige Lösung des Gruppenrepräsentantenproblems:

\begin{lstlisting}
SELECT DISTINCT ON (kunden_id)
       kunden_id, auftrag_id, datum, gesamtbetrag
FROM   Auftrag
ORDER BY kunden_id, datum DESC;
\end{lstlisting}

\begin{definition}[Semantik von \texttt{DISTINCT ON}]
\label{def:distinct-on}
Sei $\R$ eine Relation, $G=(A_1,\ldots,A_p)$ die in
\texttt{DISTINCT ON} angegebene Attributfolge und
$(B_1,\ldots,B_q)$ die in der \texttt{ORDER BY}-Klausel zusätzlich
angegebenen Attribute mit Sortierrichtungen, sodass die vollständige
Sortierordnung durch $(A_1,\ldots,A_p,B_1,\ldots,B_q)$ gegeben ist.
Dann liefert
\[
  \texttt{SELECT DISTINCT ON } (A_1,\ldots,A_p)\;
  \texttt{\ldots FROM } \R \;
  \texttt{ORDER BY } A_1,\ldots,A_p,B_1,\ldots,B_q
\]
für jeden in $\R$ auftretenden Wert $g=(t[A_1],\ldots,t[A_p])$ genau
ein Tupel: dasjenige Tupel $t_g\in\R_g$, das bezüglich der durch
$(B_1,\ldots,B_q)$ definierten Ordnung $\prec_B$ \emph{minimal} ist
(bzw.\ \emph{maximal}, falls die jeweilige Spalte \texttt{DESC}
sortiert wird).
\end{definition}

\begin{theorem}[Äquivalenz von \texttt{DISTINCT ON} und
\texttt{ROW\_NUMBER()}-Filterung]
\label{thm:distinct-on-aequivalenz}
Die Anweisung
\[
  Q_1 \;=\; \texttt{SELECT DISTINCT ON } (G)\; * \;
  \texttt{FROM } \R \;\texttt{ORDER BY } G, B_1,\ldots,B_q
\]
und die Anweisung
\begin{align*}
  Q_2 \;=\; & \texttt{SELECT * FROM (}\\
  & \quad\texttt{SELECT *, ROW\_NUMBER() OVER}\\
  & \quad\texttt{(PARTITION BY } G \texttt{ ORDER BY }
    B_1,\ldots,B_q\texttt{) AS rn}\\
  & \quad\texttt{FROM } \R\\
  & \texttt{) t WHERE rn = 1}
\end{align*}
liefern für jede Ausprägung $r$ von $\R$ dasselbe Ergebnis (bis auf
die zusätzliche Spalte \texttt{rn} in $Q_2$, die durch eine äußere
Projektion entfernt werden kann), sofern $\prec_B$ eine totale Ordnung
auf jeder Gruppe $\R_g$ induziert (d.\,h.\ keine zwei Tupel derselben
Gruppe sind bezüglich $(B_1,\ldots,B_q)$ identisch -- andernfalls ist
das Ergebnis in beiden Fällen gleichermaßen \emph{nicht-deterministisch}
bezüglich der Auswahl unter den $\prec_B$-gleichwertigen Tupeln).
\end{theorem}

\begin{proof}
Wir zeigen, dass beide Anfragen für jede Gruppe $g\in\pi_G(r)$
dasselbe Tupel $t_g$ liefern.

\emph{$Q_2$ liefert $t_g$:} Nach Definition von \texttt{ROW\_NUMBER()}
mit \texttt{PARTITION BY} $G$ \texttt{ORDER BY} $B_1,\ldots,B_q$
erhält innerhalb der Partition $\R_g$ dasjenige Tupel $t$, für das
$(t[B_1],\ldots,t[B_q])$ bezüglich $\prec_B$ minimal ist, den Rang
$\texttt{rn}=1$ zugewiesen (Definition der Fensterfunktion
\texttt{ROW\_NUMBER}: aufsteigende Nummerierung gemäß der
\texttt{ORDER BY}-Klausel der \texttt{OVER}-Spezifikation, beginnend
bei~1). Da $\prec_B$ nach Voraussetzung total auf $\R_g$ ist, ist
dieses Tupel eindeutig bestimmt und gleich $t_g=\min_{\prec_B}\R_g$
nach Definition~\ref{def:gruppenrepraesentant} (mit $\prec=\prec_B$
auf jeder Gruppe). Der äußere Filter \texttt{WHERE rn = 1} selektiert
somit genau $\{t_g \mid g\in\pi_G(r)\}$.

\emph{$Q_1$ liefert $t_g$:} Nach Definition~\ref{def:distinct-on}
sortiert die \texttt{ORDER BY}-Klausel $G,B_1,\ldots,B_q$ zunächst
alle Tupel von $r$ lexikographisch nach $G$ und führt innerhalb jeder
so entstehenden Gruppe $\R_g$ (Tupel mit gleichem $G$-Wert $g$) die
Sortierung gemäß $\prec_B$ fort. \texttt{DISTINCT ON} $(G)$ behält
sodann -- pro eindeutigem $G$-Wert $g$ -- das in dieser Gesamtordnung
\emph{erste} Tupel der Gruppe $\R_g$. Da innerhalb von $\R_g$ die
Sortierung durch $\prec_B$ bestimmt ist, ist dieses erste Tupel
$\min_{\prec_B}\R_g=t_g$.

Da $Q_1$ und $Q_2$ für jede Gruppe $g$ dasselbe Tupel $t_g$ liefern
und $\pi_G(r)$ in beiden Fällen dieselbe Indexmenge der Gruppen ist,
sind die Ergebnismengen von $Q_1$ und der um \texttt{rn} projizierten
Version von $Q_2$ identisch. \qedhere
\end{proof}

\begin{remark}[Wichtige Syntaxregel]
\label{rem:distinct-on-regel}
Aus dem Beweis von Theorem~\ref{thm:distinct-on-aequivalenz} folgt
unmittelbar die zentrale syntaktische Regel von
\texttt{DISTINCT ON}: Die in \texttt{DISTINCT ON} $(G)$ angegebenen
Ausdrücke \emph{müssen} als \emph{Präfix} der \texttt{ORDER
BY}-Klausel auftreten -- formal: Die \texttt{ORDER BY}-Liste muss von
der Form $(A_1,\ldots,A_p,B_1,\ldots,B_q)$ mit $G=(A_1,\ldots,A_p)$
sein. Andernfalls wäre die im Beweis verwendete Zerlegung der
Gesamtordnung in „zuerst nach $G$, dann innerhalb jeder $G$-Gruppe
nach $(B_1,\ldots,B_q)$“ nicht mehr gültig, und PostgreSQL kann nicht
mehr garantieren, dass das „erste Tupel pro $G$-Gruppe“ auch das
bezüglich $\prec_B$ minimale Tupel ist. PostgreSQL erzwingt diese
Regel durch einen Syntaxfehler zur Ãœbersetzungszeit.
\end{remark}

\subsection{Effizienzvergleich und Bezug zur Tiefenpfadarchitektur}
\label{sec:distinct-on-effizienz}

\begin{proposition}[Operatorbaum von \texttt{DISTINCT ON}]
\label{prop:distinct-on-plan}
Während die \texttt{ROW\_NUMBER()}-Lösung $Q_2$ aus
Theorem~\ref{thm:distinct-on-aequivalenz} im Ausführungsplan
typischerweise die Operatorfolge
\[
  \texttt{Sort} \to \texttt{WindowAgg} \to \texttt{Filter (rn = 1)}
\]
erzeugt -- wobei \texttt{WindowAgg} für \emph{jedes} Tupel von $\R$
eine Rangnummer materialisiert --, kann \texttt{DISTINCT ON} $(G)$ vom
Optimierer als
\[
  \texttt{Sort} \to \texttt{Unique (on } G\texttt{)}
\]
übersetzt werden: Der \texttt{Unique}-Operator durchläuft den
sortierten Strom und gibt -- ohne Materialisierung einer
Rangspalte -- für jede Änderung des $G$-Werts genau das erste neu
angetroffene Tupel weiter. Beide Pläne haben dieselbe asymptotische
Komplexität $O(|\R|\log|\R|)$ (dominiert durch den
\texttt{Sort}-Schritt), jedoch mit einer kleineren Konstanten für
\texttt{DISTINCT ON}, da keine zusätzliche Spalte berechnet,
gespeichert und anschließend wieder herausgefiltert werden muss.
\end{proposition}

\begin{remark}[Bezug zu Theorem~\ref{thm:selektivitaet} und
zum Tiefenpfad]
\label{rem:distinct-on-tiefenpfad}
\texttt{DISTINCT ON} entfaltet seinen größten praktischen Nutzen genau
an den Kanten des Fremdschlüsselgraphen $G_{\mathcal{S}}$
(Definition~\ref{def:er-graph}): Das in
Abschnitt~\ref{sec:row-number} und~\ref{sec:distinct-on-postgres}
durchgängig verwendete Beispiel -- „die jeweils neueste
$\textsc{Auftrag}$-Zeile pro $\textsc{Kunde}$“ -- ist die Anfrage
\emph{entlang} der Kante $(\textsc{Auftrag},\textsc{Kunde})$ des
Tiefenpfads aus Abschnitt~\ref{sec:beispiel}, präzisiert auf
„den jeweils \emph{aktuellsten Repräsentanten}“ pro Knoten der
Elternrelation. Damit ergänzt \texttt{DISTINCT ON} die in
Definition~\ref{def:fk-join} eingeführte Operation des
\emph{kanonischen Verbunds} um eine zweite, ebenso fundamentale
Operation entlang von Tiefenpfadkanten:

\begin{itemize}
  \item Der \emph{kanonische Verbund} $R_i\Join_{F=K}R_{i+1}$
    (Definition~\ref{def:fk-join}) beantwortet die Frage:
    \glqq Zu welchem $R_{i+1}$-Tupel gehört dieses
    $R_i$-Tupel?\grqq{}
  \item Die \emph{\texttt{DISTINCT ON}-Projektion}
    $\textsc{DistinctOn}_{F}(R_i,\prec_B)$ beantwortet die
    \emph{umgekehrte} und in der Praxis ebenso häufige Frage:
    \glqq Welches ist (gemäß $\prec_B$) das repräsentative
    $R_i$-Tupel zu diesem $R_{i+1}$-Tupel?\grqq{}
\end{itemize}

Da -- wie in Theorem~\ref{thm:selektivitaet} gezeigt -- im
Tiefenpfadschema die Relation $R_i$ \emph{nicht} mit den Attributen
aller anderen Entitäten vermischt ist (anders als in einem
Sternschema, in dem die für \texttt{DISTINCT ON} $(G)$
herangezogene Gruppierungsspalte $G$ häufig erst über einen
zusätzlichen Verbund mit der Hub-Relation $H$ ermittelt werden müsste),
kann \texttt{DISTINCT ON} $(F)$ \emph{direkt} auf $R_i$ angewendet
werden, mit einem Index auf $(F,B_1,\ldots,B_q)$ als idealer
Zugriffsstruktur: Ein solcher zusammengesetzter Index erlaubt es dem
\texttt{Unique}-Operator aus
Proposition~\ref{prop:distinct-on-plan}, den
\texttt{Sort}-Schritt vollständig zu entfallen, da der Index die
Tupel bereits in der für \texttt{DISTINCT ON} benötigten Reihenfolge
liefert -- der Ausführungsplan reduziert sich dann auf einen reinen
\texttt{Index Scan} mit \texttt{Unique}-Operator und besitzt Kosten
$O(|\pi_G(R_i)|+\log|R_i|)$ statt $O(|R_i|\log|R_i|)$.
\end{remark}

\begin{example}[\texttt{DISTINCT ON} im Tiefenpfadschema]
\label{ex:distinct-on-tiefenpfad}
Im Tiefenpfadschema aus Abschnitt~\ref{sec:beispiel} liefert
\begin{lstlisting}
CREATE INDEX idx_auftrag_kunde_datum
  ON Auftrag (kunden_id, datum DESC);

SELECT DISTINCT ON (kunden_id)
       kunden_id, auftrag_id, datum, gesamtbetrag
FROM   Auftrag
ORDER BY kunden_id, datum DESC;
\end{lstlisting}
für jeden Kunden den jeweils neuesten Auftrag in
$O(|\pi_{\texttt{kunden\_id}}(\textsc{Auftrag})|)$ -- also
\emph{linear in der Anzahl der Kunden}, unabhängig von der
Gesamtzahl der Aufträge. Im Sternschema, in dem
\texttt{kunden\_id} als denormalisiertes Attribut zunächst über
einen Verbund mit der Hub-Relation $H$ ermittelt werden müsste (vgl.\
Theorem~\ref{thm:selektivitaet}), kann der Index
\texttt{(kunden\_id, datum DESC)} nicht in derselben Weise direkt auf
der Quellrelation des \texttt{DISTINCT ON} angelegt werden, sodass der
\texttt{Sort}-Schritt aus Proposition~\ref{prop:distinct-on-plan} in
der Regel nicht entfällt.
\end{example}

\subsection{Weitere Anwendungsfälle entlang des Tiefenpfads}

Das Gruppenrepräsentantenproblem (Definition~\ref{def:gruppenrepraesentant})
tritt im Tiefenpfadschema systematisch an jeder Kante auf, an der eine
\emph{historisierende} oder \emph{eins-zu-viele}-Beziehung vorliegt.
Typische Anwendungen von \texttt{DISTINCT ON} im Kontext dieses Buches
sind:

\begin{itemize}
  \item \textbf{Aktuellster Status:} Für jede
    $\textsc{Auftrag}$-Zeile den aktuellsten Eintrag einer
    $\textsc{Auftragsstatusverlauf}$-Tabelle ermitteln
    (\texttt{DISTINCT ON (auftrag\_id) \ldots ORDER BY auftrag\_id,
    geaendert\_am DESC}).
  \item \textbf{Erstes Vorkommen:} Für jede $\textsc{Region}$ den
    Kunden mit der kleinsten \texttt{kunden\_id} (z.\,B.\ den
    „dienstältesten“ Kunden) ermitteln
    (\texttt{DISTINCT ON (region\_id) \ldots ORDER BY region\_id,
    kunden\_id ASC}).
  \item \textbf{Extremwert je Gruppe mit vollständigem Tupel:}
    Für jeden $\textsc{Kunde}$ den $\textsc{Auftrag}$ mit dem
    höchsten \texttt{gesamtbetrag} -- inklusive aller übrigen
    Attribute dieses Auftrags, nicht nur des Maximalwerts selbst, was
    mit einem einfachen \texttt{MAX}-Aggregat (Abschnitt 3.7) nicht
    möglich wäre, ohne erneut einen Selbstverbund oder eine
    Fensterfunktion zu verwenden.
\end{itemize}

In allen drei Fällen gilt die in
Bemerkung~\ref{rem:distinct-on-tiefenpfad} entwickelte
Beobachtung: Die Gruppierungsspalte $G$ von \texttt{DISTINCT ON} ist
im Tiefenpfadschema stets ein \emph{lokales} Attribut oder ein
\emph{unmittelbarer} Fremdschlüssel der Quellrelation -- niemals ein
über mehrere Verbundstufen $\delta>1$ (Abschnitt~\ref{sec:abfragen})
vermitteltes Attribut --, was \texttt{DISTINCT ON} zu einem
natürlichen Begleiter des in Kapitel~\ref{ch:tiefenentwurf}
entwickelten Entwurfsparadigmas macht.



Eine View ist eine benannte, gespeicherte SQL-Abfrage, die sich wie eine
Tabelle verwenden lässt:

\begin{lstlisting}
CREATE VIEW AuftragsUebersicht AS
SELECT a.auftrag_id, k.name AS kunde, p.bezeichnung AS produkt
FROM   Auftrag a
JOIN   Kunde k ON a.kunden_id = k.kunden_id
JOIN   Auftragsposition ap ON ap.auftrag_id = a.auftrag_id
JOIN   Produkt p ON p.produkt_id = ap.produkt_id;
\end{lstlisting}

Views realisieren die externe Ebene der ANSI/SPARC-Architektur (vgl.\
Abbildung~\ref{fig:ansi-sparc}). Im tiefenorientierten Entwurf
entsprechen Views typischerweise \emph{materialisierten Pfaden} im
DFS-Wald des Schemas: Eine View, die mehrere Tiefenstufen zu einer
flachen Sicht zusammenfasst, stellt für lesende Anwendungen die
Bequemlichkeit eines Sternschemas bereit, ohne dass das zugrunde
liegende Schema selbst sternförmig sein muss. Dies ist ein zentrales
Argument gegen den Einwand, tiefenorientierte Schemata seien für
Leseoperationen unbequem (vgl.\ Abschnitt~\ref{sec:beweis}).

\section{Datenmodifikationen}

Die Datenmanipulationsoperationen \texttt{INSERT}, \texttt{UPDATE} und
\texttt{DELETE} unterliegen den im Schema definierten
Integritätsbedingungen. Insbesondere erzwingt referentielle Integrität
(Definition~\ref{def:schluessel}) bei \texttt{DELETE}-Operationen auf
einem referenzierten Tupel eine der folgenden Strategien:

\begin{itemize}
  \item \texttt{ON DELETE RESTRICT}: Löschen wird verweigert, falls
    referenzierende Tupel existieren.
  \item \texttt{ON DELETE CASCADE}: Referenzierende Tupel werden
    automatisch mitgelöscht.
  \item \texttt{ON DELETE SET NULL}: Der Fremdschlüsselwert wird auf
    \texttt{NULL} gesetzt.
\end{itemize}

Die Wahl dieser Strategie entlang einer Kante des Fremdschlüsselgraphen
$G_{\mathcal{S}}$ hat unmittelbare Auswirkung auf die in
Abschnitt~\ref{sec:abfragen} analysierte
Schreibkomplexität: In einem Tiefenpfad propagiert \texttt{CASCADE}
entlang einer einzigen, klar definierten Kette, während in einem
Sternschema mit zentralem Hub eine \texttt{CASCADE}-Operation potenziell
\emph{alle} abhängigen Relationen gleichzeitig betrifft.

\section{Darstellung von Datenbankinhalten}

Zur Darstellung von Anfrageergebnissen stehen neben der tabellarischen
Standardausgabe Aggregations- und Gruppierungsfunktionen
(\texttt{COUNT}, \texttt{SUM}, \texttt{AVG}, \texttt{MIN}, \texttt{MAX})
sowie Window-Funktionen (\texttt{OVER (PARTITION BY \ldots ORDER BY
\ldots)}) zur Verfügung. Letztere erlauben es, entlang eines
Tiefenpfads kumulative Kennzahlen (z.\,B. laufende Summen pro
Auftrag) zu berechnen, ohne den natürlichen Verbund mehrfach
auszuführen.

\chapter{Datenbank-Entwurf}
\label{ch:entwurf}

\section{Anomalien in Datenbanken}
\label{sec:anomalien-kurz}

Schlecht entworfene Relationenschemata führen zu drei klassischen
Anomalietypen: Einfüge-, Lösch- und Änderungsanomalien (vgl.\
Abschnitt~\ref{sec:anomalien} für die vollständige formale
Klassifikation und den Beweis von Theorem~\ref{thm:anomaliefrei}).
Ursache aller drei Anomalietypen ist \emph{Redundanz}, die wiederum aus
der Verletzung von Normalformen resultiert.

\section{Normalformen von Relationen}
\label{sec:normalformen-kurz}

Die Normalformenhierarchie 1NF $\supset$ 2NF $\supset$ 3NF $\supset$
BCNF $\supset$ 4NF wird in Abschnitt~\ref{sec:normalformen} formal
eingeführt. Zentral für dieses Buch ist
Theorem~\ref{thm:transitivitaet} (Transitivitätsverteilung): Eine
Kette funktionaler Abhängigkeiten $K \to A_1 \to A_2 \to \cdots \to A_n$
zwingt -- soll 3NF erreicht werden -- zu einer Zerlegung in eine lineare
Kette von Relationen $R_1(K,A_1), R_2(A_1,A_2), \ldots$. Diese Kette ist
nichts anderes als ein \emph{Tiefenpfad} im Sinne von
Definition~\ref{def:masse}.

\section{Dekomposition}
\label{sec:dekomposition}

\begin{definition}[Verlustfreie Dekomposition]
\label{def:lossless}
Eine Zerlegung einer Relation $R$ mit Schema $\att(R)=X\cup Y$ in
$R_1$ mit $\att(R_1)=X$ und $R_2$ mit $\att(R_2)=Y$ heißt
\emph{verlustfrei} (lossless-join), wenn für jede zulässige
Ausprägung $r$ von $R$ gilt:
\[
  r = \pi_X(r) \Join \pi_Y(r).
\]
\end{definition}

\begin{theorem}[Lossless-Join-Kriterium]
\label{thm:lossless}
Die Zerlegung von $R$ in $R_1(X)$ und $R_2(Y)$ mit $X\cup Y=\att(R)$ ist
genau dann verlustfrei, wenn
\[
  (X\cap Y) \to X \quad\text{oder}\quad (X\cap Y) \to Y
\]
in der Menge der funktionalen Abhängigkeiten von $R$ gilt.
\end{theorem}

\begin{proof}
($\Leftarrow$) Gelte $(X\cap Y)\to X$. Sei $t \in
\pi_X(r)\Join\pi_Y(r)$. Dann existieren $t_1,t_2\in r$ mit
$t[X]=t_1[X]$ und $t[Y]=t_2[Y]$ sowie $t_1[X\cap Y]=t_2[X\cap Y]=t[X\cap
Y]$. Da $(X\cap Y)\to X$, folgt $t_1[X]=t_2[X]$ aus $t_1[X\cap
Y]=t_2[X\cap Y]$, also $t_1[X]=t[X]=t_2[X]$. Damit stimmen $t_1$ und
$t_2$ auf $X\cap Y$ und auf $X$ überein; da $X\cup Y=\att(R)$ und $t$
durch $t[X]$ und $t[Y]=t_2[Y]$ vollständig bestimmt ist, gilt $t=t_2\in
r$. Die Inklusion $r \subseteq \pi_X(r)\Join\pi_Y(r)$ gilt stets
trivial, also $r=\pi_X(r)\Join\pi_Y(r)$.
($\Rightarrow$) Der Beweis der Notwendigkeit erfolgt durch
Kontraposition mittels eines Gegenbeispiel-Tableaus, siehe
\cite{Ullman1988}, Theorem 3.18. Gilt weder $(X\cap Y)\to X$ noch
$(X\cap Y)\to Y$, so existieren zwei Tupel $t_1,t_2\in r$ mit
$t_1[X\cap Y]=t_2[X\cap Y]$, aber $t_1[X]\neq t_2[X]$ und $t_1[Y]\neq
t_2[Y]$. Der Verbund $\pi_X(r)\Join\pi_Y(r)$ enthält dann zusätzlich die
„Spurious Tuples“ $(t_1[X],t_2[Y])$ und $(t_2[X],t_1[Y])$, die nicht in
$r$ enthalten sind, womit $r \subsetneq \pi_X(r)\Join\pi_Y(r)$ und die
Zerlegung verlustbehaftet ist. \qedhere
\end{proof}

\begin{remark}
Theorem~\ref{thm:lossless} ist die formale Grundlage des in
Theorem~\ref{thm:transitivitaet} verwendeten
„Lossless-Join-Lemma“: Bei der Zerlegung entlang einer transitiven
Abhängigkeit $K\to A_i\to A_{i+1}$ gilt $A_i \to (A_i,A_{i+1})$
trivial und $A_i \to R'$ nach Voraussetzung, womit das Kriterium von
Theorem~\ref{thm:lossless} mit $X\cap Y=\{A_i\}$ erfüllt ist.
\end{remark}

\begin{definition}[Abhängigkeitserhaltung]
\label{def:dep-erhaltend}
Eine Zerlegung $\{R_1,\ldots,R_k\}$ von $R$ ist
\emph{abhängigkeitserhaltend}, wenn die Vereinigung der auf $R_1,\ldots,
R_k$ projizierten funktionalen Abhängigkeiten die ursprüngliche Menge
$\Sigma$ logisch impliziert, d.\,h.\ $(\bigcup_i \pi_{R_i}(\Sigma))^+ =
\Sigma^+$.
\end{definition}

Der in Theorem~\ref{thm:transitivitaet} bewiesene Tiefenpfad
$R_1(K,A_1), R_2(A_1,A_2),\ldots,R_{n-1}(A_{n-1},A_n)$ ist sowohl
verlustfrei (Theorem~\ref{thm:lossless}, iterativ angewandt) als auch
abhängigkeitserhaltend (Definition~\ref{def:dep-erhaltend}), da jede FD
$A_i \to A_{i+1}$ vollständig in $R_i$ enthalten bleibt. Dies ist die
\emph{stärkste} mögliche Garantie, die ein Dekompositionsalgorithmus
liefern kann, und unterscheidet die Tiefenpfad-Zerlegung von einer
naiven BCNF-Synthese, die Abhängigkeitserhaltung im Allgemeinen nicht
garantiert~\cite{Ullman1988}.

\section{Minimalcover und der Synthesealgorithmus zur Schemaminimierung}
\label{sec:minimalcover}

Die in Abschnitt~\ref{sec:dekomposition} und in
Theorem~\ref{thm:transitivitaet} verwendeten Zerlegungsschritte setzen
implizit voraus, dass die Menge der funktionalen Abhängigkeiten
$\Sigma$ in einer „aufgeräumten“ Form vorliegt, in der keine
Abhängigkeit redundant und kein Attribut auf der linken Seite
überflüssig ist. Diese Form wird als \emph{Minimalcover}
(minimale Hülle, minimal cover, kanonische Überdeckung) bezeichnet.
Der in diesem Abschnitt entwickelte \emph{Synthesealgorithmus}
verwendet das Minimalcover, um aus einer beliebigen Menge $\Sigma$
funktionaler Abhängigkeiten ein Schema mit einer
\emph{minimalen Anzahl von Relationen} zu konstruieren, das
verlustfrei (Definition~\ref{def:lossless}) und
abhängigkeitserhaltend (Definition~\ref{def:dep-erhaltend}) ist und
mindestens 3NF erfüllt. Dieser Abschnitt ist bewusst an dieser Stelle
platziert: Er verbindet die Dekompositionstheorie aus
Abschnitt~\ref{sec:dekomposition} mit der konkreten DDL-Erzeugung in
Abschnitt~\ref{sec:ddl} und liefert damit das algorithmische Bindeglied
zwischen der formalen Theorie (Korollar~\ref{cor:tiefenpfad3NF}) und
der praktischen Schemaerzeugung.

\subsection{Hüllenbildung und Äquivalenz von FD-Mengen}

\begin{definition}[Hülle einer Attributmenge]
\label{def:attributhuelle}
Sei $\Sigma$ eine Menge funktionaler Abhängigkeiten über $\att(R)$ und
$X\subseteq\att(R)$. Die \emph{Hülle} $X^+_\Sigma$ von $X$ bezüglich
$\Sigma$ ist die größte Attributmenge $Y\supseteq X$, sodass
$X\to Y$ aus $\Sigma$ mittels der Armstrong-Axiome
\cite{Armstrong1974} (Reflexivität, Augmentation, Transitivität)
herleitbar ist.
\end{definition}

\begin{definition}[Äquivalenz von FD-Mengen]
\label{def:fd-aequivalent}
Zwei Mengen funktionaler Abhängigkeiten $\Sigma_1$ und $\Sigma_2$ über
demselben Attributschema heißen \emph{äquivalent}, geschrieben
$\Sigma_1 \equiv \Sigma_2$, wenn $\Sigma_1^+ = \Sigma_2^+$, d.\,h.\
wenn jede aus $\Sigma_1$ herleitbare FD auch aus $\Sigma_2$ herleitbar
ist und umgekehrt.
\end{definition}

\subsection{Definition des Minimalcovers}

\begin{definition}[Minimalcover]
\label{def:minimalcover}
Eine Menge funktionaler Abhängigkeiten $\Sigma_{\min}$ heißt
\emph{Minimalcover} von $\Sigma$, wenn $\Sigma_{\min}\equiv\Sigma$
gilt und zusätzlich:
\begin{enumerate}[label=(M\arabic*)]
  \item \label{m1} Jede rechte Seite jeder FD in $\Sigma_{\min}$
    besteht aus genau einem Attribut (\emph{Singleton-RHS}).
  \item \label{m2} Für jede FD $X\to A\in\Sigma_{\min}$ und jedes
    $B\in X$ gilt $(X\setminus\{B\})^+_{\Sigma_{\min}} \neq
    (X\setminus\{B\})^+_{\Sigma_{\min}}\cup\{A\}$, d.\,h.\ kein
    Attribut von $X$ kann aus $X\to A$ entfernt werden, ohne die
    Äquivalenz zu $\Sigma$ zu verletzen
    (\emph{Linksreduziertheit}).
  \item \label{m3} Für jede FD $f\in\Sigma_{\min}$ gilt
    $(\Sigma_{\min}\setminus\{f\})^+ \neq \Sigma_{\min}^+$, d.\,h.\
    keine FD ist redundant (\emph{Redundanzfreiheit}).
\end{enumerate}
\end{definition}

\begin{remark}
Ein Minimalcover ist im Allgemeinen \emph{nicht eindeutig}: Verschiedene
Reihenfolgen bei der Anwendung der Eliminationsschritte~\ref{m2}
und~\ref{m3} können zu unterschiedlichen, aber äquivalenten
Minimalcovern führen \cite{Ullman1988}. Für den Synthesealgorithmus ist
dies unschädlich, da alle Minimalcover dieselbe Anzahl von
„Gruppen“ gleicher linker Seiten und damit dieselbe \emph{Anzahl}
synthetisierter Relationen liefern (Theorem~\ref{thm:synthese-minimal}).
\end{remark}

\subsection{Algorithmus zur Berechnung eines Minimalcovers}

\begin{lstlisting}[language=, caption={Berechnung eines Minimalcovers
$\Sigma_{\min}$ aus einer FD-Menge $\Sigma$}, label=alg:minimalcover,
basicstyle=\ttfamily\small, frame=single, escapeinside={(*}{*)}]
function MINIMALCOVER(Sigma):
    # Schritt 1: Singleton-RHS erzwingen
    G := {}
    for each (X -> Y) in Sigma:
        for each A in Y:
            G := G + { X -> A }

    # Schritt 2: Linksreduktion
    for each (X -> A) in G:
        for each B in X:
            if A in HUELLE(X \ {B}, G):
                # B ist in X ueberfluessig
                ersetze (X -> A) in G durch ((X \ {B}) -> A)
                X := X \ {B}

    # Schritt 3: Entfernen redundanter Abhaengigkeiten
    for each f = (X -> A) in G:
        if A in HUELLE(X, G \ {f}):
            G := G \ {f}

    return G


function HUELLE(X, Sigma):
    # Berechnet die Attributhuelle X^+ bzgl. Sigma
    Ergebnis := X
    repeat
        verandert := false
        for each (Y -> Z) in Sigma:
            if Y subseteq Ergebnis and Z not subseteq Ergebnis:
                Ergebnis := Ergebnis + Z
                verandert := true
    until not verandert
    return Ergebnis
\end{lstlisting}

\subsection{Laufzeitanalyse}

\begin{theorem}[Laufzeit von \textsc{Minimalcover}]
\label{thm:laufzeit-minimalcover}
Sei $n=|\att(R)|$ die Anzahl der Attribute und $m=|\Sigma|$ die Anzahl
der funktionalen Abhängigkeiten in $\Sigma$, wobei jede linke und
rechte Seite höchstens $n$ Attribute umfasst. Dann terminiert
Algorithmus~\ref{alg:minimalcover} und besitzt eine Laufzeit von
$O(m^3 n^2)$.
\end{theorem}

\begin{proof}
Wir analysieren die drei Schritte einzeln.

\emph{Hilfsfunktion \textsc{Huelle}:} Die \texttt{repeat}-Schleife
fügt in jedem Durchlauf, der nicht zum Abbruch führt, mindestens ein
Attribut zu \texttt{Ergebnis} hinzu (sonst wäre \texttt{verandert =
false} und die Schleife terminierte). Da $|\texttt{Ergebnis}|\leq n$,
sind höchstens $n$ Durchläufe möglich, in denen jeweils alle $m$ FDs
aus $\Sigma$ auf Anwendbarkeit geprüft werden, was $O(m)$ pro
Durchlauf kostet. Damit ist $\textsc{Huelle}\in O(mn)$.

\emph{Schritt~1 (Singleton-RHS):} Jede der $m$ FDs wird in höchstens
$n$ FDs mit Singleton-RHS aufgespalten, also $|G|\in O(mn)$ und
Schritt~1 selbst kostet $O(mn)$.

\emph{Schritt~2 (Linksreduktion):} Für jede der $|G|=O(mn)$ FDs $X\to
A$ wird für jedes der höchstens $n$ Attribute $B\in X$ ein Aufruf von
\textsc{Huelle} durchgeführt. Dies ergibt
$O(mn)\cdot O(n)\cdot O(mn)=O(m^2n^3)$ im worst case. Da jedoch $|X|\leq
n$ und in der Praxis $|X|\ll n$, geben wir die Schranke konservativ in
Abhängigkeit von $n$ und $m$ an: $O(m^2 n^3)$.

\emph{Schritt~3 (Redundanzentfernung):} Für jede der $|G|=O(mn)$ FDs
wird ein Aufruf von \textsc{Huelle} auf der um $f$ reduzierten Menge
$G\setminus\{f\}$ durchgeführt, was $O(mn)\cdot O(mn)=O(m^2n^2)$
kostet.

Die Gesamtlaufzeit ist somit
$O(mn)+O(m^2n^3)+O(m^2n^2)=O(m^2n^3)$.
Da in praktisch relevanten Schemata $m=O(n)$ gilt (die Anzahl der FDs
ist durch die Anzahl der Attribute beschränkt, da jedes Attribut
typischerweise von höchstens konstant vielen Schlüsselkandidaten
abhängt), vereinfacht sich dies zu $O(n^5)$ in $n$; in der allgemeinen,
von $n$ \emph{und} $m$ unabhängig parametrisierten Form, geben wir die
obere Schranke $O(m^2 n^3) \subseteq O(m^3 n^2)$ für $m\geq n$ konservativ
mit $O(m^3n^2)$ an, was die Behauptung zeigt. \qedhere
\end{proof}

\begin{remark}
Für die in dieser Arbeit betrachteten Tiefenpfadschemata
(Definition~\ref{def:tiefenpfad}) ist jede Relation $R_i$ klein:
$|\att(R_i)|=O(1)$ und $|\Sigma_{R_i}|=O(1)$, da jede Relation gemäß
Theorem~\ref{thm:transitivitaet} höchstens die FD
$A_{i-1}\to A_i$ sowie die Schlüsselbedingung enthält. Die in
Theorem~\ref{thm:laufzeit-minimalcover} gezeigte Laufzeit ist also
\emph{pro Relation} des Tiefenpfads konstant, und die
Gesamtlaufzeit für das gesamte Schema $\mathcal{S}$ mit $l$ Relationen
ist $O(l)$ -- linear in der Pfadlänge. Im Gegensatz dazu würde eine
naive Anwendung des Algorithmus auf eine \emph{einzige}, undekomponierte
Hub-Relation $H$ eines $k$-Sternschemas mit $|\att(H)|=O(k)$ und
$|\Sigma_H|=O(k)$ eine Laufzeit von $O(k^5)$ erfordern -- ein weiteres,
algorithmisches Argument für die in Kapitel~\ref{ch:tiefenentwurf}
entwickelte Präferenz für Tiefenpfade gegenüber Sternschemata.
\end{remark}

\subsection{Der Synthesealgorithmus zur Schemaminimierung}

Mit dem Minimalcover aus Algorithmus~\ref{alg:minimalcover} lässt sich
nun der vollständige \emph{Synthesealgorithmus} formulieren, der aus
einem Relationenschema $R$ mit Attributmenge $\att(R)$ und FD-Menge
$\Sigma$ ein \emph{minimales} Mengen von Relationen $\{R_1,\ldots,R_k\}$
erzeugt, das die Eigenschaften aus Definition~\ref{def:lossless} und
Definition~\ref{def:dep-erhaltend} erfüllt und mindestens 3NF
garantiert \cite{Ullman1988,Date2003}.

\begin{lstlisting}[language=, caption={Synthesealgorithmus zur
Erzeugung eines minimalen 3NF-Schemas}, label=alg:synthese,
basicstyle=\ttfamily\small, frame=single]
function SYNTHESE_3NF(attr(R), Sigma):
    Sigma_min := MINIMALCOVER(Sigma)

    # Schritt 1: Gruppierung nach linker Seite
    Gruppen := {}
    for each (X -> A) in Sigma_min:
        Gruppen[X] := Gruppen[X] + {A}      # FDs mit gleicher LHS X
                                              # zu einer Relation buendeln

    # Schritt 2: Eine Relation pro Gruppe
    Relationen := {}
    for each X in Gruppen.keys():
        R_X := X union Gruppen[X]
        Relationen := Relationen + { R_X }

    # Schritt 3: Schluesselrelation hinzufuegen, falls keine
    #            der erzeugten Relationen einen Schluessel
    #            von R enthaelt
    K := EIN_SCHLUESSEL(attr(R), Sigma)
    if not exists R_X in Relationen with K subseteq attr(R_X):
        Relationen := Relationen + { K }

    # Schritt 4: Redundante Relationen entfernen
    for each R_X in Relationen:
        if exists R_Y in Relationen, R_Y != R_X,
                  with attr(R_X) subseteq attr(R_Y):
            Relationen := Relationen \ { R_X }

    return Relationen
\end{lstlisting}

\begin{remark}[Schritt~3: Schlüsselermittlung]
Die Hilfsfunktion \texttt{EIN\_SCHLUESSEL} bestimmt einen
(nicht notwendig eindeutigen) Schlüssel $K\subseteq\att(R)$ gemäß
Definition~\ref{def:schluessel}, etwa durch den klassischen Algorithmus
„starte mit $K=\att(R)$ und entferne iterativ Attribute $B$, solange
$(K\setminus\{B\})^+_\Sigma=\att(R)$“. Dieser Schritt hat
Laufzeit $O(n)\cdot O(\textsc{Huelle})=O(mn^2)$ und ändert die in
Theorem~\ref{thm:laufzeit-minimalcover} bewiesene Gesamtkomplexität
nicht.
\end{remark}

\subsection{Korrektheit des Synthesealgorithmus}

\begin{theorem}[Korrektheit von \textsc{Synthese\_3NF}]
\label{thm:synthese-korrekt}
Das von Algorithmus~\ref{alg:synthese} erzeugte Schema
$\mathcal{R}=\{R_1,\ldots,R_k\}$ besitzt die folgenden drei
Eigenschaften:
\begin{enumerate}
  \item[(K1)] \textbf{Abhängigkeitserhaltung:} $(\bigcup_{i}
    \pi_{R_i}(\Sigma))^+ = \Sigma^+$.
  \item[(K2)] \textbf{Verlustfreiheit:} Für jede zulässige Ausprägung
    $r$ von $R$ gilt $r = \Join_{i=1}^{k}\pi_{R_i}(r)$.
  \item[(K3)] \textbf{3NF:} Jede Relation $R_i\in\mathcal{R}$ ist in
    dritter Normalform.
\end{enumerate}
\end{theorem}

\begin{proof}
\textbf{(K1) Abhängigkeitserhaltung.}
Nach Konstruktion enthält für jede FD $X\to A\in\Sigma_{\min}$ die
Relation $R_X=X\cup\text{Gruppen}[X]$ sowohl $X$ als auch $A$, also
gilt $X\to A \in \pi_{R_X}(\Sigma)$ für jede solche FD. Damit ist
$\Sigma_{\min}\subseteq\bigcup_i\pi_{R_i}(\Sigma)$, also
$\Sigma_{\min}^+\subseteq(\bigcup_i\pi_{R_i}(\Sigma))^+$. Da nach
Definition~\ref{def:minimalcover} $\Sigma_{\min}\equiv\Sigma$, also
$\Sigma_{\min}^+=\Sigma^+$, und da andererseits jede projizierte FD
$\pi_{R_i}(\Sigma)$ aus $\Sigma$ (und damit aus $\Sigma^+$) ableitbar
ist, gilt auch
$(\bigcup_i\pi_{R_i}(\Sigma))^+\subseteq\Sigma^+$. Beide Inklusionen
zusammen ergeben (K1).

\textbf{(K2) Verlustfreiheit.}
Nach Schritt~3 des Algorithmus enthält $\mathcal{R}$ eine Relation
$R_K\supseteq K$ für einen Schlüssel $K$ von $R$. Da $K\to\att(R)$
(Definition~\ref{def:schluessel}, jeder Schlüssel bestimmt alle
Attribute), gilt nach dem Lossless-Join-Kriterium
(Theorem~\ref{thm:lossless}), iterativ angewandt auf
$R_K$ und jede weitere Relation $R_i$: Da $K\subseteq\att(R_K)$ und
$K\to\att(R)\supseteq\att(R_i)$, gilt insbesondere
$(\att(R_K)\cap\att(R_i))\to\att(R_i)$ für jedes $i$ (denn
$K\subseteq\att(R_K)\cap\att(R_i)$ würde dies implizieren, falls
$K\subseteq \att(R_i)$; im allgemeinen Fall folgt die Aussage aus
$K^+_\Sigma=\att(R)\supseteq \att(R_K)\cap\att(R_i)$ und der
Transitivität der Hüllenbildung, sodass
$(\att(R_K)\cap\att(R_i))^+\supseteq K^+=\att(R)\supseteq\att(R_i)$).
Damit ist die Zerlegung $\{R_K\}\cup\{R_i\}$ für jedes $i$ verlustfrei
nach Theorem~\ref{thm:lossless}, und durch sukzessive Anwendung dieses
Arguments (jede weitere Relation wird verlustfrei mit dem bereits
verlustfrei rekonstruierten Teil $R_K\Join R_{i_1}\Join\cdots$
verbunden, da $K$ weiterhin in jedem Zwischenergebnis als
Determinante für $\att(R)$ enthalten ist) ist die Gesamtzerlegung
$\mathcal{R}$ verlustfrei, d.\,h.\ es gilt (K2).

\textbf{(K3) 3NF.}
Sei $R_X\in\mathcal{R}$ mit $\att(R_X)=X\cup\text{Gruppen}[X]$ und sei
$B\in\text{Gruppen}[X]$ ein Nichtschlüsselattribut von $R_X$ (d.\,h.\
$B\notin K'$ für jeden Schlüssel $K'$ von $R_X$), das transitiv von
einem Schlüssel $K'$ von $R_X$ abhängt: $K'\to C\to B$ mit $C\not\to
K'$ und $B\notin K'\cup C$. Da $K'\to C$ und $C\to B$ beide in
$\pi_{R_X}(\Sigma_{\min}^+)$ liegen müssten (3NF-Verletzungen sind FDs
über $\att(R_X)$), und da nach Konstruktion \emph{jede} FD aus
$\Sigma_{\min}$ mit linker Seite $X$ bereits vollständig in $R_X$
abgebildet ist (alle FDs mit LHS $=X$ wurden zu \texttt{Gruppen}[X]
zusammengefasst), könnte $C\to B$ nur dann in $R_X$ auftreten, wenn
entweder $C=X$ (dann wäre $X\to B$ bereits in $\Sigma_{\min}$ mit
LHS $X$, also $B\in\text{Gruppen}[X]$ kein Widerspruch zur
3NF, da dann $X\to B$ direkt vom Schlüssel $X$ abhängt -- keine
Transitivität) oder $C\subsetneq X\cup\text{Gruppen}[X]$ mit $C\neq X$
und $C\to B\in\Sigma_{\min}$. Im letzteren Fall hätte aber
$\Sigma_{\min}$ zwei FDs $X\to B$ und $C\to B$ mit $C\subsetneq
\att(R_X)$, $C\neq X$ -- dies widerspricht der Linksreduziertheit
\ref{m2} bzw.\ Redundanzfreiheit \ref{m3} von $\Sigma_{\min}$, da dann
entweder $X\to B$ aus $C\to B$ und $X\supseteq C$ (Augmentation)
herleitbar und somit redundant wäre, oder $C\to B$ würde bereits
implizieren, dass $X$ um die in $X\setminus C$ enthaltenen Attribute
linksreduziert werden könnte. Beide Fälle widersprechen der Annahme,
dass $\Sigma_{\min}$ ein Minimalcover ist. Also kann eine solche
transitive Abhängigkeit innerhalb $R_X$ nicht auftreten, und $R_X$
erfüllt 3NF. Da dies für jede Relation $R_X\in\mathcal{R}$ sowie für
die in Schritt~3 hinzugefügte Schlüsselrelation $R_K$ (die nach
Konstruktion keine Nichtschlüsselattribute mit transitiven
Abhängigkeiten enthält, da $\att(R_K)=K$ oder $K$ plus höchstens
durch $K$ direkt bestimmte Attribute) gilt, erfüllt $\mathcal{R}$
insgesamt (K3). \qedhere
\end{proof}

\begin{theorem}[Minimalität der Relationenanzahl]
\label{thm:synthese-minimal}
Die Anzahl $k=|\mathcal{R}|$ der von Algorithmus~\ref{alg:synthese}
erzeugten Relationen ist minimal unter allen abhängigkeitserhaltenden
3NF-Zerlegungen von $R$ bezüglich $\Sigma$, bis auf das
Hinzufügen einer einzelnen zusätzlichen Schlüsselrelation in
Schritt~3.
\end{theorem}

\begin{proof}
Nach Schritt~1 und~2 entspricht $k$ (vor Schritt~3 und~4) der Anzahl
der \emph{verschiedenen linken Seiten} in $\Sigma_{\min}$. Sei
$\mathcal{R}'=\{R'_1,\ldots,R'_{k'}\}$ eine beliebige
abhängigkeitserhaltende 3NF-Zerlegung von $R$. Nach (K1) muss für jede
FD $X\to A\in\Sigma_{\min}$ (und $\Sigma_{\min}\equiv\Sigma$) ein
$R'_j\in\mathcal{R}'$ existieren mit $X\cup\{A\}\subseteq\att(R'_j)$,
da $X\to A$ andernfalls in keiner projizierten FD-Menge
$\pi_{R'_j}(\Sigma)$ vollständig enthalten wäre und somit die
Abhängigkeitserhaltung verletzt würde (eine FD, deren linke und rechte
Seite auf verschiedene Relationen verteilt sind, kann nicht durch das
DBMS lokal überprüft werden, vgl.\ Abschnitt~\ref{sec:codd-rules}).
Da $\Sigma_{\min}$ nach Konstruktion redundanzfrei ist
(Eigenschaft~\ref{m3}), kann durch das in Schritt~4 von
Algorithmus~\ref{alg:synthese} durchgeführte Entfernen von
Teilmengenrelationen keine FD aus $\Sigma_{\min}$ „doppelt“
abgedeckt werden, ohne dass eine der beiden abdeckenden Relationen in
der anderen enthalten wäre -- in diesem Fall wird durch Schritt~4
genau die überflüssige Relation entfernt. Es folgt, dass
$\mathcal{R}'$ mindestens so viele Relationen benötigt, wie es
\emph{nach Entfernung von Teilmengenbeziehungen} verschiedene
„maximale“ Gruppen von FDs mit gleicher oder ineinander enthaltener
linker Seite in $\Sigma_{\min}$ gibt -- und dies ist genau die Anzahl
$k$ (vor Schritt~3) der von $\textsc{Synthese\_3NF}$ erzeugten
Relationen. Schritt~3 fügt höchstens eine weitere Relation hinzu (für
den Fall, dass kein $R_X$ bereits einen Schlüssel von $R$ enthält);
diese ist für (K2) notwendig (vgl.\ Beweis von Theorem~\ref{thm:synthese-korrekt},
(K2)) und kann daher nicht eingespart werden, ohne (K2) zu verletzen.
Damit ist $k$ minimal bis auf diese eine zusätzliche Relation. \qedhere
\end{proof}

\begin{example}[Anwendung auf den Tiefenpfad]
\label{ex:synthese-tiefenpfad}
Sei $\att(R)=\{K,A_1,\ldots,A_n\}$ mit $\Sigma=\{K\to A_1\to A_2\to
\cdots\to A_n\}$ wie in Theorem~\ref{thm:transitivitaet}. Das
Minimalcover ist
$\Sigma_{\min}=\{K\to A_1,\ A_1\to A_2,\ \ldots,\ A_{n-1}\to A_n\}$
(jede dieser FDs ist linksreduziert -- die linke Seite besteht bereits
aus einem einzigen Attribut -- und redundanzfrei, da $A_i\to A_{i+1}$
nicht aus den übrigen FDs ohne $A_i\to A_{i+1}$ selbst herleitbar ist).
Die $n$ verschiedenen linken Seiten $K,A_1,\ldots,A_{n-1}$ liefern $n$
Relationen
\[
  R_1=(K,A_1),\quad R_2=(A_1,A_2),\quad\ldots,\quad
  R_n=(A_{n-1},A_n),
\]
identisch mit der in Theorem~\ref{thm:transitivitaet} angegebenen
Zerlegungskette. Da $K$ ein Schlüssel von $R$ ist und bereits
$K\subseteq\att(R_1)$, entfällt Schritt~3. Algorithmus~\ref{alg:synthese}
reproduziert somit \emph{exakt} den Tiefenpfad
$R_1\to R_2\to\cdots\to R_n$ aus
Korollar~\ref{cor:tiefenpfad3NF} -- der Synthesealgorithmus ist also
nicht nur ein allgemeines Minimierungsverfahren, sondern liefert für
Schemata mit linearer Abhängigkeitskette \emph{automatisch} das in
dieser Arbeit favorisierte Tiefenpfadschema, ohne dass dessen Struktur
im Voraus bekannt sein müsste.
\end{example}


\label{sec:ddl}

Die Definition eines Relationenschemas erfolgt mittels \texttt{CREATE
TABLE}. Für den Tiefenpfad
$\textsc{Land} \leftarrow \textsc{Region} \leftarrow \textsc{Kunde}
\leftarrow \textsc{Auftrag} \leftarrow \textsc{Auftragsposition}$
(vgl.\ Abschnitt~\ref{sec:beispiel}) lautet die DDL:

\begin{lstlisting}
CREATE TABLE Land (
  land_id    INTEGER PRIMARY KEY,
  name       VARCHAR(60) NOT NULL
);

CREATE TABLE Region (
  region_id  INTEGER PRIMARY KEY,
  name       VARCHAR(60) NOT NULL,
  land_id    INTEGER NOT NULL
             REFERENCES Land(land_id)
);

CREATE TABLE Kunde (
  kunden_id  INTEGER PRIMARY KEY,
  name       VARCHAR(100) NOT NULL,
  region_id  INTEGER NOT NULL
             REFERENCES Region(region_id)
);

CREATE TABLE Auftrag (
  auftrag_id INTEGER PRIMARY KEY,
  datum      DATE NOT NULL,
  kunden_id  INTEGER NOT NULL
             REFERENCES Kunde(kunden_id)
);

CREATE TABLE Auftragsposition (
  position_id INTEGER PRIMARY KEY,
  auftrag_id  INTEGER NOT NULL
              REFERENCES Auftrag(auftrag_id)
              ON DELETE CASCADE,
  produkt_id  INTEGER NOT NULL,
  menge       INTEGER NOT NULL CHECK (menge > 0)
);
\end{lstlisting}

Jede \texttt{REFERENCES}-Klausel definiert eine Kante im
Fremdschlüsselgraphen $G_{\mathcal{S}}$ (Definition~\ref{def:er-graph},
Proposition~\ref{prop:er-rel-graph}). Die obige Kette von fünf
Tabellen bildet einen Tiefenpfad der Tiefe $\depth=4$ gemäß
Definition~\ref{def:masse} und realisiert exakt die in
Theorem~\ref{thm:transitivitaet} geforderte 3NF-Zerlegungskette für die
transitive Abhängigkeit $\text{Auftrag} \to \text{Kunde} \to
\text{Region} \to \text{Land}$.

Integritätsbedingungen werden zusätzlich über \texttt{CHECK}-Klauseln
(Wertebereichseinschränkungen), \texttt{UNIQUE}-Constraints
(Alternativschlüssel) und \texttt{NOT NULL} (Pflichtattribute)
formuliert. Alle diese Bedingungen sind gemäß der
Integritätsunabhängigkeit (Abschnitt~\ref{sec:codd-rules}) Teil des
Schemas und werden vom DBMS unabhängig von der Anwendung durchgesetzt.

\chapter{Datenbankbetrieb}
\label{ch:betrieb}

\section{Das Transaktionskonzept in Datenbanksystemen}

Eine \emph{Transaktion} ist eine Folge von Datenbankoperationen, die als
logische Einheit (ACID, vgl.\ Kapitel~\ref{ch:einleitung}) ausgeführt
wird:

\begin{lstlisting}
BEGIN;
UPDATE Konto SET saldo = saldo - 100 WHERE konto_id = 1;
UPDATE Konto SET saldo = saldo + 100 WHERE konto_id = 2;
COMMIT;
\end{lstlisting}

\begin{definition}[Schedule, Serialisierbarkeit]
Ein \emph{Schedule} ist eine zeitliche Verschränkung der Operationen
mehrerer Transaktionen $T_1,\ldots,T_n$. Ein Schedule heißt
\emph{serialisierbar}, wenn seine Wirkung auf die Datenbank derjenigen
einer seriellen (nacheinander ausgeführten) Reihenfolge der
Transaktionen entspricht.
\end{definition}

\section{Mehrbenutzerbetrieb}

Im Mehrbenutzerbetrieb können ohne Synchronisationsmechanismen folgende
Anomalien auftreten: \emph{Lost Update}, \emph{Dirty Read},
\emph{Non-Repeatable Read} und \emph{Phantom Read}. SQL definiert vier
Isolationsstufen (\texttt{READ UNCOMMITTED}, \texttt{READ COMMITTED},
\texttt{REPEATABLE READ}, \texttt{SERIALIZABLE}), die jeweils eine
Teilmenge dieser Anomalien ausschließen.

\begin{remark}[Sperrgranularität entlang des Tiefenpfads]
\label{rem:sperren}
In einem Tiefenpfadschema (Kapitel~\ref{ch:tiefenentwurf}) betrifft eine
Transaktion, die ein Tupel der Relation $R_i$ ändert, typischerweise nur
Sperren auf $R_i$ und gegebenenfalls $R_{i-1}$ (referentielle
Prüfung). Im Sternschema mit zentralem Hub $H$ konkurrieren hingegen
\emph{alle} Transaktionen, die irgendeine der in $H$ denormalisierten
Entitäten betreffen, um Sperren auf $H$ -- ein direkter Bezug zur in
Abschnitt~\ref{sec:wartbarkeit} bewiesenen lokalen Änderbarkeit
(Theorem~\ref{thm:schaemaevolution}), übertragen auf die Dimension der
Nebenläufigkeit: Lokale Schemaänderungen und lokale Sperrkonflikte sind
zwei Ausprägungen derselben strukturellen Eigenschaft des
Tiefenpfads.
\end{remark}

\section{Datenschutz und Zugriffsrechte}
\label{sec:zugriffsrechte}

Zugriffsrechte werden mit \texttt{GRANT} und \texttt{REVOKE} verwaltet:

\begin{lstlisting}
GRANT SELECT ON Kunde TO analyst_rolle;
GRANT SELECT, INSERT, UPDATE ON Auftrag TO sachbearbeiter_rolle;
REVOKE INSERT ON Auftrag FROM sachbearbeiter_rolle;
\end{lstlisting}

Die in Abschnitt~\ref{sec:wartbarkeit} bewiesene
„Berechtigungsgranularität“ tiefer Pfade (Proposition zur
Zugriffsrechteverwaltung) zeigt sich hier konkret: Die Rolle
\texttt{analyst\_rolle} erhält Lesezugriff exakt auf
\textsc{Kunde}, ohne implizit Zugriff auf \textsc{Auftrag} oder
\textsc{Auftragsposition} zu erlangen, da diese Information in einem
Sternschema mit gemeinsamem Hub nicht trennbar wäre.

\section{Benutzung von Systemtabellen}
\label{sec:systemtabellen}

Relationale DBMS stellen einen \emph{Systemkatalog} (Data Dictionary)
bereit, der das Schema selbst als relationale Daten beschreibt
(Beispiele: \texttt{ALL\_TABLES}, \texttt{ALL\_CONSTRAINTS},
\texttt{ALL\_TAB\_COLUMNS} bei Oracle, bzw.\
\texttt{information\_schema} im SQL-Standard).

\begin{example}[Extraktion des Fremdschlüsselgraphen aus dem Katalog]
\label{ex:katalog-graph}
Der Fremdschlüsselgraph $G_{\mathcal{S}}$ (Definition~\ref{def:er-graph})
lässt sich direkt aus dem Systemkatalog rekonstruieren:
\begin{lstlisting}
SELECT
  c.table_name  AS quelle,
  ccu.table_name AS ziel
FROM information_schema.table_constraints c
JOIN information_schema.constraint_column_usage ccu
  ON c.constraint_name = ccu.constraint_name
WHERE c.constraint_type = 'FOREIGN KEY';
\end{lstlisting}
Das Ergebnis dieser Anfrage ist die Kantenmenge $E$ von
$G_{\mathcal{S}}$. Die in Abschnitt~\ref{sec:topologien}
definierten Maße $\depth(G_{\mathcal{S}})$, $\breadth(G_{\mathcal{S}})$
und $\tau(G_{\mathcal{S}})$ können somit \emph{automatisiert} und
\emph{werkzeuggestützt} für jedes beliebige existierende
Datenbankschema berechnet werden -- ein praktisches Instrument zur
Bewertung bestehender Schemata gemäß dem in
Kapitel~\ref{ch:tiefenentwurf} entwickelten Gütekriterium.
\end{example}

\chapter{PL/SQL}
\label{ch:plsql}

\section{Blockstruktur}

PL/SQL erweitert SQL um prozedurale Elemente. Ein PL/SQL-Block gliedert
sich in:

\begin{lstlisting}
DECLARE
  -- Deklaration von Variablen, Cursorn, Ausnahmen
BEGIN
  -- Ausfuehrbare Anweisungen
EXCEPTION
  -- Fehlerbehandlung
END;
\end{lstlisting}

\section{Ein Beispiel}

\begin{lstlisting}
DECLARE
  v_kunden_id Kunde.kunden_id%TYPE := 1;
  v_name      Kunde.name%TYPE;
BEGIN
  SELECT name INTO v_name
  FROM Kunde
  WHERE kunden_id = v_kunden_id;
  DBMS_OUTPUT.PUT_LINE('Kunde: ' || v_name);
END;
\end{lstlisting}

\section{Kontrollstrukturen}

PL/SQL stellt \texttt{IF-THEN-ELSIF-ELSE}, \texttt{LOOP},
\texttt{WHILE LOOP} und \texttt{FOR LOOP} bereit.

\section{Das Cursor-Konzept}

Ein \emph{Cursor} ermöglicht die zeilenweise Verarbeitung eines
Anfrageergebnisses:

\begin{lstlisting}
DECLARE
  CURSOR c_auftraege IS
    SELECT auftrag_id, datum FROM Auftrag
    WHERE kunden_id = 1;
  v_auftrag_id Auftrag.auftrag_id%TYPE;
  v_datum      Auftrag.datum%TYPE;
BEGIN
  OPEN c_auftraege;
  LOOP
    FETCH c_auftraege INTO v_auftrag_id, v_datum;
    EXIT WHEN c_auftraege%NOTFOUND;
    DBMS_OUTPUT.PUT_LINE(v_auftrag_id || ': ' || v_datum);
  END LOOP;
  CLOSE c_auftraege;
END;
\end{lstlisting}

\section{Fehlerbehandlung}

Ausnahmen werden im \texttt{EXCEPTION}-Block behandelt, entweder
vordefiniert (\texttt{NO\_DATA\_FOUND}, \texttt{TOO\_MANY\_ROWS}) oder
benutzerdefiniert mittels \texttt{RAISE\_APPLICATION\_ERROR}.

\section{Prozeduren und Funktionen}

\begin{lstlisting}
CREATE OR REPLACE FUNCTION gesamtwert_auftrag(
  p_auftrag_id IN Auftrag.auftrag_id%TYPE
) RETURN NUMBER IS
  v_summe NUMBER := 0;
BEGIN
  SELECT SUM(menge * preis) INTO v_summe
  FROM Auftragsposition ap
  JOIN Produkt p ON p.produkt_id = ap.produkt_id
  WHERE ap.auftrag_id = p_auftrag_id;
  RETURN NVL(v_summe, 0);
END;
\end{lstlisting}

Diese Funktion realisiert eine \emph{aggregierende Traversierung}
entlang zweier Kanten des Tiefenpfads
($\textsc{Auftragsposition}\to\textsc{Auftrag}$ und
$\textsc{Auftragsposition}\to\textsc{Produkt}$). Da beide Relationen
gemäß Theorem~\ref{thm:selektivitaet} indiziert über ihren
Primärschlüssel zugreifbar sind, ist die Join-Tiefe $\delta=1$ pro
Aggregationsschritt -- konstant und unabhängig von der Gesamtgröße der
Datenbank.

\section{Trigger}

\begin{lstlisting}
CREATE OR REPLACE TRIGGER trg_auftrag_audit
AFTER INSERT OR UPDATE OR DELETE ON Auftrag
FOR EACH ROW
BEGIN
  INSERT INTO Auftrag_Audit(auftrag_id, aktion, zeitpunkt)
  VALUES (
    COALESCE(:NEW.auftrag_id, :OLD.auftrag_id),
    CASE
      WHEN INSERTING THEN 'INSERT'
      WHEN UPDATING  THEN 'UPDATE'
      WHEN DELETING  THEN 'DELETE'
    END,
    SYSTIMESTAMP
  );
END;
\end{lstlisting}

\begin{remark}[Trigger und lokale Änderbarkeit]
Der obige Trigger ist an genau eine Relation, $\textsc{Auftrag}$,
gebunden. Im Tiefenpfadschema entspricht dies der in
Theorem~\ref{thm:schaemaevolution} bewiesenen lokalen
Schemaoperation: Das Hinzufügen eines Audit-Triggers für
$\textsc{Kunde}$ erfordert keinerlei Änderung an
$\textsc{Auftrag}$-Triggern oder umgekehrt. Im Sternschema mit
denormalisierten Kundendaten in der Hub-Relation $H$ müsste ein
analoger Audit-Trigger auf $H$ \emph{alle} Änderungen an \emph{allen}
denormalisierten Entitäten protokollieren und im Triggercode
unterscheiden, was die zyklomatische Komplexität des Triggers linear
mit der Anzahl der Spokes wachsen lässt.
\end{remark}

\chapter{Datenbankzugriff in höheren Programmiersprachen}
\label{ch:hostsprachen}

\section{Das Problem der Impedanzdiskrepanz}
\label{sec:impedanz}

Relationale Datenbanken arbeiten mengenorientiert: Eine Anfrage liefert
eine Relation, also eine \emph{Menge} von Tupeln, als Ergebnis. Höhere
Programmiersprachen wie C, Java, Perl oder Python hingegen sind
\emph{prozedural} bzw.\ objektorientiert und verarbeiten Daten in der
Regel sequentiell, Element für Element, in Form von Variablen, Arrays,
Listen oder Objekten. Diese begriffliche Kluft wird in der Literatur
als \emph{Impedanzdiskrepanz} (impedance mismatch) bezeichnet.

\begin{definition}[Impedanzdiskrepanz]
\label{def:impedanz}
Die Impedanzdiskrepanz zwischen einem relationalen Datenmodell
$\mathcal{S}=\{R_1,\ldots,R_m\}$ und einer Hostsprache $\mathcal{H}$
besteht aus drei Teilproblemen:
\begin{enumerate}
  \item \textbf{Mengen- vs.\ Elementorientierung:} Das Ergebnis einer
    \texttt{SELECT}-Anweisung ist eine Relation; $\mathcal{H}$ kennt
    jedoch primär einzelne Werte oder lineare Datenstrukturen.
  \item \textbf{Typsystemdiskrepanz:} SQL-Datentypen
    (\texttt{NUMERIC}, \texttt{DATE}, \texttt{NULL}\ldots) besitzen
    keine direkte Entsprechung in $\mathcal{H}$; insbesondere
    \texttt{NULL} als „fehlender Wert“ (vgl.\
    Abschnitt~\ref{sec:codd-rules}) muss explizit auf ein Konstrukt von
    $\mathcal{H}$ (z.\,B.\ \texttt{Optional}, Zeiger \texttt{NULL},
    Sentinel-Werte) abgebildet werden.
  \item \textbf{Strukturdiskrepanz:} Die Tiefenpfadstruktur
    $G_{\mathcal{S}}$ (Definition~\ref{def:er-graph}) eines
    normalisierten Schemas wird in $\mathcal{H}$ häufig als Menge
    \emph{verschachtelter} Objekte modelliert (z.\,B.\ ein
    \texttt{Auftrag}-Objekt mit eingebettetem \texttt{Kunde}-Objekt),
    obwohl die Datenbank diese Information auf mehrere Relationen
    verteilt speichert.
\end{enumerate}
\end{definition}

Die in den folgenden Abschnitten behandelten Zugriffstechniken --
Embedded SQL, Cursor-basierte APIs, DBI und objektrelationales Mapping
-- lassen sich alle als unterschiedliche \emph{Strategien zur
Auflösung} der drei Teilprobleme aus Definition~\ref{def:impedanz}
auffassen. Ein zentrales Ergebnis dieses Kapitels ist, dass die
Tiefenpfadstruktur des zugrunde liegenden Schemas (Kapitel~\ref{ch:tiefenentwurf})
nicht nur die Datenbankseite, sondern auch die Architektur des
Anwendungscodes maßgeblich vereinfacht, da -- wie in
Abschnitt~\ref{sec:modularisierung} formal begründet wird -- die
Struktur $G_{\mathcal{S}}$ direkt auf die Modul- bzw.\
Klassenstruktur der Anwendung abgebildet werden kann.

\section{Embedded SQL}
\label{sec:embedded-sql}

Embedded SQL erlaubt das Einbetten von SQL-Anweisungen direkt in
Quellcode einer Hostsprache, eingeleitet durch \texttt{EXEC SQL}. Ein
Präcompiler (z.\,B.\ \texttt{proc} bei Oracle, \texttt{ecpg} bei
PostgreSQL) übersetzt die \texttt{EXEC SQL}-Direktiven in Aufrufe der
Client-Bibliothek des DBMS, bevor der reguläre Compiler der Hostsprache
zum Einsatz kommt.

\subsection{Deklarationsabschnitt und Hostvariablen}

Variablen der Hostsprache, die in SQL-Anweisungen verwendet werden
sollen (\emph{Hostvariablen}), müssen in einem
\texttt{DECLARE}-Abschnitt bekannt gemacht werden:

\begin{lstlisting}[language=C]
EXEC SQL BEGIN DECLARE SECTION;
  int   kunden_id;
  char  name[101];
  char  region_name[61];
  short name_ind;        /* Indikatorvariable fuer NULL */
EXEC SQL END DECLARE SECTION;
\end{lstlisting}

Die \emph{Indikatorvariable} (\texttt{name\_ind}) löst Teilproblem~2
aus Definition~\ref{def:impedanz}: Da C keinen nativen
\texttt{NULL}-Wert für \texttt{char[]} kennt, signalisiert das DBMS
über die Indikatorvariable, ob der zugehörige Spaltenwert
\texttt{NULL} war ($-1$), abgeschnitten wurde (positiver Wert) oder
regulär vorliegt ($0$).

\subsection{Einzeltupel-Zugriff}

\begin{lstlisting}[language=C]
EXEC SQL SELECT k.name, r.name
  INTO :name :name_ind, :region_name
  FROM Kunde k JOIN Region r ON k.region_id = r.region_id
  WHERE k.kunden_id = :kunden_id;

if (sqlca.sqlcode == 100) {
  /* SQLCODE 100: keine Zeile gefunden */
} else if (sqlca.sqlcode < 0) {
  /* Fehlerbehandlung, vgl. Abschnitt 7.2.4 */
}
\end{lstlisting}

Diese Anfrage realisiert den kanonischen Verbund
$\textsc{Kunde}\Join_{\texttt{region\_id}=\texttt{region\_id}}\textsc{Region}$
gemäß Definition~\ref{def:fk-join} -- also eine einzelne Kante des
Tiefenpfads -- und überträgt das Ergebnis in genau zwei Hostvariablen.
Solange Anfragen genau ein Tupel liefern (was bei einer Anfrage über
einen Primärschlüssel garantiert ist, vgl.\
Definition~\ref{def:schluessel}), genügt die einfache
\texttt{INTO}-Klausel; Teilproblem~1 aus
Definition~\ref{def:impedanz} stellt sich hier nicht.

\subsection{Mengenorientierter Zugriff: Cursor in Embedded SQL}
\label{sec:embedded-cursor}

Liefert eine Anfrage mehrere Tupel, muss Teilproblem~1
(Mengen- vs.\ Elementorientierung) explizit aufgelöst werden. Hierzu
dient ein \emph{Cursor}, der die mengenorientierte Relation in einen
sequentiell durchlaufbaren Strom von Tupeln überführt:

\begin{lstlisting}[language=C]
EXEC SQL DECLARE c_auftraege CURSOR FOR
  SELECT a.auftrag_id, a.datum
  FROM Auftrag a
  WHERE a.kunden_id = :kunden_id
  ORDER BY a.datum;

EXEC SQL OPEN c_auftraege;

EXEC SQL BEGIN DECLARE SECTION;
  int  v_auftrag_id;
  char v_datum[11];
EXEC SQL END DECLARE SECTION;

while (1) {
  EXEC SQL FETCH c_auftraege INTO :v_auftrag_id, :v_datum;
  if (sqlca.sqlcode == 100) break;  /* keine weiteren Zeilen */
  verarbeite_auftrag(v_auftrag_id, v_datum);
}
EXEC SQL CLOSE c_auftraege;
\end{lstlisting}

\begin{remark}[Cursor als Iterator über einen Pfadabschnitt]
\label{rem:cursor-iterator}
Ein Cursor realisiert in der Terminologie der Hostsprache einen
\emph{Iterator} über das Ergebnis einer Anfrage. Im Tiefenpfadschema
(Kapitel~\ref{ch:tiefenentwurf}) entspricht ein solcher Cursor
typischerweise der Iteration über alle Kindknoten eines festen
Elternknotens entlang einer Kante von $G_{\mathcal{S}}$ -- im Beispiel
oben: alle $\textsc{Auftrag}$-Tupel, die auf einen festen
$\textsc{Kunde}$ verweisen. Diese Eins-zu-eins-Korrespondenz zwischen
„Cursor über $R_i$ gefiltert nach Fremdschlüssel auf $R_{i-1}$“ und
„Kante $(R_i,R_{i-1})$ in $G_{\mathcal{S}}$“ ist die Grundlage der in
Abschnitt~\ref{sec:repository-pattern} eingeführten
Repository-Architektur.
\end{remark}

\subsection{Transaktionssteuerung und Fehlerbehandlung}
\label{sec:embedded-transaktionen}

Transaktionsgrenzen (Kapitel~\ref{ch:betrieb}) werden in Embedded SQL
explizit über \texttt{EXEC SQL COMMIT WORK} bzw.\ \texttt{EXEC SQL
ROLLBACK WORK} gesetzt. Die \texttt{SQLCA} (SQL Communication Area)
liefert nach jeder \texttt{EXEC SQL}-Anweisung einen Statuscode
(\texttt{sqlca.sqlcode}), über den Fehler -- etwa Verletzungen von
\texttt{CHECK}- oder Fremdschlüsselbedingungen
(Abschnitt~\ref{sec:ddl}) -- erkannt werden:

\begin{lstlisting}[language=C]
EXEC SQL WHENEVER SQLERROR GOTO error_handler;
EXEC SQL WHENEVER NOT FOUND CONTINUE;

EXEC SQL INSERT INTO Auftragsposition
  (position_id, auftrag_id, produkt_id, menge)
  VALUES (:pos_id, :auftrag_id, :produkt_id, :menge);

EXEC SQL COMMIT WORK;
goto fertig;

error_handler:
  printf("Fehler %ld: %s\n", sqlca.sqlcode, sqlca.sqlerrm.sqlerrmc);
  EXEC SQL ROLLBACK WORK;
fertig:
  ;
\end{lstlisting}

\begin{remark}[Lokalität von Integritätsverletzungen im Tiefenpfad]
Eine \texttt{INSERT}-Anweisung auf $\textsc{Auftragsposition}$ kann
nur gegen Integritätsbedingungen verstoßen, die $\textsc{Auftragsposition}$
selbst betreffen, sowie gegen die Fremdschlüsselbedingungen auf den
unmittelbar benachbarten Knoten $\textsc{Auftrag}$ und
$\textsc{Produkt}$ in $G_{\mathcal{S}}$. Der Fehlerbehandlungscode kann
sich somit auf eine kleine, im Voraus bekannte Menge möglicher
Fehlerursachen beschränken -- ein direkter Anwendungsfall der in
Theorem~\ref{thm:schaemaevolution} bewiesenen Lokalität von
Schemaoperationen, übertragen auf die Lokalität der
Fehlerbehandlung.
\end{remark}

\section{Cursor- und API-basierter Zugriff: ODBC und JDBC}
\label{sec:odbc-jdbc}

Neben Embedded SQL, das einen Präcompiler erfordert, existieren
\emph{Bibliotheks-APIs}, die ohne Präcompiler auskommen und daher
plattform- und sprachübergreifend einsetzbar sind. Die wichtigsten
Vertreter sind ODBC (Open Database Connectivity, ursprünglich für C)
und JDBC (Java Database Connectivity).

\subsection{Grundstruktur einer JDBC-Anwendung}

\begin{lstlisting}[language=Java]
String sql =
  "SELECT a.auftrag_id, a.datum, k.name " +
  "FROM Auftrag a JOIN Kunde k ON a.kunden_id = k.kunden_id " +
  "WHERE k.region_id = ?";

try (Connection con = DriverManager.getConnection(url, user, pass);
     PreparedStatement ps = con.prepareStatement(sql)) {

  ps.setInt(1, regionId);

  try (ResultSet rs = ps.executeQuery()) {
    while (rs.next()) {
      int    auftragId = rs.getInt("auftrag_id");
      Date   datum     = rs.getDate("datum");
      String kunde     = rs.getString("name");
      verarbeite(auftragId, datum, kunde);
    }
  }
} catch (SQLException ex) {
  // Fehlerbehandlung, vgl. Abschnitt 7.3.3
}
\end{lstlisting}

\subsection{Vorbereitete Anweisungen (Prepared Statements)}
\label{sec:prepared-statements}

Ein \texttt{PreparedStatement} trennt die \emph{Struktur} einer
SQL-Anweisung (das „Was“) von ihren \emph{Parametern} (das „Womit“).
Dies hat zwei Konsequenzen:

\begin{enumerate}
  \item \textbf{Sicherheit:} Parameterwerte werden nie textuell in die
    SQL-Anweisung eingefügt, wodurch SQL-Injection (vgl.\
    Kapitel~\ref{ch:www}) strukturell ausgeschlossen wird.
  \item \textbf{Effizienz:} Das DBMS kann den Ausführungsplan für die
    Struktur der Anweisung einmalig bestimmen (vgl.\
    Abschnitt~\ref{sec:abfragen}) und für wiederholte Ausführungen mit
    unterschiedlichen Parameterwerten wiederverwenden.
\end{enumerate}

\begin{proposition}[Wiederverwendbarkeit entlang eines Tiefenpfads]
\label{prop:prepared-pfad}
Sei $R_1 \to R_2 \to \cdots \to R_l$ ein Tiefenpfad gemäß
Definition~\ref{def:tiefenpfad}. Für jede Kante $(R_i,R_{i+1})$ existiert
genau \emph{eine} kanonische Verbundanfrage gemäß
Definition~\ref{def:fk-join}, parametrisiert über den Primärschlüssel
von $R_{i+1}$. Eine Anwendung, die das gesamte Schema $\mathcal{S}$
bedient, benötigt somit höchstens $|E(G_{\mathcal{S}})|$ verschiedene
\texttt{PreparedStatement}-Vorlagen für Einzelschritt-Traversierungen --
unabhängig von der Anzahl der Tupel in $\mathcal{S}$.
\end{proposition}

\begin{proof}
Nach Definition~\ref{def:fk-join} ist der kanonische Verbund entlang
einer Kante $(R_i,R_j)$ durch die Verbundbedingung $F=K$ eindeutig
bestimmt, wobei $F\subseteq\att(R_i)$ der Fremdschlüssel und
$K\subseteq\att(R_j)$ der Primärschlüssel von $R_j$ ist; diese
Bedingung hängt nur vom Schema $\mathcal{S}$, nicht von einer
konkreten Instanz ab. Für jede der $|E(G_{\mathcal{S}})|$ Kanten gibt
es somit genau eine solche Bedingung, also höchstens
$|E(G_{\mathcal{S}})|$ unterschiedliche Anfragestrukturen. Da
$|E(G_{\mathcal{S}})| = \depth(G_{\mathcal{S}})$ für einen reinen
Tiefenpfad (Definition~\ref{def:masse}) und
$\depth(G_{\mathcal{S}})$ unabhängig von der Datenmenge ist, folgt die
Behauptung. \qedhere
\end{proof}

In einem $k$-Sternschema (Definition~\ref{def:stern}) hingegen
existieren $k$ Kanten von der Hub-Relation zu den Spoke-Relationen,
\emph{zusätzlich} aber häufig $\binom{k}{2}$ potenzielle
Mehrfachverbund-Kombinationen für Anfragen, die mehrere Spokes
gleichzeitig betreffen, da diese nicht über kurze Kettenanfragen,
sondern über den zentralen Hub vermittelt werden müssen. Die Anzahl
der benötigten \texttt{PreparedStatement}-Vorlagen wächst hier im
Allgemeinen überlinear in $k$, während sie im Tiefenpfad nach
Proposition~\ref{prop:prepared-pfad} linear in $\depth(G_{\mathcal{S}})$
bleibt.

\subsection{Fehlerbehandlung und Ressourcenverwaltung}

JDBC signalisiert Fehler über die Ausnahmeklasse
\texttt{SQLException}, die einen herstellerspezifischen
\texttt{SQLState}-Code (nach SQL-Standard genormt) enthält. Klasse
\texttt{23xxx} kennzeichnet beispielsweise Verletzungen von
Integritätsbedingungen (\texttt{23503} = Verletzung einer
Fremdschlüsselbedingung). Die in Java~7 eingeführte
\texttt{try-with-resources}-Syntax (siehe obiges Beispiel) garantiert,
dass \texttt{Connection}, \texttt{PreparedStatement} und
\texttt{ResultSet} auch im Fehlerfall korrekt geschlossen werden --
eine Notwendigkeit, da offene Cursor (Bemerkung~\ref{rem:cursor-iterator})
serverseitige Ressourcen und gegebenenfalls Sperren
(Bemerkung~\ref{rem:sperren}) belegen.

\section{Die DBI-Schnittstelle für Perl}
\label{sec:perl-dbi}

Die \texttt{DBI}-Schnittstelle (Database Interface) abstrahiert den
Datenbankzugriff in Perl unabhängig vom konkreten DBMS über
austauschbare \texttt{DBD}-Treiber (Database Driver):

\begin{lstlisting}[language=Perl]
use DBI;
my $dbh = DBI->connect("dbi:Pg:dbname=vertrieb", $user, $pass,
                        { RaiseError => 1, AutoCommit => 0 });

my $sth = $dbh->prepare(
  "SELECT a.auftrag_id, a.datum FROM Auftrag a WHERE a.kunden_id = ?"
);
$sth->execute($kunden_id);

while (my $row = $sth->fetchrow_hashref) {
  print "$row->{auftrag_id}: $row->{datum}\n";
}
$sth->finish;
$dbh->commit;
$dbh->disconnect;
\end{lstlisting}

\subsection{Fehlerbehandlung mit \texttt{RaiseError} und \texttt{eval}}

Mit \texttt{RaiseError => 1} wirft jeder DBI-Methodenaufruf bei einem
Fehler eine Perl-Exception, die mit \texttt{eval}/\texttt{\$@}
abgefangen werden kann:

\begin{lstlisting}[language=Perl]
eval {
  $dbh->do(
    "INSERT INTO Auftragsposition (position_id, auftrag_id, " .
    "produkt_id, menge) VALUES (?, ?, ?, ?)",
    undef, $pos_id, $auftrag_id, $produkt_id, $menge
  );
  $dbh->commit;
};
if ($@) {
  warn "Fehler beim Einfuegen: $@";
  $dbh->rollback;
}
\end{lstlisting}

\subsection{Transaktionsklammern und AutoCommit}

Mit \texttt{AutoCommit => 0} wird jede Folge von DML-Anweisungen
(Abschnitt~\ref{sec:select} ff.) implizit zu einer Transaktion gemäß
Kapitel~\ref{ch:betrieb}, bis explizit \texttt{commit} oder
\texttt{rollback} aufgerufen wird. Dies ist insbesondere bei
mehrstufigen Einfügeoperationen entlang eines Tiefenpfads relevant: Das
Anlegen eines neuen Auftrags mit mehreren Positionen erfordert
\texttt{INSERT}-Anweisungen auf $\textsc{Auftrag}$ und mehrfach auf
$\textsc{Auftragsposition}$; erst der gemeinsame
\texttt{commit} stellt sicher, dass entweder \emph{alle} diese
Operationen sichtbar werden oder \emph{keine} (Atomarität, vgl.\
Kapitel~\ref{ch:einleitung}).

\section{Portabilität der Zugriffsverfahren}
\label{sec:portabilitaet}

Sowohl Embedded SQL als auch API-basierte Schnittstellen (ODBC, JDBC,
DBI) abstrahieren von der konkreten DBMS-Implementierung, sofern
ausschließlich Standard-SQL-Konstrukte verwendet werden. Diese
Portabilität wird durch die in Abschnitt~\ref{sec:codd-rules}
geforderte Strukturunabhängigkeit ermöglicht. Tabelle~\ref{tab:vergleich-zugriff}
fasst die in diesem Kapitel behandelten Zugriffsverfahren vergleichend
zusammen.

\begin{table}[H]
\centering
\small
\begin{tabular}{@{}lp{3.45cm}p{3.45cm}p{3.45cm}@{}}
\toprule
& \textbf{Embedded SQL} & \textbf{ODBC/JDBC} & \textbf{DBI (Perl)} \\
\midrule
Übersetzung & Präcompiler nötig & keine, reine API & keine, reine API \\
Mengenzugriff & Cursor (Abschn.~\ref{sec:embedded-cursor}) &
  \texttt{ResultSet}-Iterator & \texttt{fetchrow\_*}-Iterator \\
Parametrisierung & Hostvariablen & \texttt{PreparedStatement} &
  Platzhalter \texttt{?} \\
\texttt{NULL}-Behandlung & Indikatorvariablen & \texttt{wasNull()} &
  \texttt{undef} \\
Fehlerbehandlung & \texttt{SQLCA}, \texttt{WHENEVER} &
  \texttt{SQLException} & \texttt{RaiseError}/\texttt{eval} \\
\bottomrule
\end{tabular}
\caption{Vergleich der in Kapitel~\ref{ch:hostsprachen} behandelten
Zugriffsverfahren hinsichtlich der drei Teilprobleme der
Impedanzdiskrepanz (Definition~\ref{def:impedanz}).}
\label{tab:vergleich-zugriff}
\end{table}

\section{Objektrelationales Mapping (ORM) und das Tiefenpfadparadigma}
\label{sec:orm}

Objektrelationale Mapper (ORM, z.\,B.\ Hibernate, JPA, SQLAlchemy,
Doctrine) lösen Teilproblem~3 aus Definition~\ref{def:impedanz}
(Strukturdiskrepanz) durch deklarative \emph{Mapping-Beschreibungen},
die Relationen auf Klassen und Fremdschlüsselbeziehungen auf
Objektreferenzen abbilden:

\begin{lstlisting}[language=Java]
@Entity
class Auftrag {
  @Id int auftragId;
  Date datum;

  @ManyToOne
  @JoinColumn(name = "kunden_id")
  Kunde kunde;

  @OneToMany(mappedBy = "auftrag")
  List<Auftragsposition> positionen;
}
\end{lstlisting}

\subsection{Lazy Loading entlang des Tiefenpfads}
\label{sec:lazy-loading}

ORMs implementieren häufig \emph{Lazy Loading}: Eine assoziierte
Entität (z.\,B.\ \texttt{auftrag.kunde}) wird erst bei tatsächlichem
Zugriff aus der Datenbank nachgeladen, nicht bereits beim Laden des
\texttt{Auftrag}-Objekts.

\begin{proposition}[Lazy Loading als Pfadtraversierung]
\label{prop:lazy-loading}
Sei $R_1\to R_2\to\cdots\to R_l$ ein Tiefenpfad und sei ein Objekt
$o_1$ vom Typ $R_1$ geladen. Der Zugriff auf die assoziierte Entität
$o_2$ vom Typ $R_2$ über Lazy Loading entspricht genau der Ausführung
des kanonischen Verbunds $R_1\Join_{F=K}R_2$ aus
Definition~\ref{def:fk-join}, beschränkt auf das Tupel $o_1$ -- also
einer Anfrage mit Join-Tiefe $\delta=1$ im Sinne von
Abschnitt~\ref{sec:abfragen}. Eine vollständige Traversierung des
Pfads bis $R_l$ durch sukzessive Lazy-Loading-Zugriffe
$o_1\to o_2\to\cdots\to o_l$ erzeugt somit $l-1$ Einzelanfragen mit
$\delta=1$, anstatt einer einzigen Anfrage mit $\delta=l-1$
(Proposition~\ref{prop:jointiefe}).
\end{proposition}

\begin{remark}[Das $N{+}1$-Problem]
\label{rem:n-plus-1}
Wird Proposition~\ref{prop:lazy-loading} unreflektiert auf eine
\emph{Liste} von $n$ Objekten $o_1^{(1)},\ldots,o_1^{(n)}$ angewandt
(z.\,B.\ alle Aufträge eines Tages, mit anschließendem Zugriff auf
\texttt{kunde} für jeden Auftrag), entstehen $1+n$ Anfragen: eine zum
Laden der $n$ Aufträge, sowie $n$ weitere Einzelanfragen für die
jeweils zugehörigen Kunden. Dies ist das in der Praxis als
„$N{+}1$-Problem“ bekannte Performance-Antimuster. Es ist
\emph{kein} Argument gegen Tiefenpfadschemata an sich, sondern gegen
die unreflektierte Verwendung von Lazy Loading bei Listenzugriffen:
Die Lösung besteht darin, den kanonischen Verbund explizit
anzufordern (\texttt{JOIN FETCH} in JPQL bzw.\ \texttt{eager
loading}), wodurch wieder die in
Proposition~\ref{prop:jointiefe} beschriebene Situation mit
$\delta=1$ für die \emph{gesamte} Liste (statt $\delta=1$ pro Element)
entsteht. Der Tiefenpfad selbst bleibt von dieser
Optimierungsentscheidung unberührt; sie betrifft ausschließlich die
\emph{Strategie} des ORM beim Ãœbersetzen von Objektgraph-Zugriffen in
SQL-Anfragen.
\end{remark}

\subsection{Kaskadierende Operationen im ORM}

Analog zu \texttt{ON DELETE CASCADE} auf Datenbankebene
(Abschnitt~\ref{sec:select}, Datenmodifikationen) bieten ORMs
\emph{Cascade-Optionen} (z.\,B.\ \texttt{cascade =
CascadeType.REMOVE}) an, die das Löschen eines Objekts auf assoziierte
Objekte propagieren. Im Tiefenpfadschema betrifft eine solche Kaskade
genau den Teilbaum unterhalb des gelöschten Knotens in
$G_{\mathcal{S}}$ (vgl.\ Definition~\ref{def:dfswald}); im
Sternschema mit zentralem Hub müsste eine analoge
Cascade-Konfiguration auf der Hub-Klasse \emph{alle} Spoke-Assoziationen
auflisten, was die Wartbarkeit der Mapping-Konfiguration im Sinne von
Theorem~\ref{thm:schaemaevolution} beeinträchtigt.

\section{Modularisierung von Anwendungscode entlang des Fremdschlüsselgraphen}
\label{sec:modularisierung}

\subsection{Das Repository-Pattern}
\label{sec:repository-pattern}

Ein in der Praxis weit verbreitetes Architekturmuster ist das
\emph{Repository-Pattern}: Für jede Relation $R_i\in\mathcal{S}$
existiert ein Modul bzw.\ eine Klasse \texttt{$R_i$Repository}, das
sämtliche Datenbankzugriffe auf $R_i$ kapselt (CRUD-Operationen,
Cursor gemäß Bemerkung~\ref{rem:cursor-iterator}, Prepared Statements
gemäß Proposition~\ref{prop:prepared-pfad}).

\begin{theorem}[Eindeutigkeit der Repository-Zerlegung im Tiefenpfad]
\label{thm:repository-zerlegung}
Sei $\mathcal{S}=\{R_1,\ldots,R_l\}$ ein Tiefenpfadschema gemäß
Definition~\ref{def:tiefenpfad} mit Fremdschlüsselgraph
$G_{\mathcal{S}}=R_1\to R_2\to\cdots\to R_l$. Dann existiert eine
Zerlegung des Anwendungscodes in $l$ Repository-Module
$M_1,\ldots,M_l$ mit $M_i \leftrightarrow R_i$, sodass
\begin{enumerate}
  \item jede SQL-Anweisung des Anwendungscodes, die ausschließlich
    Attribute von $R_i$ liest oder schreibt, vollständig in $M_i$
    enthalten ist,
  \item jede SQL-Anweisung, die einen kanonischen Verbund
    $R_i\Join_{F=K}R_{i+1}$ (Definition~\ref{def:fk-join}) realisiert,
    entweder in $M_i$ oder in $M_{i+1}$ enthalten ist, und
  \item kein Modul $M_i$ direkten Code-Zugriff auf eine Relation
    $R_j$ mit $|i-j|>1$ benötigt.
\end{enumerate}
\end{theorem}

\begin{proof}
Wir definieren $M_i$ als die Menge aller SQL-Anweisungen, deren
\texttt{FROM}- bzw.\ \texttt{JOIN}-Klauseln ausschließlich Relationen
aus $\{R_{i-1},R_i,R_{i+1}\}\cap\mathcal{S}$ referenzieren und die
mindestens $R_i$ enthalten, zugeordnet zu $R_i$ mit der kleinsten
Indexsumme im Falle einer Mehrdeutigkeit (d.\,h.\ ein Verbund
$R_i\Join R_{i+1}$ wird $M_i$ zugeordnet). Eigenschaft~1 gilt
unmittelbar nach Konstruktion, da eine Anweisung über $R_i$ allein
notwendig in $\{R_{i-1},R_i,R_{i+1}\}\cap\mathcal{S}$ liegt und somit
$M_i$ zugeordnet wird (oder $M_{i-1}$/$M_{i+1}$, falls die
Mehrdeutigkeitsregel anders entscheidet -- in jedem Fall aber
\emph{einem} Modul vollständig). Eigenschaft~2 gilt nach Konstruktion
direkt aus der Definition von $M_i$. Für Eigenschaft~3 betrachten wir
ein Modul $M_i$ und eine Relation $R_j$ mit $|i-j|>1$. Da
$G_{\mathcal{S}}$ ein Tiefenpfad ist (Definition~\ref{def:tiefenpfad}),
existiert nach Definition~\ref{def:er-graph} keine Kante
$(R_i,R_j)$ für $|i-j|>1$, also auch kein Fremdschlüssel zwischen
$R_i$ und $R_j$. Folglich kann nach
Definition~\ref{def:fk-join} kein \emph{kanonischer} Verbund zwischen
$R_i$ und $R_j$ formuliert werden, ohne die dazwischenliegenden
Relationen $R_{i+1},\ldots,R_{j-1}$ (bzw.\ $R_{j+1},\ldots,R_{i-1}$) zu
durchlaufen -- eine solche mehrstufige Anfrage wird nach Konstruktion
nicht $M_i$, sondern keinem der so definierten Module direkt
zugeordnet, sondern -- wie in der Praxis üblich -- einer
übergeordneten \emph{Service}-Schicht, die mehrere Repositories
komponiert, ohne selbst auf $R_j$ zuzugreifen. Damit benötigt $M_i$
selbst keinen direkten Zugriff auf $R_j$. \qedhere
\end{proof}

\begin{remark}[Sternschema verletzt Theorem~\ref{thm:repository-zerlegung}]
Für ein $k$-Sternschema (Definition~\ref{def:stern}) mit Hub $H$ und
Spokes $S_1,\ldots,S_k$ existiert \emph{keine} Zerlegung in $k+1$
Module mit den Eigenschaften~1--3: Da jede Spoke-Relation $S_j$ nur
über $H$ erreichbar ist und $H$ mit \emph{allen} $S_j$ über
Fremdschlüsselkanten verbunden ist (Definition~\ref{def:stern}, $\deg(H)=k$),
müsste das Modul $M_H$ für $H$ nach Eigenschaft~2 sämtliche
$k$ kanonischen Verbunde $H\Join S_j$ enthalten. $M_H$ wächst damit
linear mit $k$ und verletzt für $k\geq 2$ die Idee einer auf
\emph{eine} Relation fokussierten Modulgrenze. Dies ist die
anwendungsarchitektonische Entsprechung der in
Theorem~\ref{thm:schaemaevolution} bewiesenen mangelnden lokalen
Änderbarkeit von Sternschemata.
\end{remark}

\subsection{Service-Schicht für Mehrfach-Pfad-Operationen}

Operationen, die mehrere nicht benachbarte Knoten von
$G_{\mathcal{S}}$ betreffen -- etwa „Berechne den Gesamtumsatz aller
Kunden einer Region“ ($\textsc{Region}\leftarrow\textsc{Kunde}
\leftarrow\textsc{Auftrag}\leftarrow\textsc{Auftragsposition}$, also
Pfadlänge~3 --, werden gemäß Theorem~\ref{thm:repository-zerlegung}
in einer \emph{Service}-Schicht oberhalb der Repositories
implementiert. Eine solche Service-Methode komponiert die Aufrufe
mehrerer Repositories oder formuliert -- aus Effizienzgründen, vgl.\
Abschnitt~\ref{sec:abfragen} -- eine einzelne mehrstufige
SQL-Anfrage mit $\delta=3$. In beiden Fällen bleibt die in
Theorem~\ref{thm:repository-zerlegung} bewiesene Zerlegung der
\emph{Datenzugriffsschicht} unangetastet: Die Service-Schicht
\emph{verwendet} die Module $M_{\text{Region}},\ldots,
M_{\text{Auftragsposition}}$, ohne deren interne Struktur zu
duplizieren.

\section{Zusammenfassung}

Dieses Kapitel hat gezeigt, dass die drei Teilprobleme der
Impedanzdiskrepanz (Definition~\ref{def:impedanz}) -- Mengen- vs.\
Elementorientierung, Typsystemdiskrepanz und Strukturdiskrepanz -- von
Embedded SQL, ODBC/JDBC, Perl-DBI und ORMs auf unterschiedliche, aber
strukturell verwandte Weise gelöst werden (Cursor, Prepared
Statements, Lazy Loading). In allen Fällen erweist sich die
Tiefenpfadstruktur $G_{\mathcal{S}}$ des zugrunde liegenden Schemas als
der zentrale Bauplan für die Architektur des Anwendungscodes:
Theorem~\ref{thm:repository-zerlegung} zeigt, dass ein Tiefenpfadschema
eine \emph{eindeutige}, lokal konsistente Zerlegung des
Datenzugriffscodes in Repository-Module induziert, während ein
Sternschema diese Eigenschaft verletzt. Damit setzt sich die in
Kapitel~\ref{ch:tiefenentwurf} formal etablierte Ãœberlegenheit
tiefenorientierter Schemata konsequent bis in die
Anwendungsarchitektur fort.

\chapter{WWW-Integration von Datenbanken}
\label{ch:www}

\section{Kommandozeilen- und graphische Benutzeroberflächen}

Klassische Benutzeroberflächen (CLI- oder GUI-basiert) greifen direkt
über Treiber (ODBC, JDBC) auf das DBMS zu. Mit der Verbreitung des
World Wide Web wird zunehmend eine dreischichtige Architektur
eingesetzt: Präsentationsschicht (Browser), Anwendungsschicht
(Web-Server, Skriptsprache) und Datenhaltungsschicht (DBMS).

\section{Benutzeroberflächen im World Wide Web}

Eine Web-Oberfläche zur Pflege des Tiefenpfadschemas aus
Abschnitt~\ref{sec:beispiel} besteht typischerweise aus
einer Seite pro Relation $R_i$ -- also pro Tiefenstufe -- mit Links zu
den jeweils referenzierten bzw.\ referenzierenden Datensätzen.

\section{Architektur einer WWW-Oberfläche mit CGI-Skripten}

\begin{figure}[H]
\centering
\begin{tikzpicture}[
  box/.style={draw, rectangle, minimum width=3cm, minimum height=1cm},
  node distance=1.2cm
]
\node[box] (browser) {Browser};
\node[box, right=of browser] (webserver) {Web-Server (CGI)};
\node[box, right=of webserver] (dbms) {DBMS};
\draw[->] (browser) -- node[above,font=\tiny]{HTTP-Request} (webserver);
\draw[->] (webserver) -- node[above,font=\tiny]{SQL} (dbms);
\draw[<-, bend left=30] (browser.north) to node[above,font=\tiny]{HTML} (webserver.north -| browser.east |- webserver.north);
\end{tikzpicture}
\caption{Dreischichtige Architektur: Browser, Web-Server mit CGI-Skript, DBMS}
\end{figure}

\section{Sammeln von Eingaben mit Web-Formularen}

Web-Formulare (\texttt{<form>}) übermitteln Benutzereingaben per
\texttt{GET} oder \texttt{POST} an ein CGI-Skript, das die Eingaben
validiert und in parametrisierte SQL-Anweisungen einsetzt.

\section{Auswertung von Formulareingaben mit CGI-Skripten}

\begin{lstlisting}[language=Perl]
my $cgi = CGI->new;
my $kunden_id = $cgi->param('kunden_id');

# Parametrisierte Anfrage gegen SQL-Injection
my $sth = $dbh->prepare(
  "SELECT * FROM Kunde WHERE kunden_id = ?"
);
$sth->execute($kunden_id);
\end{lstlisting}

\begin{remark}[Sicherheit durch Parametrisierung]
Die Verwendung parametrisierter Anfragen (Platzhalter \texttt{?}) ist
unabhängig von der Schematopologie zwingend erforderlich, um
SQL-Injection zu verhindern. Im Tiefenpfadschema betrifft eine solche
Anfrage typischerweise genau eine Relation $R_i$, wodurch die
Angriffsfläche pro Formular auf das Schema von $R_i$ begrenzt bleibt.
\end{remark}

\section{Erzeugen von Formularen mit CGI-Skripten}

Formulare zur Eingabe neuer Datensätze für $R_i$ können automatisiert
aus dem Systemkatalog (Abschnitt~\ref{sec:systemtabellen}) generiert
werden: Für jedes Attribut von $R_i$ wird ein Eingabefeld erzeugt, für
jeden Fremdschlüssel ein Auswahlfeld (Dropdown), dessen Optionen aus
der referenzierten Relation $R_{i+1}$ stammen.

\section{Interaktion mit Web-Formularen}

Mehrstufige Formulare (Wizards) bilden Tiefenpfade direkt ab: Schritt
$1$ erfasst Daten für $R_1$, Schritt $2$ für $R_2$ unter Verwendung des
in Schritt~1 erzeugten Primärschlüssels als Fremdschlüssel, usw.

\section{Effizienz von CGI-Skripten}

Der wiederholte Prozessstart klassischer CGI-Skripte verursacht
Overhead. Alternativen sind \texttt{mod\_perl}, \texttt{FastCGI} oder
moderne Anwendungsserver, die Datenbankverbindungen poolen.

\section{CGI-Skripten im Mehrbenutzerbetrieb}

Verbindungs-Pooling muss mit den in
Abschnitt~\ref{sec:zugriffsrechte} und
Bemerkung~\ref{rem:sperren} diskutierten Aspekten von Sperren und
Zugriffsrechten zusammenspielen: Pool-Verbindungen sollten nicht mit
zu weitreichenden Rechten ausgestattet sein; stattdessen empfiehlt
sich eine Anwendungsschicht, die pro Benutzerrolle entitätspräzise
Rechte (Proposition zur Berechtigungsgranularität,
Abschnitt~\ref{sec:wartbarkeit}) durchsetzt.

\section{Implementierung umfangreicher Anwendungen}

Für umfangreiche WWW-Anwendungen empfiehlt sich eine modulare
Architektur, in der -- wie in Kapitel~\ref{ch:hostsprachen}
beschrieben -- jedes Modul einer Relation des Tiefenpfads entspricht.
Diese Modulgrenzen erlauben es, einzelne Teile der Anwendung (z.\,B.
die Stammdatenverwaltung von $\textsc{Land}$ und $\textsc{Region}$)
unabhängig von den transaktionsintensiven Teilen
($\textsc{Auftrag}$, $\textsc{Auftragsposition}$) zu skalieren und zu
warten -- eine direkte praktische Konsequenz des in
Kapitel~\ref{ch:tiefenentwurf} entwickelten Paradigmas.

\label{ch:tiefenentwurf}

\label{sec:einleitung}
% ==============================================================

Der Entwurf relationaler Datenbankschemata ist eine der zentralen Aufgaben der
Softwareentwicklung und der angewandten Informatik. Obwohl die theoretischen
Grundlagen -- insbesondere Codd's relationales Modell \cite{Codd1970} sowie die
daraus abgeleiteten Normalformen \cite{Codd1974, Date2003} -- seit Jahrzehnten
bekannt sind, zeigen praktische Systeme immer wieder charakteristische
Entwurfsfehler. Einer der häufigsten und gleichzeitig folgenreichsten Fehler ist
die sogenannte \emph{sternförmige Schemaanordnung}: Eine zentrale
,,Hub\grqq{}-Relation referenziert direkt eine Vielzahl unabhängiger
,,Spoke\grqq{}-Relationen über Fremdschlüssel, ohne dass die natürliche semantische
Tiefe der Domäne abgebildet wird.

Dieser Beitrag stellt die These auf und belegt sie formal:
\begin{quote}
\emph{Ein tiefenorientierter Datenbankentwurf -- charakterisiert durch lineare
oder baumartige Fremdschlüsselketten -- ist dem sternförmigen Entwurf in allen
wesentlichen Gütekriterien der Datenbanktheorie überlegen.}
\end{quote}

Die Argumentation gliedert sich wie folgt: Abschnitt~\ref{sec:grundlagen}
etabliert die formalen Grundlagen. Abschnitt~\ref{sec:topologien} definiert
die Topologien ,,Tiefe\grqq{} und ,,Breite\grqq{} präzise.
Abschnitt~\ref{sec:normalformen} zeigt den Zusammenhang mit der
Normalisierungstheorie. Abschnitt~\ref{sec:anomalien} behandelt
Anomalievermeidung, Abschnitt~\ref{sec:abfragen} die Abfragekomplexität,
Abschnitt~\ref{sec:wartbarkeit} Wartbarkeit und Erweiterbarkeit.
Abschnitt~\ref{sec:beweis} liefert den konsolidierten formalen Nachweis.
Abschnitt~\ref{sec:beispiel} illustriert die Theorie an einem durchgängigen
Praxisbeispiel. Kapitel~\ref{ch:ausblick} fasst die Ergebnisse zusammen
und gibt einen ausführlichen Ausblick.

% ==============================================================
\section{Formale Grundlagen}
\label{sec:grundlagen}
% ==============================================================

\subsection{Relationales Modell}

\begin{definition}[Relation]
Eine \emph{Relation} $R$ über einem Attributschema $\att(R) = \{A_1, A_2, \ldots, A_n\}$
ist eine endliche Menge von Tupeln $t: \att(R) \to \bigcup_{i} \mathrm{dom}(A_i)$,
wobei $t(A_i) \in \mathrm{dom}(A_i)$ für alle $i$ \cite{Codd1970}.
\end{definition}

\begin{definition}[Datenbankschema]
Ein \emph{Datenbankschema} $\mathcal{S} = (\mathcal{R}, \mathcal{F}, \mathcal{K})$
besteht aus einer endlichen Menge von Relationenschemata $\mathcal{R} = \{R_1, \ldots, R_m\}$,
einer Menge funktionaler Abhängigkeiten $\mathcal{F}$ sowie einer Menge von
Integritätsbedingungen $\mathcal{K}$ (insbesondere Primär- und Fremdschlüsselbedingungen)
\cite{Ullman1988}.
\end{definition}

\begin{definition}[Funktionale Abhängigkeit]
Sei $R$ ein Relationenschema und $X, Y \subseteq \att(R)$.
Die \emph{funktionale Abhängigkeit} $X \to Y$ gilt in $R$, wenn für alle
Instanzen $r$ von $R$ und alle Tupel $t_1, t_2 \in r$ gilt:
$t_1(X) = t_2(X) \implies t_1(Y) = t_2(Y)$ \cite{Armstrong1974}.
\end{definition}

\begin{definition}[Fremdschlüssel]
Sei $R_i, R_j \in \mathcal{R}$ mit $K_j \subseteq \att(R_j)$ dem Primärschlüssel
von $R_j$. Eine Attributmenge $F \subseteq \att(R_i)$ heißt
\emph{Fremdschlüssel in $R_i$ referenzierend $R_j$}, wenn für alle Datenbankinstanzen
$d$ gilt: $\pi_F(r_i) \subseteq \pi_{K_j}(r_j)$ \cite{Date2003}.
\end{definition}

\subsection{Referenzgraph}

\begin{definition}[Referenzgraph]
\label{def:refgraph}
Sei $\mathcal{S}$ ein Datenbankschema. Der \emph{Referenzgraph}
$G(\mathcal{S}) = (V, E)$ hat als Knotenmenge $V = \mathcal{R}$ und als
Kantenmenge $E = \{(R_i, R_j) \mid \exists\, F \subseteq \att(R_i)\colon
F \text{ ist FK in } R_i \text{ referenzierend } R_j\}$.
\end{definition}

Der Referenzgraph ist das zentrale Werkzeug zur Unterscheidung von
Tiefenstruktur und Breitenstruktur eines Schemas.

\subsection{Komplexitätsmaße}

\begin{definition}[Schemakomplexitätsmaße]
\label{def:masse}
Sei $G(\mathcal{S}) = (V, E)$ der Referenzgraph eines Schemas $\mathcal{S}$.
\begin{enumerate}[label=(\roman*)]
  \item Der \emph{maximale Grad} $\Delta(G) = \max_{v \in V} \deg(v)$ misst
        die maximale Anzahl von Fremdschlüsselreferenzen eines Knotens.
  \item Die \emph{maximale Pfadlänge} $\lambda(G) = \max_{u,v \in V} d(u,v)$
        (Diameter) misst die längste Referenzkette im Schema.
  \item Das \emph{Tiefe-Breite-Verhältnis} $\tau(G) = \lambda(G) / \Delta(G)$
        charakterisiert die strukturelle Orientierung des Schemas.
\end{enumerate}
\end{definition}

Ein Schema heißt \emph{tiefenorientiert}, wenn $\tau(G) \gg 1$, und
\emph{breitenorientiert (sternförmig)}, wenn $\tau(G) \ll 1$.

% ==============================================================
\section{Topologien: Tiefe vs.\ Breite}
\label{sec:topologien}
% ==============================================================

\subsection{Formale Charakterisierung}

\begin{definition}[Sternschema]
\label{def:stern}
Ein Schema $\mathcal{S}$ heißt \emph{$k$-Sternschema}, wenn es eine
,,Hub\grqq{}-Relation $H \in \mathcal{R}$ gibt mit $\deg(H) = k$ und
$\lambda(G(\mathcal{S})) \leq 2$. Die $k$ von $H$ referenzierten Relationen
heißen \emph{Spoke}-Relationen.
\end{definition}

Das Sternschema ist insbesondere aus dem Data-Warehouse-Bereich bekannt
(Kimball-Methodik, \cite{Kimball1996}), wo es für OLAP-Abfragen optimiert
wird. Im OLTP-Kontext hingegen offenbart es gravierende Schwächen
(vgl.\ Abschnitt~\ref{sec:anomalien}).

\begin{definition}[Tiefenpfadschema]
\label{def:tiefenpfad}
Ein Schema $\mathcal{S}$ heißt \emph{Tiefenpfadschema der Länge $l$},
wenn $G(\mathcal{S})$ einen Pfad $R_1 \to R_2 \to \cdots \to R_l$
enthält mit $\Delta(G) \leq c$ für eine kleine Konstante $c$ (typisch $c \leq 3$)
und $\lambda(G(\mathcal{S})) = l - 1$.
\end{definition}

\begin{figure}[H]
\centering
\begin{tikzpicture}[
  node distance=1.4cm and 2.0cm,
  box/.style={rectangle, draw, rounded corners=3pt, minimum width=2.2cm,
              minimum height=0.8cm, font=\small\ttfamily, fill=blue!8},
  hub/.style={rectangle, draw, rounded corners=3pt, minimum width=2.5cm,
              minimum height=0.9cm, font=\small\ttfamily\bfseries, fill=orange!20},
  arr/.style={-{Stealth[length=7pt]}, thick},
  darr/.style={-{Stealth[length=7pt]}, thick, blue!70}
]

% === Sternschema (links) ===
\node[hub] (H) at (0,0) {Hub $H$};
\foreach \i/\name/\angle in {
  1/Spoke$_1$/90,
  2/Spoke$_2$/30,
  3/Spoke$_3$/-30,
  4/Spoke$_4$/-90,
  5/Spoke$_5$/150,
  6/Spoke$_6$/210
}{
  \node[box] (S\i) at ([shift={(\angle:2.8cm)}]H) {\name};
  \draw[arr, orange!60!black] (H) -- (S\i);
}
\node[above=0.15cm of H, font=\small\bfseries] {(a) Sternschema};

% === Tiefenpfadschema (rechts) ===
\begin{scope}[xshift=8.5cm]
  \node[box] (R1) at (0, 2.4)   {$R_1$};
  \node[box] (R2) at (0, 1.0)   {$R_2$};
  \node[box] (R3) at (0,-0.4)   {$R_3$};
  \node[box] (R4) at (0,-1.8)   {$R_4$};
  \node[box] (R5) at (0,-3.2)   {$R_5$};

  \draw[darr] (R1) -- (R2);
  \draw[darr] (R2) -- (R3);
  \draw[darr] (R3) -- (R4);
  \draw[darr] (R4) -- (R5);

  \node[above=0.15cm of R1, font=\small\bfseries] {(b) Tiefenpfadschema};
\end{scope}

\end{tikzpicture}
\caption{Gegenüberstellung von (a) Sternschema mit Hub-Relation $H$
und 6 Spoke-Relationen ($\Delta = 6$, $\lambda = 1$, $\tau = 0{,}17$) und
(b) Tiefenpfadschema mit linearer Kette
($\Delta = 1$, $\lambda = 4$, $\tau = 4$).}
\label{fig:topologien}
\end{figure}

\subsection{Semantische Motivation}

Die Wahl der Topologie ist kein rein ästhetisches, sondern ein
\emph{semantisch erzwungenes} Entwurfsprinzip. Reale Domänen weisen
fast immer hierarchische oder transitive Abhängigkeitsstrukturen auf:
Eine \texttt{Bestellung} hängt von einem \texttt{Kunden} ab, der wiederum
einer \texttt{Region} zugeordnet ist, die einem \texttt{Land} angehört.
Diese natürliche Tiefe der Abhängigkeiten erzwingt eine tiefe Pfadstruktur
und sollte nicht künstlich in eine flache Sternstruktur gezwungen werden
\cite{Date2003, Elmasri2015}.

% ==============================================================
\section{Normalisierungstheorie und Schematopologie}
\label{sec:normalformen}
% ==============================================================

\subsection{Normalformen und ihre Anforderungen}

Wir rekapitulieren die für den Nachweis relevanten Normalformen \cite{Codd1974, Fagin1977}:

\begin{definition}[Erste Normalform -- 1NF]
Eine Relation $R$ befindet sich in \emph{1NF}, wenn alle Attributdomänen
atomar sind, d.\,h.\ keine mengenwertigen oder strukturierten Attribute
enthält \cite{Codd1970}.
\end{definition}

\begin{definition}[Zweite Normalform -- 2NF]
Eine Relation $R$ in 1NF befindet sich in \emph{2NF}, wenn jedes
Nichtschlüsselattribut voll funktional abhängig vom Primärschlüssel ist,
d.\,h.\ keine partielle Abhängigkeit von einem echten Teilschlüssel existiert
\cite{Codd1974}.
\end{definition}

\begin{definition}[Dritte Normalform -- 3NF]
Eine Relation $R$ in 2NF befindet sich in \emph{3NF}, wenn kein
Nichtschlüsselattribut transitiv vom Primärschlüssel abhängig ist
\cite{Codd1974}.
\end{definition}

\begin{definition}[Boyce-Codd-Normalform -- BCNF]
Eine Relation $R$ befindet sich in \emph{BCNF}, wenn für jede
nicht-triviale funktionale Abhängigkeit $X \to A$ gilt: $X$ ist ein
Superschlüssel von $R$ \cite{Boyce1974}.
\end{definition}

\begin{definition}[Vierte Normalform -- 4NF]
Eine Relation $R$ in BCNF befindet sich in \emph{4NF}, wenn für jede
nicht-triviale mehrwertige Abhängigkeit $X \twoheadrightarrow Y$ gilt:
$X$ ist ein Superschlüssel von $R$ \cite{Fagin1977}.
\end{definition}

\subsection{Transitivität und Schematopologie}

\begin{theorem}[Transitivitätsverteilung]
\label{thm:transitivitaet}
Sei $R$ eine Relation mit $\att(R) = \{K, A_1, A_2, \ldots, A_n\}$ und
funktionalen Abhängigkeiten $K \to A_1 \to A_2 \to \cdots \to A_n$.
Dann verletzt $R$ die 3NF für alle $A_i$ mit $i \geq 2$.
Die einzige verlustfreie, abhängigkeitserhaltende Zerlegung in 3NF ist
die Kette $R_1(K, A_1),\; R_2(A_1, A_2),\; \ldots,\; R_{n-1}(A_{n-1}, A_n)$.
\end{theorem}

\begin{proof}
Sei $K \to A_i \to A_{i+1}$ für ein $i \geq 1$ eine transitive Kette.
Da $A_i$ kein Schlüssel ist, ist $A_{i+1}$ transitiv von $K$ abhängig und
$R$ verletzt 3NF. Die Zerlegung $R' = R \setminus \{A_{i+1}\}$ und
$R'' = (A_i, A_{i+1})$ ist verlustfrei nach dem Lossless-Join-Lemma
(da $A_i \to A_{i+1}$ gilt und $A_i = R' \cap R''$ sowie $A_i \to R''$ gilt)
\cite{Ullman1988}. Rekursive Anwendung für alle transitiven Abhängigkeiten
liefert die vollständige Zerlegungskette. Die Abhängigkeitserhaltung folgt
daraus, dass jede FD $A_i \to A_{i+1}$ in der entsprechenden Relation
$R_i(A_{i-1}, A_i, A_{i+1})$ bzw.\ $R_i(A_i, A_{i+1})$ präsent bleibt.
\end{proof}

\begin{corollary}[Tiefenpfad als 3NF-Normalform]
\label{cor:tiefenpfad3NF}
Ein Datenbankschema, das alle transitiven Abhängigkeiten einer Domäne
korrekt repräsentiert, hat zwingend einen tiefenorientierten Referenzgraphen.
Sternförmige Schemata, die dieselbe Domäne abbilden, verletzen die 3NF,
sofern transitive Abhängigkeiten in der Domäne existieren.
\end{corollary}

Korollar~\ref{cor:tiefenpfad3NF} ist das zentrale formale Ergebnis dieses
Abschnitts. Es besagt, dass die Normalisierungstheorie \emph{tiefe Pfade}
direkt erzwingt, wenn die Domäne transitive Abhängigkeiten aufweist --
was in praktisch allen nicht-trivialen realen Domänen der Fall ist.

\subsection{Mehrwertige Abhängigkeiten und 4NF}

\begin{proposition}[Mehrwertige Abhängigkeiten im Sternschema]
\label{prop:mvd_stern}
In einem $k$-Sternschema mit Hub $H$ und Spokes $S_1, \ldots, S_k$ gilt:
Falls zwei Spoke-Attribute $B_1 \in \att(S_1)$ und $B_2 \in \att(S_2)$
inhaltlich unabhängig sind und beide durch $K_H$ identifiziert werden, so
liegt in $H$ eine nicht-triviale mehrwertige Abhängigkeit $K_H \twoheadrightarrow B_1$
vor, wenn beide Attribute direkt in $H$ gehalten werden. Dies verletzt 4NF.
\end{proposition}

Die korrekte Auflösung ist -- erneut -- die Extraktion in separate Relationen,
was die Tiefenstruktur des Schemas weiter verstärkt.

% ==============================================================
\section{Anomalievermeidung}
\label{sec:anomalien}
% ==============================================================

\subsection{Klassifikation von Anomalien}

Die Datenbanktheorie unterscheidet drei klassische Anomalietypen \cite{Date2003}:

\begin{enumerate}[label=(\roman*)]
  \item \textbf{Einfügeanomalie}: Daten können nicht eingetragen werden,
        ohne gleichzeitig nicht vorhandene Daten erfinden zu müssen.
  \item \textbf{Löschanomalie}: Das Löschen eines Tupels vernichtet
        unbeabsichtigt weitere Informationen.
  \item \textbf{Änderungsanomalie}: Eine inhaltliche Änderung muss in
        mehreren Tupeln durchgeführt werden, da die Information redundant gespeichert ist.
\end{enumerate}

\begin{theorem}[Anomaliefreiheit durch Tiefenpfade]
\label{thm:anomaliefrei}
Sei $\mathcal{S}$ ein Datenbankschema in BCNF mit einem Tiefenpfadgraphen
$G(\mathcal{S})$. Dann ist $\mathcal{S}$ frei von Einfüge-, Lösch- und
Änderungsanomalien bzgl.\ der durch $G(\mathcal{S})$ modellierten
transitiven Abhängigkeiten.
\end{theorem}

\begin{proof}
BCNF garantiert, dass in jeder Relation $R_i$ der Kette jede nicht-triviale
FD $X \to A$ mit $X$ als Superschlüssel gilt. Somit ist jede Information
genau einmal gespeichert (keine Redundanz). Ohne Redundanz:
\begin{itemize}
  \item \emph{Änderungsanomalie}: Entfällt, da jede Eigenschaft in genau einer
        Relation steht. Eine Änderung betrifft genau ein Tupel.
  \item \emph{Löschanomalie}: Da eine Entität $e_i$ in $R_i$ unabhängig von
        ihrer Referenz durch $R_{i-1}$ existiert (referentielle Integrität
        via FK), gehen beim Löschen eines $R_{i-1}$-Tupels keine $R_i$-Daten
        verloren (ON DELETE SET NULL / RESTRICT).
  \item \emph{Einfügeanomalie}: Eine neue Entität in $R_i$ kann unabhängig
        von $R_{i-1}$ eingefügt werden, da FK-Referenz nur in der
        referenzierenden Relation besteht.
\end{itemize}
Im Sternschema hingegen erzeugt die direkte Kodierung transitiver Attribute
in $H$ Redundanz, da Attributwerte von Spoke-Entitäten wiederholt in $H$
gespeichert sind -- dies ermöglicht alle drei Anomalietypen.
\end{proof}

\subsection{Quantitative Analyse der Redundanz}

\begin{proposition}[Redundanzfaktor im Sternschema]
\label{prop:redundanz}
Sei $H$ eine Hub-Relation mit $n_H$ Tupeln und Spoke-Relation $S_1$ mit
$n_{S_1}$ Tupeln und $|A_{S_1}|$ Attributen. Werden die Attribute von $S_1$
direkt in $H$ denormalisiert, so beträgt der Redundanzfaktor:
\[
  \rho = \frac{n_H}{n_{S_1}} \cdot |A_{S_1}|
\]
Für $n_H \gg n_{S_1}$ (typisch bei 1:n-Beziehungen) gilt $\rho \gg 1$.
\end{proposition}

In einer typischen Bestelldatenbank mit $n_H = 10^6$ Bestellungen und
$n_{S_1} = 10^4$ Kunden ergibt sich $\rho = 100 \cdot |A_{\text{Kunde}}|$.
Bei 10 Kundenattributen werden also 1\,000-fach mehr Daten gespeichert
als notwendig -- ein massiver Speicher- und Konsistenzaufwand.

% ==============================================================
\section{Abfragekomplexität und Optimierbarkeit}
\label{sec:abfragen}
% ==============================================================

\subsection{Join-Komplexität}

Ein häufiges Argument gegen tiefe Pfade ist der höhere Join-Aufwand.
Wir zeigen, dass dieses Argument zwar formal korrekt, aber praktisch
irrelevant ist, wenn moderne Datenbankoptimierung berücksichtigt wird.

\begin{definition}[Join-Tiefe]
Die \emph{Join-Tiefe} $\delta(\mathcal{Q})$ einer Abfrage $\mathcal{Q}$ über
Schema $\mathcal{S}$ ist die Anzahl der JOIN-Operationen in der optimalen
logischen Ausführung von $\mathcal{Q}$.
\end{definition}

\begin{proposition}[Join-Tiefe und Pfadlänge]
\label{prop:jointiefe}
Für eine vollständige Traversal-Abfrage über einen Tiefenpfad
$R_1 \to R_2 \to \cdots \to R_l$ gilt $\delta = l - 1$.
Für ein $k$-Sternschema gilt für dieselbe semantische Abfrage $\delta = k$.
Da $l \ll k$ in tiefen Schemata (wenige, längere Ketten statt vieler
direkter Abhängigkeiten), ist die Join-Tiefe im Tiefenpfad geringer
für partielle Traversals \cite{Ramakrishnan2003}.
\end{proposition}

\begin{remark}
Selbst wenn $l > k$ für globale Traversals gilt, kompensiert dies der
Abfrageoptimierer durch:
\begin{itemize}
  \item \emph{Index-Only Scans} auf FK-Spalten (eliminiert Heap-Zugriffe),
  \item \emph{Hash-Joins} mit amortisierter $O(n)$-Komplexität \cite{Graefe1993},
  \item \emph{Materialized Views} für häufige tiefe Traversals,
  \item \emph{Predicate Pushdown} entlang des Pfades.
\end{itemize}
\end{remark}

\subsection{Selektivität und Filtereffizienz}

\begin{theorem}[Selektivitätsvorteil tiefer Pfade]
\label{thm:selektivitaet}
Sei $\mathcal{Q}$ eine Abfrage mit Filterprädikat $\sigma_P$ auf Attributen
einer Entität $E$ der Tiefe $d$ im Pfad. Im Tiefenpfadschema kann $\sigma_P$
direkt auf der Relation $R_d$ angewandt werden (Predicate Pushdown).
Im Sternschema hingegen sind transitive Attribute von $E$ in der Hub-Relation
$H$ de\-normalisiert, was bedeutet, dass der Filter auf der gesamten,
$n_H$-tupligen Hub-Relation angewandt werden muss.
\[
  \text{Kosten}_{\text{Stern}}(\sigma_P) = O(n_H) \cdot c_{\text{IO}}
  \qquad\text{vs.}\qquad
  \text{Kosten}_{\text{Tief}}(\sigma_P) = O(n_{R_d}) \cdot c_{\text{IO}}
\]
Da $n_H \geq n_{R_d}$ für alle $d$, ist der Tiefenpfad mindestens so
effizient wie das Sternschema, typisch deutlich effizienter.
\end{theorem}

\begin{proof}
Im normalisierten Tiefenpfadschema liegen die Attribute von Entität $E$ in
$R_d$ mit $|R_d| = n_{R_d}$ Tupeln. Ein Index auf dem Primärschlüssel von
$R_d$ ermöglicht $O(\log n_{R_d} + k')$ für $k'$ Ergebnistupel.
Im Sternschema befinden sich diese Attribute in $H$ vermischt mit allen anderen
Attributen; ohne Partial Index muss der gesamte Heap von $H$ gescannt werden
($O(n_H)$). Da eine 1:n-Beziehung $|E| \leq n_H$ impliziert, gilt
$n_{R_d} \leq n_H$, und die Ungleichung ist bewiesen.
\end{proof}

\subsection{Schreibkomplexität und Konsistenzerhaltung}

\begin{proposition}[Schreibkostenvorteil]
Im Tiefenpfadschema erfordert eine Änderung an Attributen der Entität $E$
genau \emph{eine} UPDATE-Operation auf $R_d$. Im denormalisierten
Sternschema müssen alle Tupel in $H$, die auf diese Entität verweisen,
aktualisiert werden: im Worst Case $O(n_H / n_{R_d})$ Operationen --
proportional zum Redundanzfaktor $\rho$ aus Proposition~\ref{prop:redundanz}.
\end{proposition}

% ==============================================================
\section{Wartbarkeit und Erweiterbarkeit}
\label{sec:wartbarkeit}
% ==============================================================

\subsection{Schemaevolution}

Ein zentrales Qualitätsmerkmal eines Datenbankschemas ist seine
Fähigkeit, sich an veränderte Anforderungen anzupassen -- bekannt als
\emph{Schemaevolution} \cite{Roddick1995}.

\begin{theorem}[Lokale Änderbarkeit im Tiefenpfad]
\label{thm:schaemaevolution}
Im Tiefenpfadschema $R_1 \to R_2 \to \cdots \to R_l$ ist das Hinzufügen,
Entfernen oder Ändern von Attributen der Entität $R_i$ eine
\emph{lokale Schemaoperation}, die ausschließlich $R_i$ betrifft.
Im $k$-Sternschema mit denormalisierten Attributen in $H$ ist dieselbe
Operation \emph{global} und betrifft $H$ sowie potenziell alle
abhängigen Abfragen, Views und Applikationsschichten.
\end{theorem}

\begin{proof}
Da im Tiefenpfadschema jedes Attribut in genau einer Relation vorkommt
(BCNF-Eigenschaft), hat eine Schemaänderung an $R_i$ keinen direkten
Einfluss auf $R_j$ mit $j \neq i$. Nur FK-Referenzen
$R_{i-1} \to R_i$ sind betroffen -- diese beschränken sich auf eine
einzige Kante im Referenzgraphen. Im Sternschema sind Attribute aus
multiple Entitäten in $H$ gemischt; eine Änderung an $H$ betrifft
alle darauf basierenden Abfragen, was zu einem kaskadierten
Änderungsaufwand führt \cite{Ambler2003}.
\end{proof}

\subsection{Zugriffsrechteverwaltung}

Ein weiterer Vorteil tiefer Pfade zeigt sich bei der Granularität der
Zugriffsrechteverwaltung:

\begin{proposition}[Berechtigungsgranularität]
Im Tiefenpfadschema können Zugriffsrechte \emph{entitätspräzise} vergeben
werden: Ein Nutzer erhält SELECT-Rechte auf $R_d$, ohne Zugriff auf andere
Relationen zu erhalten. Im Sternschema mit Hub $H$ bedeutet Zugriff auf
$H$ stets Zugriff auf alle in $H$ denormalisierten Entitätsdaten --
eine grobe Granularität, die dem Prinzip der minimalen Rechtevergabe
(Principle of Least Privilege) widerspricht \cite{Saltzer1975}.
\end{proposition}

\subsection{Testbarkeit und Datenqualität}

Tiefenpfadschemata sind inhärent leichter zu testen und zu validieren,
da:
\begin{itemize}
  \item Jede Relation eine \emph{klar umgrenzte semantische Einheit} darstellt
        (Single Responsibility für Relationen),
  \item Integritätsbedingungen pro Relation formuliert und geprüft werden,
  \item Fremdschlüsselbedingungen entlang des Pfades eine implizite
        Validierungskaskade bilden.
\end{itemize}

% ==============================================================
\section{Konsolidierter formaler Nachweis}
\label{sec:beweis}
% ==============================================================

\subsection{Gütekriterien}

Wir formalisieren nun den umfassenden Vergleich anhand eines Vektors von
Gütekriterien:

\begin{definition}[Gütevektor]
Sei $\mathcal{S}$ ein Datenbankschema. Der \emph{Gütevektor}
$\mathbf{q}(\mathcal{S}) \in \mathbb{R}^6$ ist definiert als:
\[
  \mathbf{q}(\mathcal{S}) = \bigl(
    q_{\NF},\; q_{\text{Red}},\; q_{\text{Anom}},\;
    q_{\text{Sch}},\; q_{\text{Wart}},\; q_{\text{Sek}}
  \bigr)
\]
wobei die Komponenten Normalformgrad, Redundanzfreiheit,
Anomaliefreiheit, Schreibkosteneffizienz, Wartbarkeit und
Selektivitätseffizienz kodieren (alle normiert auf $[0,1]$,
höher ist besser).
\end{definition}

\begin{theorem}[Ãœberlegenheit des Tiefenpfadschemas]
\label{thm:hauptsatz}
Sei $\mathcal{S}_{\text{Tief}}$ ein Tiefenpfadschema in BCNF und
$\mathcal{S}_{\text{Stern}}$ ein semantisch äquivalentes $k$-Sternschema
($k \geq 3$) mit denormalisierten transitiven Attributen. Dann gilt:
\[
  \mathbf{q}(\mathcal{S}_{\text{Tief}}) \succ \mathbf{q}(\mathcal{S}_{\text{Stern}})
\]
komponentenweise, d.\,h.\ $q_j(\mathcal{S}_{\text{Tief}}) \geq q_j(\mathcal{S}_{\text{Stern}})$
für alle $j$, mit strikter Ungleichung für mindestens vier Komponenten.
\end{theorem}

\begin{proof}
Wir zeigen die strikten Ungleichungen komponentenweise:

\medskip
(i) \textbf{Normalformgrad} $q_{\NF}$: Nach Korollar~\ref{cor:tiefenpfad3NF}
befindet sich $\mathcal{S}_{\text{Tief}}$ in BCNF, während
$\mathcal{S}_{\text{Stern}}$ transitive Abhängigkeiten in $H$ enthält
und damit 3NF verletzt. Somit $q_{\NF}(\mathcal{S}_{\text{Tief}}) > q_{\NF}(\mathcal{S}_{\text{Stern}})$.

(ii) \textbf{Redundanzfreiheit} $q_{\text{Red}}$: Nach Proposition~\ref{prop:redundanz}
ist $\rho(\mathcal{S}_{\text{Stern}}) \gg 1$, während
$\rho(\mathcal{S}_{\text{Tief}}) = 1$ (jede Information einmalig gespeichert).

(iii) \textbf{Anomaliefreiheit} $q_{\text{Anom}}$: Folgt direkt aus
Satz~\ref{thm:anomaliefrei}.

(iv) \textbf{Schreibkosteneffizienz} $q_{\text{Sch}}$: Aus der Redundanz
$\rho > 1$ im Sternschema folgt ein entsprechend höherer Schreibaufwand.

(v) \textbf{Wartbarkeit} $q_{\text{Wart}}$: Nach Satz~\ref{thm:schaemaevolution}
sind Schemaänderungen im Tiefenpfad lokal beschränkt.

(vi) \textbf{Selektivitätseffizienz} $q_{\text{Sek}}$: Nach
Satz~\ref{thm:selektivitaet} gilt stets
$\text{Kosten}_{\text{Tief}} \leq \text{Kosten}_{\text{Stern}}$.

Da für (i)--(v) strikte Ungleichungen gelten, ist die Behauptung bewiesen.
\end{proof}

% ==============================================================
\section{Graphentheoretische Fundierung: DFS, BFS und die Asymmetrie der Suchalgorithmen}
\label{sec:dfs_bfs}
% ==============================================================

Die bisherigen Argumente stützten sich auf die Normalformtheorie und die
Anomalielehre. Es gibt jedoch eine tiefere, algorithmische Begründung für
die Ãœberlegenheit tiefer Strukturen, die in der Theorie der Graphentraversierung
verwurzelt ist. Die Beobachtung ist fundamental:

\begin{quote}
\emph{Tiefensuche (DFS) und Breitensuche (BFS) über den Referenzgraphen eines
Datenbankschemas sind \textbf{nicht symmetrisch}. Ihre Ergebnisstrukturen --
und damit die induzierten Schemastrukturen -- unterscheiden sich qualitativ
und in ihrer Komplexität grundlegend.}
\end{quote}

\subsection{Tiefensuche und der DFS-Wald}

\begin{definition}[DFS-Baum und DFS-Wald, nach Cormen et al.~\cite{Cormen2009}]
\label{def:dfswald}
Sei $G = (V, E)$ ein gerichteter Graph. Die \emph{Tiefensuche} (Depth-First
Search, DFS) traversiert $G$, indem sie von jedem besuchten Knoten $u$ aus
rekursiv in die Nachfolger abtaucht, bevor sie zum Ausgangspunkt zurückkehrt.
Die dabei entstehenden Baumkanten bilden einen \emph{DFS-Wald}
$T_{\mathrm{DFS}} = (V, E_{\mathrm{tree}})$ mit
$E_{\mathrm{tree}} \subseteq E$. Jede Zusammenhangskomponente des DFS-Walds
ist ein \emph{DFS-Baum}.
\end{definition}

Der DFS-Wald hat bemerkenswerte Eigenschaften, die ihn für den
Datenbankschemaentwurf auszeichnen:

\begin{theorem}[Struktureigenschaften des DFS-Walds, \cite{Cormen2009, Tarjan1972}]
\label{thm:dfswald}
Sei $T_{\mathrm{DFS}}$ der DFS-Wald über dem Referenzgraphen
$G(\mathcal{S})$. Dann gilt:
\begin{enumerate}[label=(\roman*)]
  \item $T_{\mathrm{DFS}}$ ist ein \emph{Wald} (azyklischer Graph, Menge
        von Bäumen) -- also eine azyklische, hierarchische Struktur.
  \item Die Tiefe $\depth(T_{\mathrm{DFS}})$ ist maximal unter allen
        aufspannenden Wäldern von $G(\mathcal{S})$.
  \item Rückwärtskanten (back edges) $e \notin E_{\mathrm{tree}}$ identifizieren
        Zyklen und damit potenzielle \emph{zirkuläre Abhängigkeiten} im Schema.
  \item Die \emph{topologische Sortierung} des Schemas -- notwendig für
        eine anomaliefreie Einfüge- und Löschreihenfolge -- ist
        genau die DFS-Postorder-Reihenfolge.
\end{enumerate}
\end{theorem}

\begin{corollary}[DFS-Wald als natürliches Schemaabbild]
Der DFS-Wald $T_{\mathrm{DFS}}$ über $G(\mathcal{S})$ ist das
\emph{natürliche strukturelle Abbild} eines korrekt normalisierten
Tiefenpfadschemas. Die Baumkanten entsprechen Fremdschlüsselbeziehungen,
die Baumtiefe der semantischen Abhängigkeitskette.
Jedes BCNF-Schema erzeugt -- wenn als Graph betrachtet -- einen azyklischen,
tiefenorientierten Referenzgraphen, der einem DFS-Wald isomorph ist.
\end{corollary}

\subsection{Breitensuche und das Fehlen eines BFS-Walds}

Die Breitensuche (Breadth-First Search, BFS) traversiert den Graphen
schichtweise: Alle Knoten der Distanz $d$ werden vollständig besucht, bevor
Knoten der Distanz $d+1$ betreten werden.

\begin{definition}[BFS-Baum, nach Cormen et al.~\cite{Cormen2009}]
\label{def:bfsbaum}
Die Breitensuche über $G = (V, E)$ ausgehend von einem Startknoten $s$
erzeugt einen \emph{BFS-Baum} $T_{\mathrm{BFS}}$, der die kürzesten Pfade
von $s$ zu allen erreichbaren Knoten enthält. Die Breite des BFS-Baums auf
Tiefe $d$ ist $b_d = |\{v \mid d_G(s,v) = d\}|$.
\end{definition}

\begin{theorem}[BFS erzeugt keinen Wald]
\label{thm:bfskeinwald}
Die Breitensuche über einem unzusammenhängenden Graphen $G$ mit $k$
Zusammenhangskomponenten $C_1, \ldots, C_k$ erzeugt im Allgemeinen
\emph{keinen Wald}, sondern eine Menge von $k$ BFS-Bäumen
$T_{\mathrm{BFS}}^{(1)}, \ldots, T_{\mathrm{BFS}}^{(k)}$, die jeweils
\emph{sternförmig} um ihren Startknoten $s_i \in C_i$ angeordnet sind.
Insbesondere:
\begin{enumerate}[label=(\roman*)]
  \item Für einen $k$-Stern\-graph $G_\star$ mit Hub $H$ und Knoten
        $S_1, \ldots, S_k$ ist der BFS-Baum mit Start $H$ identisch
        mit $G_\star$ selbst: $T_{\mathrm{BFS}} = G_\star$.
  \item Die BFS-Tiefe beträgt $\depth(T_{\mathrm{BFS}}) = 1$, unabhängig
        von $k$.
  \item Die BFS über einem zusammenhängenden Graphen mit $n$ Knoten
        erzeugt genau \emph{einen} Baum, nicht einen Wald -- der Begriff
        ,,BFS-Wald\grqq{} ist daher für zusammenhängende Graphen
        nicht definiert.
\end{enumerate}
\end{theorem}

\begin{proof}
(i) Beim Startsknoten $H$ werden in der ersten BFS-Schicht alle direkten
Nachbarn $S_1, \ldots, S_k$ besucht. Da diese keine weiteren Nachbarn
besitzen (Stern-Definition), terminiert BFS nach Schicht 1.
Die Baumkanten sind genau $\{(H, S_i) \mid i = 1, \ldots, k\}$, was dem
Stern entspricht.

(ii) $\depth(T_{\mathrm{BFS}}) = \max_{v} d_{T}(s, v) = 1$.

(iii) BFS ist eine \emph{Single-Source}-Traversierung; bei zusammenhängendem
$G$ wird genau ein Baum mit Wurzel $s$ erzeugt. Mehrere Wurzeln (Wald)
entstehen nur durch Multi-Source-BFS, was der Standard-BFS-Definition
widerspricht. Im Gegensatz dazu erzeugt DFS über einem unzusammenhängenden
Graphen automatisch einen Wald, indem für jede noch unbesuchte Zusammenhangskomponente
eine neue DFS-Traversierung gestartet wird.
\end{proof}

\subsection{Die fundamentale Asymmetrie von DFS und BFS}

\begin{theorem}[Asymmetrie von DFS und BFS]
\label{thm:asymmetrie}
DFS und BFS sind \emph{algorithmisch asymmetrisch} in folgenden Dimensionen:
\begin{enumerate}[label=(\roman*)]
  \item \textbf{Ergebnisstruktur}: DFS erzeugt stets einen \emph{Wald}
        (Menge von Bäumen mit maximaler Tiefe), BFS erzeugt einen oder
        mehrere \emph{flache Bäume} (Tiefe $\leq$ Durchmesser der
        Zusammenhangskomponente, typisch minimal).
  \item \textbf{Komplexität der Ergebnisstruktur}: Der DFS-Wald über
        einem Graphen $G = (V, E)$ hat Tiefe $\Theta(|V|)$ im
        schlechtesten Fall (Pfadgraph), der BFS-Baum hat Tiefe
        $O(\log |V|)$ bei ausgeglichenen Graphen -- eine fundamentale
        Größenordnungsdifferenz.
  \item \textbf{Kantenklassifikation}: DFS unterscheidet vier Kantentypen
        (Baumkanten, Rückwärtskanten, Vorwärtskanten, Querkanten), was eine
        reichere Strukturanalyse erlaubt \cite{Cormen2009}. BFS unterscheidet
        nur Baumkanten und Querkanten.
  \item \textbf{Informationsgehalt}: Die DFS-Stempel (Entdeckungs- und
        Abschlusszeiten) enkodieren die vollständige topologische Struktur
        des Graphen. BFS-Distanzen enkodieren nur metrische, nicht aber
        hierarchische Information.
  \item \textbf{Waldeigenschaft}: DFS erzeugt für beliebige Graphen
        automatisch einen Wald -- dies ist eine inhärente Eigenschaft der
        Rekursionsstruktur. BFS besitzt diese Eigenschaft nicht.
\end{enumerate}
\end{theorem}

\begin{figure}[H]
\centering
\begin{tikzpicture}[
  node distance=1.1cm and 1.5cm,
  box/.style={circle, draw, minimum size=0.75cm, font=\small\ttfamily,
              fill=blue!10},
  hub/.style={circle, draw, minimum size=0.85cm, font=\small\ttfamily\bfseries,
              fill=orange!25},
  treearr/.style={-{Stealth[length=6pt]}, thick, blue!70},
  starr/.style={-{Stealth[length=6pt]}, thick, orange!70!black},
  dashed/.style={-{Stealth[length=6pt]}, dashed, gray}
]

%% ---- DFS-Wald (links) ----
\begin{scope}[xshift=0cm]
  \node[box] (A) at (0, 0)    {$A$};
  \node[box] (B) at (0,-1.4)  {$B$};
  \node[box] (C) at (0,-2.8)  {$C$};
  \node[box] (D) at (0,-4.2)  {$D$};
  \node[box] (E) at (0,-5.6)  {$E$};

  \draw[treearr] (A) -- (B);
  \draw[treearr] (B) -- (C);
  \draw[treearr] (C) -- (D);
  \draw[treearr] (D) -- (E);

  \node[above=0.3cm of A, align=center, font=\small\bfseries]
    {DFS-Wald\\(Tiefenpfad)};
  \node[below=0.2cm of E, font=\scriptsize\color{blue!60}]
    {Tiefe $= 4$, $\Delta = 1$};
\end{scope}

%% ---- BFS-Baum / Sternschema (rechts) ----
\begin{scope}[xshift=6.5cm]
  \node[hub] (H) at (0, -2.8) {$H$};
  \foreach \i/\angle in {1/90, 2/30, 3/-30, 4/-90, 5/150, 6/210}{
    \node[box] (N\i) at ([shift={(\angle:2.4cm)}]H) {$S_{\i}$};
    \draw[starr] (H) -- (N\i);
  }
  \node[above=2.5cm of H, align=center, font=\small\bfseries]
    {BFS-Baum\\(Sternschema)};
  \node[below=2.7cm of H, font=\scriptsize\color{orange!70!black}]
    {Tiefe $= 1$, $\Delta = 6$};
\end{scope}

\end{tikzpicture}
\caption{Gegenüberstellung der durch DFS induzierten Tiefenpfadstruktur
(links, Tiefe $= 4$, $\Delta = 1$) und des durch BFS induzierten
Sternschemas (rechts, Tiefe $= 1$, $\Delta = 6$).
Die Asymmetrie ist offensichtlich: DFS erzeugt einen Wald mit maximaler
Tiefe und minimaler Breite; BFS erzeugt eine flache, sternförmige Struktur.}
\label{fig:dfs_bfs}
\end{figure}

\subsection{Datenbanktheoretische Konsequenz der Asymmetrie}

\begin{corollary}[Datenbanktheoretische Konsequenz]
\label{cor:konsequenz_asymmetrie}
Die algorithmische Asymmetrie von DFS und BFS hat eine direkte
datenbanktheoretische Entsprechung:

\begin{enumerate}[label=(\roman*)]
  \item Der \emph{DFS-Wald} über dem Referenzgraphen eines Schemas ist
        das strukturelle Abbild des \emph{normalisierten Tiefenpfadschemas}.
        Er existiert für beliebige Graphen automatisch und ist azyklisch,
        hierarchisch und eindeutig (bis auf Reihenfolge).
  \item Der \emph{BFS-Baum} über dem Referenzgraphen eines Schemas entspricht
        dem \emph{sternförmigen, denormalisierten Schema}. Er existiert nur
        für Single-Source-Traversierungen und kodiert keine hierarchische
        Tiefenstruktur.
  \item Da DFS automatisch einen \emph{Wald} erzeugt, BFS hingegen
        \emph{keinen Wald}, ist die Ãœberlegenheit des Tiefenpfadschemas
        nicht nur normalisierungstheoretisch, sondern auch algorithmisch
        fundamental: Der DFS-Wald ist die \emph{kanonische, strukturell
        reichere} Darstellung jedes Referenzgraphen.
  \item Die Komplexitätsdifferenz ist qualitativ, nicht nur quantitativ:
        DFS liefert topologische Sortierbarkeit, Zykeldetektierbarkeit
        und vollständige hierarchische Information; BFS liefert keine
        dieser Eigenschaften.
\end{enumerate}
\end{corollary}

\begin{remark}[Praktische Implikation]
Wenn ein Datenbankarchitekt ein Schema entwirft, vollzieht er -- bewusst
oder unbewusst -- eine Traversierung des semantischen Abhängigkeitsgraphen
der Domäne. Wählt er implizit eine DFS-Strategie (erst in die Tiefe einer
Entitätshierarchie, dann zur nächsten), entsteht ein tiefes, normalisiertes
Schema. Wählt er eine BFS-Strategie (alle direkten Attribute einer zentralen
Entität sammeln, ohne in Hierarchien abzusteigen), entsteht ein flaches,
fehleranfälliges Sternschema. Die Asymmetrie der Algorithmen \emph{spiegelt
sich direkt in der Qualität des resultierenden Schemas wider}.
\end{remark}

% ==============================================================
\section{Durchgängiges Praxisbeispiel}
\label{sec:beispiel}
% ==============================================================

\subsection{Domäne: Auftragsabwicklung}

Wir betrachten eine klassische Auftragsabwicklungs-Domäne mit folgenden
Entitäten und natürlichen transitiven Abhängigkeiten:

\[
  \texttt{Bestellung} \xrightarrow{\text{FK}} \texttt{Kunde}
  \xrightarrow{\text{FK}} \texttt{Ort}
  \xrightarrow{\text{FK}} \texttt{Region}
  \xrightarrow{\text{FK}} \texttt{Land}
\]

\subsection{Sternschema (anti-pattern)}

\begin{lstlisting}[caption={Anti-Pattern: Sternschema mit denormalisierten Attributen in Hub-Relation},label=lst:stern]
-- Hub-Relation: enthaelt direkt alle abhaengigen Attribute
CREATE TABLE Bestellung (
    BestellNr     INT PRIMARY KEY,
    Datum         DATE,
    Betrag        DECIMAL(10,2),
    -- Kundenattribute (transitive Abhaengigkeit Stufe 1)
    KundeNr       INT,
    KundeName     VARCHAR(100),
    KundeEmail    VARCHAR(150),
    -- Ortattribute (transitive Abhaengigkeit Stufe 2)
    OrtID         INT,
    OrtName       VARCHAR(80),
    PLZ           CHAR(5),
    -- Regionattribute (transitive Abhaengigkeit Stufe 3)
    RegionID      INT,
    RegionName    VARCHAR(80),
    -- Landattribute (transitive Abhaengigkeit Stufe 4)
    LandID        INT,
    LandName      VARCHAR(60),
    ISO_Code      CHAR(2)
    -- Verletzt 2NF, 3NF und BCNF!
);
\end{lstlisting}

Probleme: Aendert sich der Name eines Kunden, muessen alle seine
Bestellungstupel aktualisiert werden. Existiert ein Kunde noch ohne
Bestellung, kann er nicht eingetragen werden (Einfuegeanomalie).

\subsection{Tiefenpfadschema (BCNF-konform)}

\begin{lstlisting}[caption={Korrekte Tiefenpfadstruktur: BCNF-konformes Schema},label=lst:tief]
CREATE TABLE Land (
    LandID    INT PRIMARY KEY,
    LandName  VARCHAR(60) NOT NULL,
    ISO_Code  CHAR(2)     NOT NULL UNIQUE
);

CREATE TABLE Region (
    RegionID   INT PRIMARY KEY,
    RegionName VARCHAR(80) NOT NULL,
    LandID     INT NOT NULL,
    FOREIGN KEY (LandID) REFERENCES Land(LandID)
);

CREATE TABLE Ort (
    OrtID    INT PRIMARY KEY,
    OrtName  VARCHAR(80) NOT NULL,
    PLZ      CHAR(5),
    RegionID INT NOT NULL,
    FOREIGN KEY (RegionID) REFERENCES Region(RegionID)
);

CREATE TABLE Kunde (
    KundeNr   INT PRIMARY KEY,
    KundeName VARCHAR(100) NOT NULL,
    Email     VARCHAR(150) UNIQUE,
    OrtID     INT NOT NULL,
    FOREIGN KEY (OrtID) REFERENCES Ort(OrtID)
);

CREATE TABLE Bestellung (
    BestellNr INT PRIMARY KEY,
    Datum     DATE         NOT NULL,
    Betrag    DECIMAL(10,2) NOT NULL,
    KundeNr   INT          NOT NULL,
    FOREIGN KEY (KundeNr) REFERENCES Kunde(KundeNr)
);
-- Pfadtiefe: 4 (Bestellung->Kunde->Ort->Region->Land)
-- Jede Relation in BCNF. Keine Anomalien moeglich.
\end{lstlisting}

Der Referenzgraph dieses Schemas ist ein reiner Pfad der Länge 4 --
exakt die durch den DFS-Wald induzierte Tiefenstruktur gemäß
Satz~\ref{thm:dfswald}.

\subsection{Vergleichstabelle}

\begin{table}[H]
\centering
\renewcommand{\arraystretch}{1.35}
\begin{tabular}{lcc}
\toprule
\textbf{Kriterium} & \textbf{Sternschema} & \textbf{Tiefenpfad (BCNF)} \\
\midrule
Normalformgrad            & $<$ 3NF   & BCNF / 4NF \\
Redundanzfaktor $\rho$    & $\gg 1$   & $= 1$ \\
Anomaliefreiheit          & nein      & ja \\
Schreibaufwand (Update)   & $O(n_H)$  & $O(1)$ \\
Suchkosteneffizienz       & $O(n_H)$  & $O(n_{R_d})$ \\
Schemaevolution           & global    & lokal \\
Berechtigungsgranularität & grob      & fein \\
DFS-Wald-Isomorphismus    & nein      & ja \\
Topolog.\ Sortierbarkeit  & begrenzt  & vollständig \\
Zyklusfreiheit            & ungeprüft & inhärent \\
\bottomrule
\end{tabular}
\caption{Zusammenfassender Vergleich Sternschema vs.\ Tiefenpfadschema
nach den formalen Kriterien dieser Arbeit.}
\label{tab:vergleich}
\end{table}

% ==============================================================

\chapter{Fazit und Ausblick}
\label{ch:ausblick}

\section{Zusammenfassung der Ergebnisse}
\label{sec:fazit}

Diese Arbeit hat formal nachgewiesen, dass der tiefenorientierte
Datenbankentwurf dem sternförmigen Entwurf in allen wesentlichen
Gütekriterien überlegen ist. Die Argumentation verlief auf drei
komplementären Ebenen:

\textbf{Normalisierungstheoretisch} (Abschnitte~\ref{sec:normalformen}
und~\ref{sec:anomalien} in Kapitel~\ref{ch:tiefenentwurf}, vertieft in
Kapitel~\ref{ch:entwurf}): Transitive Abhängigkeiten erzwingen tiefe
Pfade, da ihre korrekte Auflösung gemäß
Theorem~\ref{thm:transitivitaet} genau eine lineare Zerlegungskette
liefert, die nach Theorem~\ref{thm:lossless} verlustfrei und nach
Definition~\ref{def:dep-erhaltend} abhängigkeitserhaltend ist.
Sternförmige Schemata verletzen zwingend die 3NF, sobald transitive
Abhängigkeiten in der Domäne existieren (Korollar~\ref{cor:tiefenpfad3NF}).

\textbf{Algorithmisch} (Abschnitt~\ref{sec:dfs_bfs}): Die Tiefensuche
(DFS) und die Breitensuche (BFS) sind \emph{fundamental asymmetrisch}.
DFS erzeugt stets einen Wald -- eine hierarchische, azyklische,
topologisch sortierbare Struktur mit maximaler Tiefe
(Theorem~\ref{thm:dfswald}). BFS erzeugt keinen Wald, sondern flache,
sternförmige Bäume ohne hierarchische Tiefenstruktur
(Theorem~\ref{thm:bfskeinwald}). Diese algorithmische Asymmetrie
(Theorem~\ref{thm:asymmetrie}) spiegelt sich direkt in der
Schemaqualität wider: Der DFS-Wald ist die kanonische Darstellung eines
korrekt normalisierten Schemas, der BFS-Baum die Darstellung eines
denormalisierten Sternschemas (Korollar~\ref{cor:konsequenz_asymmetrie}).

\textbf{Praktisch} (Abschnitte~\ref{sec:abfragen},~\ref{sec:wartbarkeit}
und~\ref{sec:beispiel} in Kapitel~\ref{ch:tiefenentwurf}, illustriert
in den Kapiteln~\ref{ch:sql},~\ref{ch:betrieb},~\ref{ch:plsql},
\ref{ch:hostsprachen} und~\ref{ch:www}): Tiefenpfadschemata erzielen
überlegene Selektivität (Theorem~\ref{thm:selektivitaet}), geringeren
Schreibaufwand, feingranulare Zugriffsrechte
(Abschnitt~\ref{sec:zugriffsrechte}), günstigere
Sperrgranularität (Bemerkung~\ref{rem:sperren}) und bessere
Wartbarkeit (Theorem~\ref{thm:schaemaevolution}). Ergänzend wurde mit
\texttt{DISTINCT ON} (Abschnitt~\ref{sec:distinct-on}) ein
PostgreSQL-spezifisches Sprachmittel formal eingeführt, das das
Gruppenrepräsentantenproblem (Definition~\ref{def:gruppenrepraesentant})
in einem einzigen, deklarativen Schritt löst und gerade entlang der
Kanten eines Tiefenpfads -- wo die Gruppierungsspalte stets ein
lokaler Fremdschlüssel ist (Bemerkung~\ref{rem:distinct-on-tiefenpfad})
-- einen reinen Index-Scan ohne separate Sortierphase ermöglicht
(Proposition~\ref{prop:distinct-on-plan}).

Das zentrale Prinzip, das aus allen drei Ebenen gleichlautend folgt,
lautet:

\begin{quote}
\large\bfseries
Beim Entwurf relationaler Datenbankschemata ist stets die Tiefe der
Breite vorzuziehen. Relationen müssen in die Tiefe gehen, nicht in die
Breite.
\end{quote}

Darüber hinaus hat dieses Buch gezeigt, dass dieses Prinzip nicht isoliert
in der Entwurfsphase verbleibt, sondern sich konsistent durch
\emph{alle} Phasen des Lebenszyklus eines Datenbanksystems zieht: von
der konzeptuellen Modellierung (Kapitel~\ref{ch:modelle}) über die
Anfragesprache (Kapitel~\ref{ch:sql}) und den physischen Entwurf
(Kapitel~\ref{ch:entwurf}), den operativen Betrieb
(Kapitel~\ref{ch:betrieb}), die prozedurale Erweiterung
(Kapitel~\ref{ch:plsql}), die Anwendungsentwicklung
(Kapitel~\ref{ch:hostsprachen}) bis hin zur WWW-Integration
(Kapitel~\ref{ch:www}).

\section{Ausblick}
\label{sec:ausblick}

Die in diesem Buch entwickelte Theorie eröffnet eine Reihe
weiterführender Forschungs- und Anwendungsrichtungen, die im Folgenden
ausführlich diskutiert werden.

\subsection{Automatisierte Schemaanalyse und Refactoring-Werkzeuge}
\label{sec:ausblick-tooling}

Beispiel~\ref{ex:katalog-graph} hat gezeigt, dass der
Fremdschlüsselgraph $G_{\mathcal{S}}$ eines beliebigen relationalen
Schemas vollständig automatisiert aus dem Systemkatalog extrahiert
werden kann. Damit lassen sich die in
Definition~\ref{def:masse} eingeführten Maße $\depth(G_{\mathcal{S}})$,
$\breadth(G_{\mathcal{S}})$ und $\tau(G_{\mathcal{S}})$
(Tiefen-Breiten-Verhältnis) für jedes existierende Datenbankschema
berechnen. Ein naheliegendes und praktisch hochrelevantes
Forschungsvorhaben besteht in der Entwicklung eines
\emph{Schema-Linters}, der

\begin{itemize}
  \item für jede Relation $R$ eines Schemas $\mathcal{S}$ den
    \emph{lokalen Sterngrad} $\Delta(R)$ (Anzahl eingehender
    Fremdschlüsselkanten) berechnet und Relationen mit $\Delta(R)$
    oberhalb eines konfigurierbaren Schwellwerts als
    \emph{Hub-Kandidaten} markiert,
  \item für jeden Hub-Kandidaten automatisiert prüft, ob unter den auf
    ihn verweisenden Spoke-Relationen transitive funktionale
    Abhängigkeiten gemäß Theorem~\ref{thm:transitivitaet} bestehen, die
    eine weitere Zerlegung nahelegen,
  \item auf Basis dieser Analyse \emph{Refactoring-Vorschläge} in Form
    von konkreten \texttt{ALTER TABLE}- und \texttt{CREATE TABLE}-
    Anweisungen generiert, welche die betroffene Hub-Relation entlang
    der erkannten transitiven Abhängigkeiten in einen Tiefenpfad
    zerlegen, unter Wahrung der in Definition~\ref{def:lossless} und
    Definition~\ref{def:dep-erhaltend} geforderten Eigenschaften
    (Verlustfreiheit und Abhängigkeitserhaltung),
  \item begleitend dazu \texttt{CREATE VIEW}-Anweisungen erzeugt, welche
    die ursprüngliche, von Anwendungen erwartete flache Sicht auf die
    Hub-Relation als View über den neu erzeugten Tiefenpfad
    rekonstruieren -- im Sinne der in
    Abschnitt~\ref{sec:select} (\glqq Virtuelle Tabellen\grqq{})
    diskutierten Kompensation der scheinbaren Bequemlichkeit von
    Sternschemata durch materialisierte Pfade.
\end{itemize}

Ein solches Werkzeug würde die in dieser Arbeit rein formal geführte
Argumentation in eine praktisch anwendbare Methodik überführen und
ließe sich insbesondere in CI/CD-Pipelines für
Datenbankmigrationsskripte integrieren, sodass jede neue Migration
automatisch auf eine Verschlechterung von $\tau(G_{\mathcal{S}})$
geprüft wird, bevor sie in Produktion übernommen wird.

\subsection{Erweiterung auf NoSQL- und NewSQL-Systeme}
\label{sec:ausblick-nosql}

Die in Kapitel~\ref{ch:tiefenentwurf} entwickelten Maße $\depth$,
$\breadth$ und $\tau$ sind rein graphentheoretischer Natur und setzen
nicht zwingend ein relationales Datenmodell voraus. Eine
vielversprechende Erweiterungsrichtung besteht darin, das
tiefenorientierte Entwurfsparadigma auf dokumentenorientierte
(z.\,B.\ MongoDB), spaltenorientierte (z.\,B.\ Cassandra) und
graphbasierte (z.\,B.\ Neo4j) Datenbanksysteme zu übertragen:

\begin{itemize}
  \item In dokumentenorientierten Systemen entspricht die
    \emph{Einbettung} (Embedding) verwandter Dokumente einer
    Denormalisierung im Sinne von Proposition~\ref{prop:redundanz},
    während \emph{Referenzierung} zwischen Collections einer
    Fremdschlüsselkante im Sinne von Definition~\ref{def:er-graph}
    entspricht. Es stellt sich die Frage, ob sich
    Theorem~\ref{thm:anomaliefrei} (Anomaliefreiheit durch
    Tiefenpfade) sinngemäß auf Referenzierungsstrukturen in
    Dokumentenmodellen übertragen lässt, und unter welchen Bedingungen
    Embedding dennoch -- als bewusste, lokal begrenzte Ausnahme --
    gerechtfertigt ist.
  \item In graphbasierten Systemen ist der Fremdschlüsselgraph
    $G_{\mathcal{S}}$ nicht mehr nur ein analytisches Hilfsmittel,
    sondern die primäre Datenstruktur selbst. Hier wäre zu
    untersuchen, ob die in Abschnitt~\ref{sec:dfs_bfs} bewiesene
    Asymmetrie von DFS und BFS (Theorem~\ref{thm:asymmetrie}) auch für
    die Effizienz von Graphdatenbankabfragen (Cypher, Gremlin)
    quantitativ nachweisbar ist, insbesondere im Hinblick auf
    Pfadabfragen entlang tiefer Ketten im Vergleich zu
    sternförmigen Nachbarschaftsabfragen.
  \item In NewSQL-Systemen mit verteilter Speicherung (z.\,B.\
    CockroachDB, Spanner) hat die Topologie des Fremdschlüsselgraphen
    direkten Einfluss auf die \emph{Sharding-Strategie}: Tiefenpfade
    legen eine natürliche, hierarchische Partitionierung
    (Co-Location verwandter Tupel auf demselben Shard) nahe, während
    Sternschemata mit einem zentralen Hub zu einem
    \emph{Hot-Spot}-Shard führen können. Eine formale Analyse des
    Zusammenhangs zwischen $\tau(G_{\mathcal{S}})$ und der
    Lastverteilung über Shards ist eine vielversprechende
    Forschungsrichtung.
\end{itemize}

\subsection{Quantitative empirische Validierung}
\label{sec:ausblick-empirie}

Die in dieser Arbeit geführten Beweise sind rein formal-mathematischer
Natur. Eine konsequente Weiterführung besteht in einer breit angelegten
empirischen Studie, die

\begin{enumerate}
  \item eine repräsentative Stichprobe realer, produktiver
    Datenbankschemata aus offenen Datenquellen (z.\,B.\ Open-Source-
    Projekten auf GitHub/GitLab) hinsichtlich $\depth(G_{\mathcal{S}})$,
    $\breadth(G_{\mathcal{S}})$ und $\tau(G_{\mathcal{S}})$ vermisst,
  \item für jedes Schema die tatsächliche Verteilung von
    Anfragelaufzeiten (gemäß Theorem~\ref{thm:selektivitaet}) und
    Schreiblatenzen (gemäß der Schreibkostenproposition in
    Abschnitt~\ref{sec:abfragen}) unter realistischer Last misst,
  \item die Korrelation zwischen $\tau(G_{\mathcal{S}})$ und
    qualitativen Software-Engineering-Metriken untersucht, etwa der
    Häufigkeit von Schema-Migrationen (als Proxy für die in
    Theorem~\ref{thm:schaemaevolution} formalisierte
    Schemaevolutionskosten), der Anzahl von Bugtickets, die auf
    Dateninkonsistenzen zurückzuführen sind (als Proxy für
    Anomaliefreiheit gemäß Theorem~\ref{thm:anomaliefrei}), sowie der
    Onboarding-Zeit neuer Entwickler (als Proxy für
    Verständlichkeit und Testbarkeit, Abschnitt~\ref{sec:wartbarkeit}).
\end{enumerate}

Eine solche Studie würde es erlauben, die in dieser Arbeit rein
strukturell-formal hergeleiteten Vorteile tiefenorientierter Schemata
mit konkreten, in der Praxis beobachtbaren Effektgrößen zu
unterlegen und damit den Ãœbergang von einem theoretischen
Gütekriterium zu einer empirisch validierten Entwurfsrichtlinie zu
vollziehen.

\subsection{Integration in den Datenbankoptimierer}
\label{sec:ausblick-optimizer}

Abschnitt~\ref{sec:abfragen} hat gezeigt, dass moderne
Abfrageoptimierer durch Index-Only Scans, Hash-Joins, materialisierte
Views und Predicate Pushdown den theoretisch höheren Join-Aufwand
tiefer Pfade (Proposition~\ref{prop:jointiefe}) kompensieren können.
Eine konsequente Weiterentwicklung bestünde darin, dem Optimierer
\emph{explizites Wissen} über die Tiefenpfadstruktur des Schemas zur
Verfügung zu stellen, etwa in Form eines Schema-Annotationssystems, das
für jede Relation $R_i$ ihre Tiefe $i$ im Fremdschlüsselgraphen
$G_{\mathcal{S}}$ als Metadatum im Systemkatalog
(Abschnitt~\ref{sec:systemtabellen}) hinterlegt. Der Optimierer könnte
diese Information nutzen, um:

\begin{itemize}
  \item bei mehrstufigen Verbunden entlang eines Tiefenpfads die
    Verbundreihenfolge so zu wählen, dass zuerst die Relation mit der
    höchsten Tiefe (typischerweise der kleinsten Kardinalität) als
    Build-Seite eines Hash-Joins gewählt wird,
  \item für Abfragen, die ausschließlich Attribute einer Relation
    $R_d$ mit $d > 0$ filtern (vgl.\ Theorem~\ref{thm:selektivitaet}),
    automatisch zu erkennen, dass kein Zugriff auf $R_0$ (die Wurzel
    des Tiefenpfads) erforderlich ist, und den entsprechenden
    Verbund vollständig zu eliminieren (Join-Elimination),
  \item für \texttt{CASCADE}-Operationen (Abschnitt~\ref{ch:sql},
    Datenmodifikationen) den erwarteten Propagationsaufwand entlang
    des Tiefenpfads vorab zu schätzen und bei Überschreitung eines
    Schwellwerts eine Bestätigung des Datenbankadministrators
    einzufordern.
\end{itemize}

\subsection{Lehre und curriculare Integration}
\label{sec:ausblick-lehre}

Schließlich hat dieses Buch -- in Anlehnung an die didaktische
Konzeption von Kleinschmidt und Rank~\cite{Kleinschmidt2002} -- gezeigt,
dass sich ein anspruchsvolles formales Entwurfsprinzip vollständig in
den klassischen Kanon einer einführenden Datenbankvorlesung integrieren
lässt, ohne dass dafür zusätzliche Vorlesungszeit in Form eines
separaten Theoriekapitels erforderlich wäre. Stattdessen wurde das
Prinzip \emph{horizontal} durch alle Kapitel hindurch verwoben: bei der
Diskussion von Views (Kapitel~\ref{ch:sql}), bei der
Schlüsseldefinition (Kapitel~\ref{ch:entwurf}), bei
Sperrgranularität und Zugriffsrechten (Kapitel~\ref{ch:betrieb}), bei
Triggern (Kapitel~\ref{ch:plsql}), bei der Modularisierung von
Anwendungscode (Kapitel~\ref{ch:hostsprachen}) und bei der
Architektur von WWW-Anwendungen (Kapitel~\ref{ch:www}).

Ein naheliegender nächster Schritt für die Lehre besteht darin, diese
horizontale Integration durch ein \emph{durchgängiges Projekt} zu
ergänzen: Studierende entwerfen zu Beginn des Semesters ein eigenes
Sternschema für eine selbstgewählte Domäne, berechnen
$\depth$, $\breadth$ und $\tau$ gemäß Definition~\ref{def:masse},
identifizieren transitive Abhängigkeiten gemäß
Theorem~\ref{thm:transitivitaet} und führen im Verlauf des Semesters --
parallel zur Behandlung der entsprechenden Kapitel -- ein vollständiges
Refactoring zu einem Tiefenpfadschema durch, einschließlich
SQL-DDL (Kapitel~\ref{ch:entwurf}), PL/SQL-Triggern
(Kapitel~\ref{ch:plsql}) und einer einfachen WWW-Oberfläche
(Kapitel~\ref{ch:www}). Die am Ende des Semesters erreichten Werte von
$\tau(G_{\mathcal{S}})$ vor und nach dem Refactoring liefern dabei ein
unmittelbar greifbares, quantitatives Lernziel, das die in diesem Buch
entwickelte Theorie mit der praktischen Entwurfskompetenz der
Studierenden verbindet.

\subsection{Herstellerspezifische Erweiterungen am Beispiel \texttt{DISTINCT ON}: Spannungsfeld zu Portabilität und Tiefenpfadarchitektur}
\label{sec:ausblick-distinct-on}

Abschnitt~\ref{sec:distinct-on} hat mit \texttt{DISTINCT ON} ein
Sprachmittel formal untersucht, das -- anders als die in
Kapitel~\ref{ch:hostsprachen} betonte Portabilität
(Abschnitt~\ref{sec:portabilitaet}) -- \emph{nicht} Teil des
SQL-Standards ist, sondern eine Eigenheit von PostgreSQL darstellt.
Dies wirft eine Reihe von Fragen auf, die über den Rahmen dieses
Buches hinausweisen und die im Folgenden als Ausblick skizziert
werden.

\textbf{Portabilitätsanalyse für tiefenpfadtypische Anfragemuster.}
Theorem~\ref{thm:distinct-on-aequivalenz} hat gezeigt, dass
\texttt{DISTINCT ON} semantisch äquivalent zu einer
\texttt{ROW\_NUMBER()}-basierten Filterung ist, die wiederum
Standard-SQL ist und auf praktisch allen relationalen DBMS
(einschließlich der in Abschnitt~\ref{sec:codd-rules} und
Kapitel~\ref{ch:hostsprachen} diskutierten Systeme) verfügbar ist.
Eine naheliegende Erweiterung dieser Arbeit bestünde darin, das in
Abschnitt~\ref{sec:repository-pattern} eingeführte
Repository-Pattern um eine \emph{Übersetzungsschicht} zu ergänzen,
die für jede Kante $(R_i,R_{i+1})$ des Tiefenpfads, an der ein
Gruppenrepräsentantenproblem auftritt
(Definition~\ref{def:gruppenrepraesentant}), automatisch zwischen der
PostgreSQL-Formulierung mittels \texttt{DISTINCT ON} und der
portablen \texttt{ROW\_NUMBER()}-Formulierung umschalten kann -- je
nachdem, welches DBMS die Anwendung gemäß
Abschnitt~\ref{sec:portabilitaet} ansteuert. Eine solche
Übersetzungsschicht ließe sich systematisch aus der in
Bemerkung~\ref{rem:distinct-on-regel} formalisierten Syntaxregel
(Präfixbedingung zwischen \texttt{DISTINCT ON}-Liste und
\texttt{ORDER BY}-Liste) ableiten, da diese Regel bereits die exakte
Korrespondenz zwischen $G$ (der \texttt{PARTITION BY}-Liste von
\texttt{ROW\_NUMBER()}) und $(B_1,\ldots,B_q)$ (der
\texttt{ORDER BY}-Liste innerhalb der Partition) festlegt.

\textbf{Erweiterung des Schema-Linters um \texttt{DISTINCT
ON}-Indexempfehlungen.}
Der in Abschnitt~\ref{sec:ausblick-tooling} skizzierte
Schema-Linter ließe sich um eine weitere Analyseklasse erweitern: Für
jede Kante $(R_i,R_{i+1})$ des Fremdschlüsselgraphen
$G_{\mathcal{S}}$ könnte das Werkzeug den Anwendungs-Quellcode (bzw.\
ein Anfrageprotokoll, vgl.\ \texttt{pg\_stat\_statements}) nach
Anfragen der Form
\texttt{DISTINCT ON (}$F$\texttt{) \ldots ORDER BY }$F,B_1,\ldots,B_q$
durchsuchen, wobei $F\subseteq\att(R_i)$ der Fremdschlüssel auf
$R_{i+1}$ ist. Wird ein solches Muster erkannt, aber kein
zusammengesetzter Index $(F,B_1,\ldots,B_q)$ auf $R_i$ gefunden, so
kann das Werkzeug -- analog zu
Beispiel~\ref{ex:distinct-on-tiefenpfad} und gestützt auf
Proposition~\ref{prop:distinct-on-plan} -- automatisiert eine
\texttt{CREATE INDEX}-Anweisung vorschlagen, die den
\texttt{Sort}-Schritt des Ausführungsplans eliminiert und die
asymptotischen Kosten von $O(|R_i|\log|R_i|)$ auf
$O(|\pi_{F}(R_i)|+\log|R_i|)$ reduziert.

\textbf{Allgemeineres Prinzip: Herstellererweiterungen als
„Kompressionen“ von Standardmustern.}
Die in diesem Buch an \texttt{DISTINCT ON} demonstrierte Methodik --
formale Definition (Definition~\ref{def:distinct-on}), Nachweis der
Äquivalenz zu einem Standard-SQL-Konstrukt
(Theorem~\ref{thm:distinct-on-aequivalenz}) und anschließende
Effizienzanalyse im Kontext des Tiefenpfadparadigmas
(Proposition~\ref{prop:distinct-on-plan},
Bemerkung~\ref{rem:distinct-on-tiefenpfad}) -- ist nicht auf
\texttt{DISTINCT ON} beschränkt. Zahlreiche weitere
herstellerspezifische Erweiterungen -- etwa \texttt{LATERAL}-Joins,
\texttt{FILTER}-Klauseln für Aggregatfunktionen, oder
\texttt{INSERT \ldots ON CONFLICT} (\glqq Upsert\grqq{}) -- lassen
sich nach demselben Schema untersuchen: Zunächst die Formalisierung
des gelösten Problems als algebraische bzw.\
mengentheoretische Definition (analog zu
Definition~\ref{def:gruppenrepraesentant}), dann der Nachweis der
Äquivalenz zu einer (meist umständlicheren) Standard-SQL-Formulierung,
und schließlich die Untersuchung, an welchen Stellen des
Fremdschlüsselgraphen $G_{\mathcal{S}}$ eines Tiefenpfadschemas diese
Erweiterung den größten Effizienz- oder Lesbarkeitsgewinn erbringt.
Eine systematische, nach diesem Schema aufgebaute
\glqq Erweiterungs-Bibliothek\grqq{} für tiefenpfadorientierte
Schemata wäre eine wertvolle Ergänzung sowohl für die Lehre
(Abschnitt~\ref{sec:ausblick-lehre}) als auch für die in
Abschnitt~\ref{sec:ausblick-tooling} und hier skizzierten
Tooling-Vorhaben.

\subsection{Temporale Datenbanken und die zeitliche Dimension des Tiefenpfads}
\label{sec:ausblick-temporal}

Sämtliche bisherigen Betrachtungen dieses Buches gingen implizit von
einem \emph{Schnappschuss}-Modell aus: Eine Relation $\R$ repräsentiert
den aktuellen Zustand der Welt, und frühere Zustände werden -- wenn
überhaupt -- über separate Audit- oder Historientabellen
(Abschnitt~\ref{ch:plsql}, Trigger-Beispiel) nachgehalten. Die
Datenbanktheorie kennt jedoch mit den \emph{temporalen Datenbanken}
einen eigenständigen, gut entwickelten Zweig, der die Zeit als
Erstklassenattribut in das Relationenmodell integriert. Eine
ausführliche Weiterführung dieses Buches in Richtung temporaler
Datenbanken ist aus mehreren Gründen naheliegend.

\textbf{Bitemporale Erweiterung des Tiefenpfadschemas.} In der
temporalen Datenbanktheorie unterscheidet man zwischen
\emph{Gültigkeitszeit} (valid time, der Zeitraum, in dem ein Faktum in
der modellierten Welt zutrifft) und \emph{Transaktionszeit}
(transaction time, der Zeitraum, in dem ein Faktum in der Datenbank
als bekannt gespeichert war). Eine \emph{bitemporale} Relation $R$
führt zu jedem Tupel $t$ zwei Intervalle
$[t_{vs},t_{ve})$ (valid time) und $[t_{ts},t_{te})$ (transaction
time) mit. Für ein Tiefenpfadschema $R_1\to R_2\to\cdots\to R_l$ stellt
sich die Frage, wie sich diese zeitliche Dimension mit der
Tiefenstruktur verträgt:

\begin{itemize}
  \item \textbf{Definition eines zeitparametrisierten
    Fremdschlüsselgraphen.} Man definiere
    $G_{\mathcal{S}}(\tau)=(V,E_\tau)$ als denjenigen
    Fremdschlüsselgraphen, der entsteht, wenn man jede Relation $R_i$
    auf ihren zum Zeitpunkt $\tau$ gültigen Schnappschuss
    $R_i(\tau)=\{t\in R_i \mid t_{vs}\leq\tau<t_{ve}\}$ einschränkt.
    Eine zentrale, in dieser Arbeit unbeantwortete Frage lautet: Bleibt
    die Tiefe $\depth(G_{\mathcal{S}}(\tau))$ für alle $\tau$
    konstant, oder kann sich die \emph{Topologie} des
    Fremdschlüsselgraphen über die Zeit ändern (etwa durch
    \emph{Reparenting}, also die nachträgliche Änderung eines
    Fremdschlüsselwerts, z.\,B.\ wenn ein Kunde einer anderen Region
    zugeordnet wird)? Falls ja, müsste Definition~\ref{def:masse}
    ($\depth$, $\breadth$, $\tau(G)$) zu zeitabhängigen Funktionen
    $\depth(\tau)$, $\breadth(\tau)$, $\tau_{\text{Verh}}(\tau)$
    erweitert werden, und Theorem~\ref{thm:transitivitaet}
    (Transitivitätsverteilung) müsste für \emph{zeitlich variable}
    funktionale Abhängigkeiten -- sogenannte \emph{temporale
    funktionale Abhängigkeiten} (TFDs) -- neu bewiesen werden.

  \item \textbf{Slowly Changing Dimensions als getarnte
    Tiefenpfaderweiterung.} Im Data-Warehouse-Bereich
    (Kimball-Methodik, vgl.\ Definition~\ref{def:stern} und
    \cite{Kimball1996}) ist das Muster der \emph{Slowly Changing
    Dimension} (SCD) vom Typ~2 weit verbreitet: Anstatt ein
    Dimensionsattribut bei einer Änderung zu überschreiben, wird ein
    neues Tupel mit einem neuen Gültigkeitsintervall angelegt, und das
    alte Tupel wird mit einem Enddatum versehen. Eine bemerkenswerte
    Beobachtung -- die in einer Erweiterung dieser Arbeit formal
    auszuarbeiten wäre -- ist, dass SCD2 \emph{strukturell} eine
    \emph{Tiefenerweiterung} im Sinne von
    Korollar~\ref{cor:tiefenpfad3NF} darstellt: Die ursprüngliche
    funktionale Abhängigkeit $K\to A$ (z.\,B.\
    $\texttt{kunden\_id}\to\texttt{region\_id}$) wird durch die
    Einführung der Zeitdimension zu einer
    \emph{zeitparametrisierten} Abhängigkeit
    $(K,\tau)\to A$, was bedeutet, dass $(K,\tau)$ -- nicht mehr $K$
    allein -- der (zusammengesetzte) Schlüssel der Relation ist. Damit
    wird aus einer scheinbar „flachen“ Dimensionstabelle implizit ein
    Tiefenpfad der Form $\textsc{Kunde\_Version} \to \textsc{Kunde}$,
    bei dem $\textsc{Kunde}$ die zeitinvarianten Attribute (z.\,B.\
    die natürliche Schlüssel-ID) und $\textsc{Kunde\_Version}$ die
    zeitveränderlichen Attribute samt Gültigkeitsintervall enthält.
    Eine vollständige formale Ausarbeitung dieser Beobachtung würde
    Theorem~\ref{thm:transitivitaet} um einen Fall „transitiver
    Abhängigkeit über die Zeitachse“ erweitern und damit zeigen, dass
    \emph{auch} die zeitliche Dimension dem Prinzip „Tiefe vor Breite“
    unterliegt.

  \item \textbf{\texttt{DISTINCT ON} als natürliches Werkzeug für
    temporale Schnappschüsse.} Abschnitt~\ref{sec:distinct-on} hat
    \texttt{DISTINCT ON} als Lösung des
    Gruppenrepräsentantenproblems (Definition~\ref{def:gruppenrepraesentant})
    eingeführt. Für eine bitemporale Relation $R_i$ mit
    Transaktionszeit-Intervallen liefert
    \begin{align*}
      &\texttt{SELECT DISTINCT ON (}K\texttt{) * FROM } R_i\\
      &\texttt{WHERE } t_{vs}\leq\tau \texttt{ AND } \tau < t_{ve}\\
      &\texttt{ORDER BY } K, t_{ts} \texttt{ DESC}
    \end{align*}
    für jeden Schlüsselwert $K$ genau die zum
    Transaktionszeitpunkt $\tau$ \emph{aktuellste} Version des zum
    Gültigkeitszeitpunkt $\tau$ gültigen Tupels -- eine direkte
    Anwendung von Theorem~\ref{thm:distinct-on-aequivalenz} auf das
    bitemporale Modell. Eine ausführliche Weiterentwicklung dieses
    Buches könnte einen eigenen Abschnitt „Bitemporale Anfragen mit
    \texttt{DISTINCT ON} und Bereichstypen (\texttt{tstzrange})“
    enthalten, der PostgreSQLs native Unterstützung für
    Intervalltypen und Exclusion-Constraints
    (\texttt{EXCLUDE USING gist}) mit der in
    Bemerkung~\ref{rem:distinct-on-tiefenpfad} entwickelten
    Indexstrategie kombiniert.

  \item \textbf{SQL:2011-Standard für temporale Tabellen.} Der
    SQL:2011-Standard (vgl.\ Abschnitt~\ref{sec:sql-standards})
    definiert \texttt{PERIOD FOR}-Klauseln sowie
    \texttt{SYSTEM\_VERSIONING} für transaktionszeit-versionierte
    Tabellen. Eine ausführliche Behandlung dieses Standards --
    einschließlich der Frage, wie sich
    Theorem~\ref{thm:repository-zerlegung} (Eindeutigkeit der
    Repository-Zerlegung) auf system-versionierte Tabellen überträgt,
    bei denen \emph{jede} Lesezugriffsmethode implizit um einen
    \texttt{AS OF}-Zeitpunkt parametrisiert werden muss -- wäre ein
    eigenständiges, umfangreiches Kapitel wert.
\end{itemize}

\subsection{Verteilte Transaktionen, Konsistenzmodelle und der Tiefenpfad in Microservice-Architekturen}
\label{sec:ausblick-verteilt}

Kapitel~\ref{ch:betrieb} hat das Transaktionskonzept innerhalb
\emph{eines} DBMS behandelt. In modernen
Microservice-Architekturen ist es jedoch üblich, dass
unterschiedliche Teile eines ehemals monolithischen Schemas
$\mathcal{S}$ auf unterschiedliche Datenbankinstanzen -- mit
unterschiedlichen DBMS-Produkten -- aufgeteilt werden
(\emph{Database-per-Service}-Muster). Dies wirft fundamentale Fragen
auf, die eine erhebliche Erweiterung dieser Arbeit rechtfertigen.

\textbf{Schnitte durch den Fremdschlüsselgraphen.} Sei
$G_{\mathcal{S}}=(V,E)$ der Fremdschlüsselgraph des
Gesamtschemas. Eine Aufteilung auf $p$ Microservices entspricht einer
Partition $V=V_1\,\dot\cup\,\cdots\,\dot\cup\,V_p$. Jede Kante
$(R_i,R_j)\in E$ mit $R_i\in V_a$, $R_j\in V_b$, $a\neq b$, wird zu
einer \emph{verteilten Fremdschlüsselbeziehung}, deren referentielle
Integrität (Definition~\ref{def:schluessel}) nicht mehr durch das DBMS
selbst, sondern nur noch durch Anwendungslogik (z.\,B.\
Saga-Pattern, Event-Sourcing mit Kompensationstransaktionen)
sichergestellt werden kann.

\begin{itemize}
  \item \textbf{Tiefenpfad-erhaltende Schnitte.} Eine naheliegende
    Designhypothese -- die in einer Erweiterung dieser Arbeit formal
    zu untersuchen wäre -- lautet: Schnitte durch
    $G_{\mathcal{S}}$, die \emph{entlang} eines Tiefenpfads verlaufen
    (d.\,h.\ jede Schnittkante $(R_i,R_{i+1})$ trennt zwei
    aufeinanderfolgende Knoten eines Pfades), erzeugen
    \emph{weniger} verteilte Fremdschlüsselbeziehungen pro
    Microservice als Schnitte, die einen Hub-Knoten eines
    Sternschemas (Definition~\ref{def:stern}) durchtrennen -- da ein
    Hub-Knoten $H$ mit $\deg(H)=k$ bei einer Aufteilung, die $H$ von
    $m<k$ seiner Spokes trennt, sofort $m$ verteilte
    Fremdschlüsselbeziehungen erzeugt, während ein Schnitt entlang
    eines Tiefenpfads stets nur \emph{eine} Kante durchtrennt,
    unabhängig von der Pfadlänge $l$. Ein formaler Beweis dieser
    Hypothese, etwa in Form eines Theorems
    „Tiefenpfad-Schnitte minimieren die Anzahl verteilter
    Fremdschlüsselbeziehungen pro Partitionsgrenze“, wäre ein
    bedeutender Brückenschlag zwischen der in
    Kapitel~\ref{ch:tiefenentwurf} entwickelten Theorie und der
    Praxis verteilter Systeme.

  \item \textbf{Sagas als verteilte Tiefenpfadtraversierung.} Das
    \emph{Saga-Pattern} ersetzt eine klassische ACID-Transaktion
    (Kapitel~\ref{ch:einleitung}) durch eine Folge lokaler
    Transaktionen $T_1,\ldots,T_n$ mit zugehörigen
    Kompensationstransaktionen $C_1,\ldots,C_n$, sodass im Fehlerfall
    $C_{n-1},\ldots,C_1$ in umgekehrter Reihenfolge ausgeführt
    werden. Für eine Operation, die einen Tiefenpfad
    $R_1\to R_2\to\cdots\to R_l$ traversiert (z.\,B.\ „lege einen
    neuen Auftrag mit Positionen an, reserviere Lagerbestand,
    initiiere Zahlung“), entspricht jeder Schritt $T_i$ einer
    lokalen Transaktion auf demjenigen Microservice, der $R_i$
    verwaltet. Eine ausführliche Untersuchung könnte zeigen, dass die
    \emph{Reihenfolge} der Saga-Schritte $T_1,\ldots,T_l$ und die
    Reihenfolge der Knoten im Tiefenpfad $R_1,\ldots,R_l$
    übereinstimmen sollten, damit die Kompensationsreihenfolge
    $C_l,\ldots,C_1$ exakt der \emph{umgekehrten}
    Fremdschlüsselrichtung entspricht -- d.\,h.\ Kompensationen
    propagieren von „Kindknoten“ zu „Elternknoten“ des Tiefenpfads,
    analog zu \texttt{ON DELETE CASCADE} (Abschnitt~\ref{sec:select}),
    jedoch auf Anwendungsebene und über Servicegrenzen hinweg.

  \item \textbf{CAP-Theorem und lokale vs.\ globale Konsistenz
    entlang des Pfads.} Das CAP-Theorem (Konsistenz,
    Verfügbarkeit, Partitionstoleranz -- es können höchstens zwei der
    drei Eigenschaften gleichzeitig garantiert werden) erzwingt in
    verteilten Systemen häufig den Übergang zu
    \emph{eventueller Konsistenz}. Eine umfangreiche Erweiterung
    dieses Buches könnte untersuchen, ob sich entlang eines
    Tiefenpfads $R_1\to\cdots\to R_l$ eine \emph{abgestufte}
    Konsistenzgarantie formulieren lässt: Knoten $R_1$ (z.\,B.\
    Stammdaten wie $\textsc{Land}$, $\textsc{Region}$) ändern sich
    selten und können stark konsistent (mit Lese-Repliken)
    repliziert werden, während Knoten $R_l$ (z.\,B.\
    $\textsc{Auftragsposition}$) hochfrequent geschrieben werden und
    eventuelle Konsistenz tolerieren. Die Frage, ob ein
    \emph{Konsistenzgradient} entlang der Tiefe $\depth(G_{\mathcal{S}})$
    -- von „stark konsistent“ an der Wurzel zu „eventuell konsistent“
    an den Blättern -- eine allgemeingültige Entwurfsregel darstellt,
    wäre ein eigenständiges Forschungsprogramm.
\end{itemize}

\subsection{Sicherheits- und Datenschutzaspekte: Verschlüsselung, Mandantenfähigkeit und Anonymisierung entlang des Tiefenpfads}
\label{sec:ausblick-sicherheit}

Abschnitt~\ref{sec:zugriffsrechte} hat die klassische
SQL-Rechteverwaltung (\texttt{GRANT}/\texttt{REVOKE}) im Kontext des
Tiefenpfadparadigmas behandelt. Eine sehr viel umfangreichere
Behandlung von Sicherheits- und Datenschutzaspekten wäre für eine
erweiterte Auflage dieses Buches angezeigt, insbesondere vor dem
Hintergrund der Datenschutz-Grundverordnung (DSGVO) und vergleichbarer
Regelwerke.

\textbf{Spaltenverschlüsselung und Tiefenlokalität.} Werden
einzelne Attribute (z.\,B.\ $\textsc{Kunde}.\texttt{name}$,
$\textsc{Kunde}.\texttt{adresse}$) anwendungsseitig oder
transparent (Transparent Data Encryption, TDE) verschlüsselt, so
stellt sich die Frage, wie sich dies mit den in
Abschnitt~\ref{sec:abfragen} bewiesenen Selektivitätsvorteilen
(Theorem~\ref{thm:selektivitaet}) verträgt: Ein Index auf einer
verschlüsselten Spalte ist im Allgemeinen nicht mehr ordnungserhaltend
nutzbar, sodass \texttt{DISTINCT ON}-Indizes
(Abschnitt~\ref{sec:distinct-on-effizienz}) auf verschlüsselten
Spalten ihre Effizienzvorteile verlieren können. Eine ausführliche
Untersuchung könnte systematisch klassifizieren, welche Attribute
eines Tiefenpfadschemas (Schlüssel, Fremdschlüssel,
Sortier-/Filterattribute, rein darstellende Attribute) für welche
Verschlüsselungsklasse (deterministisch, randomisiert,
ordnungserhaltend, homomorph) geeignet sind, und dabei explizit auf
die Rolle jedes Attributs im Fremdschlüsselgraphen $G_{\mathcal{S}}$
Bezug nehmen: Fremdschlüsselattribute (Definition~\ref{def:schluessel})
müssen z.\,B.\ \emph{deterministisch} verschlüsselt werden, damit
der kanonische Verbund (Definition~\ref{def:fk-join}) weiterhin auf
Gleichheit prüfen kann, während rein darstellende Attribute (z.\,B.\
\texttt{name}) randomisiert verschlüsselt werden können.

\textbf{Mandantenfähigkeit (Multi-Tenancy) als zusätzliche
Tiefenstufe.} In Software-as-a-Service-Anwendungen teilen sich
mehrere Mandanten (Tenants) dieselbe Datenbankinstanz. Die gängigen
Architekturmuster -- \emph{Shared Database, Shared Schema}
(\texttt{tenant\_id} als zusätzliche Spalte in jeder Tabelle),
\emph{Shared Database, Separate Schemas} und \emph{Separate
Databases} -- lassen sich im Lichte des Tiefenpfadparadigmas wie
folgt einordnen: Im Muster „Shared Schema“ wird \texttt{tenant\_id}
zu einem zusätzlichen, \emph{universellen} Fremdschlüsselattribut, das
in \emph{jeder} Relation $R_i\in\mathcal{S}$ vorkommt und auf eine
zentrale $\textsc{Mandant}$-Relation verweist. Formal entsteht damit
ein neuer Knoten $\textsc{Mandant}$ in $G_{\mathcal{S}}$ mit einer
Kante zu \emph{jedem} anderen Knoten -- also $\deg(\textsc{Mandant})=
|V|-1$, was $\textsc{Mandant}$ nach Definition~\ref{def:stern} zu
einer extremen Hub-Relation macht. Eine ausführliche Behandlung
müsste klären, ob dieser scheinbare Widerspruch zur in
Korollar~\ref{cor:tiefenpfad3NF} geforderten Tiefenorientierung
tatsächlich einen Widerspruch darstellt, oder ob $\textsc{Mandant}$
als \emph{orthogonale} Dimension (eine Art „globale Wurzel“, die
\emph{vor} dem eigentlichen Tiefenpfad jedes Mandanten steht)
aufzufassen ist -- mit der Konsequenz, dass für \emph{jeden}
einzelnen Mandanten der durch \texttt{tenant\_id}$=$const gefilterte
Teilgraph von $G_{\mathcal{S}}$ wieder ein reiner Tiefenpfad ist und
$\textsc{Mandant}$ selbst außerhalb der in
Theorem~\ref{thm:repository-zerlegung} betrachteten
Repository-Zerlegung steht (jedes Repository-Modul $M_i$ filtert
implizit nach \texttt{tenant\_id}, ohne dass dies die Modulgrenzen
verändert).

\textbf{Anonymisierung, Pseudonymisierung und das Recht auf
Löschung.} Art.~17 DSGVO („Recht auf Löschung“) verlangt, dass
personenbezogene Daten auf Anfrage vollständig entfernt werden. Im
Tiefenpfadschema (z.\,B.\ $\textsc{Land}\leftarrow\textsc{Region}
\leftarrow\textsc{Kunde}\leftarrow\textsc{Auftrag}\leftarrow
\textsc{Auftragsposition}$) stellt sich die Frage, wie eine
\texttt{DELETE}-Operation auf $\textsc{Kunde}$
(Abschnitt~\ref{sec:select}, Datenmodifikationen) mit den
betriebswirtschaftlichen und gegebenenfalls handels- bzw.\
steuerrechtlichen Aufbewahrungspflichten für
$\textsc{Auftrag}$-Daten in Einklang zu bringen ist. Eine umfangreiche
Erweiterung dieses Buches könnte ein formales Modell der
\emph{„Löschbarkeitstiefe“} entwickeln: Für jeden Knoten $R_i$ eines
Tiefenpfads wird eine Funktion $\ell(R_i)\in\{\text{sofort},
\text{nach Frist }d_i,\text{nie}\}$ definiert, und es wird untersucht,
unter welchen Bedingungen an $\ell$ die referentielle Integrität
(Definition~\ref{def:schluessel}) entlang des Pfads aufrechterhalten
werden kann, wenn ein „Löschen“ tatsächlich als
\emph{Pseudonymisierung} (Ersetzen personenbezogener Attribute durch
Platzhalterwerte, unter Beibehaltung des Primärschlüssels und damit
aller Fremdschlüsselbeziehungen) realisiert wird. Dies stünde in
direktem Bezug zu Theorem~\ref{thm:anomaliefrei}
(Anomaliefreiheit durch Tiefenpfade): Da im Tiefenpfadschema jedes
Attribut genau einmal gespeichert ist, betrifft eine
Pseudonymisierungsoperation auf $\textsc{Kunde}$ \emph{ausschließlich}
die Relation $\textsc{Kunde}$ -- eine direkte Anwendung der in
Theorem~\ref{thm:schaemaevolution} bewiesenen Lokalität, übertragen
von Schemaänderungen auf Datenänderungsoperationen mit
Compliance-Charakter.

\subsection{Formale Verifikation von Schemaeigenschaften mit Beweisassistenten}
\label{sec:ausblick-verifikation}

Sämtliche Theoreme dieses Buches (z.\,B.\
Theorem~\ref{thm:transitivitaet},
Theorem~\ref{thm:lossless},
Theorem~\ref{thm:anomaliefrei},
Theorem~\ref{thm:synthese-korrekt},
Theorem~\ref{thm:distinct-on-aequivalenz}) wurden in der für die
mathematische Fachliteratur üblichen Form „informeller, aber
rigoroser“ Beweise (vgl.\ \cite{Cormen2009}) geführt. Eine
weitreichende und methodisch anspruchsvolle Erweiterung dieser Arbeit
bestünde darin, diese Beweise in einem interaktiven
Beweisassistenten -- etwa Coq, Isabelle/HOL oder Lean -- zu
formalisieren und maschinell zu verifizieren.

\begin{itemize}
  \item \textbf{Formalisierung des Relationenmodells.} Als
    Ausgangspunkt müsste das Relationenmodell
    (Definition~\ref{def:relation}) als abhängiger Typ formalisiert
    werden: Eine Relation über dem Schema $(A_1{:}D_1,\ldots,A_n{:}D_n)$
    entspricht einer endlichen Menge (bzw.\ einem endlichen
    Multimenge-Typ, falls Duplikate vor Anwendung von
    \texttt{DISTINCT}, vgl.\ Definition~\ref{def:distinct-standard},
    zugelassen werden) von $n$-Tupeln über den entsprechenden
    Trägertypen. Funktionale Abhängigkeiten
    (Abschnitt~\ref{sec:normalformen}) und die Hüllenoperation
    $X^+_\Sigma$ (Definition~\ref{def:attributhuelle}) ließen sich
    als induktiv definierte Prädikate bzw.\ als Fixpunkt einer
    monotonen Funktion auf der Potenzmenge der Attribute formalisieren
    -- Letzteres ist besonders attraktiv, da Beweisassistenten über
    ausgereifte Bibliotheken für Fixpunkttheorie auf vollständigen
    Verbänden verfügen, was einen maschinenüberprüften Beweis der
    Terminierung der \texttt{HUELLE}-Funktion aus
    Algorithmus~\ref{alg:minimalcover} ermöglichen würde.

  \item \textbf{Maschinenüberprüfter Beweis von
    Theorem~\ref{thm:lossless}.} Das
    Lossless-Join-Kriterium ist ein zentrales technisches Lemma, auf
    dem die gesamte Dekompositionstheorie
    (Abschnitt~\ref{sec:dekomposition}) aufbaut. Eine Formalisierung
    müsste den natürlichen Verbund $\Join$
    (Abschnitt~\ref{sec:relalg}) als Operation auf den oben
    eingeführten Relationstypen definieren und die im
    Beweis von Theorem~\ref{thm:lossless} verwendete
    Tupelkonstruktion (insbesondere die Existenz „verirrter
    Tupel“ -- spurious tuples -- im Falle der Notwendigkeitsrichtung)
    explizit als Gegenbeispiel-Konstruktion formalisieren. Eine
    erfolgreiche Formalisierung würde die Korrektheit von
    Algorithmus~\ref{alg:synthese} (Synthesealgorithmus) auf eine
    \emph{maschinenprüfbare} Grundlage stellen.

  \item \textbf{Extraktion eines zertifizierten
    Schema-Linters.} Beweisassistenten wie Coq erlauben die
    \emph{Extraktion} von verifizierten Algorithmen in
    ausführbaren Code (z.\,B.\ OCaml oder Haskell). Eine
    weitreichende praktische Konsequenz einer formalen Verifikation
    von Algorithmus~\ref{alg:minimalcover} und
    Algorithmus~\ref{alg:synthese} wäre die Möglichkeit, den in
    Abschnitt~\ref{sec:ausblick-tooling} skizzierten Schema-Linter
    als \emph{zertifizierten} Algorithmus zu extrahieren, dessen
    Korrektheit (im Sinne von
    Theorem~\ref{thm:synthese-korrekt} und
    Theorem~\ref{thm:synthese-minimal}) nicht nur durch
    Testfälle, sondern durch einen maschinenüberprüften Beweis
    garantiert ist -- ein Qualitätsniveau, das in der
    Datenbank-Tooling-Landschaft bislang nicht erreicht wird.

  \item \textbf{Modellprüfung (Model Checking) von
    Transaktionsschedules.} Die in Kapitel~\ref{ch:betrieb}
    behandelten Konzepte -- Serialisierbarkeit, Isolationsstufen,
    die in Bemerkung~\ref{rem:sperren} diskutierte
    Sperrgranularität entlang des Tiefenpfads -- sind klassische
    Anwendungsgebiete der \emph{Modellprüfung} (model checking) mit
    Werkzeugen wie TLA+ oder SPIN. Eine umfangreiche Erweiterung
    könnte für ein konkretes Tiefenpfadschema (z.\,B.\ das in
    Abschnitt~\ref{sec:beispiel} verwendete
    $\textsc{Land}\leftarrow\textsc{Region}\leftarrow\textsc{Kunde}
    \leftarrow\textsc{Auftrag}\leftarrow\textsc{Auftragsposition}$)
    ein TLA+-Modell der nebenläufigen Transaktionen entwickeln und
    formal verifizieren, dass unter \texttt{REPEATABLE READ}
    (Kapitel~\ref{ch:betrieb}) keine der drei in
    Abschnitt~\ref{sec:anomalien-kurz} klassifizierten Anomalien
    auftreten kann -- und dabei explizit zeigen, dass die Anzahl der
    zu prüfenden Zustände (state space) für ein Tiefenpfadschema
    aufgrund der in Theorem~\ref{thm:repository-zerlegung} bewiesenen
    Lokalität signifikant kleiner ist als für ein äquivalentes
    Sternschema -- ein weiteres, diesmal aus der Modellprüfung
    stammendes quantitatives Argument für „Tiefe vor Breite“.
\end{itemize}

\subsection{Maschinelles Lernen auf Tiefenpfadschemata}
\label{sec:ausblick-ml}

Ein wachsendes Anwendungsfeld relationaler Datenbanken ist die
Bereitstellung von Trainingsdaten für maschinelles Lernen,
insbesondere für \emph{Graph Neural Networks} (GNNs) und
\emph{relationales Deep Learning}, die direkt auf der Struktur
$G_{\mathcal{S}}$ operieren, anstatt zunächst eine flache
„Feature-Tabelle“ zu erzeugen.

\textbf{Relationales Deep Learning und der Fremdschlüsselgraph als
Berechnungsgraph.} Ein GNN definiert für jeden Knoten $v$ eines
Graphen eine Repräsentation $h_v$, die iterativ durch Aggregation der
Repräsentationen seiner Nachbarn aktualisiert wird:
$h_v^{(k+1)} = f(h_v^{(k)}, \{h_u^{(k)} \mid u \in
N(v)\})$. Wendet man dies direkt auf $G_{\mathcal{S}}$ an, so
entspricht ein \emph{Knoten} des GNN nicht einer \emph{Relation}
$R_i$, sondern einem \emph{Tupel} $t\in R_i$, und eine
\emph{Kante} des GNN entspricht einem konkreten Fremdschlüsselverweis
zwischen zwei Tupeln (also einer Instanz des in
Definition~\ref{def:fk-join} definierten kanonischen Verbunds, auf
Tupelebene statt auf Relationenebene). Eine ausführliche Erweiterung
dieser Arbeit könnte folgende Fragen behandeln:

\begin{itemize}
  \item \textbf{Empfangsfeld (Receptive Field) und Pfadtiefe.}
    Nach $k$ GNN-Schichten hängt $h_v^{(k)}$ von allen Tupeln ab, die
    über einen Pfad der Länge $\leq k$ in $G_{\mathcal{S}}$ von $v$
    aus erreichbar sind. Für ein Tiefenpfadschema mit
    $\depth(G_{\mathcal{S}})=l$ bedeutet dies, dass $k=l$ Schichten
    nötig sind, um Information von der Wurzel bis zu den Blättern
    (oder umgekehrt) zu propagieren -- eine direkte
    Ãœbertragung der in Proposition~\ref{prop:jointiefe}
    definierten Join-Tiefe $\delta$ in die Sprache der GNN-Architektur:
    $\delta(\mathcal{Q})$ und die für ein GNN benötigte Schichtzahl
    $k$ sind für dieselbe semantische „Reichweite“ identisch. Für ein
    $k$-Sternschema (Definition~\ref{def:stern}) hingegen genügt
    bereits \emph{eine} Schicht, um vom Hub $H$ aus alle Spokes zu
    erreichen -- jedoch um den Preis, dass $H$ in dieser einen
    Schicht Information von $k$ strukturell verschiedenen
    Entitätstypen aggregieren muss, was in der GNN-Literatur als
    \emph{Over-Squashing} bekanntes Problem auftritt. Eine
    formale Untersuchung des Zusammenhangs zwischen
    $\tau(G_{\mathcal{S}})$ (Definition~\ref{def:masse}) und dem
    Over-Squashing-Phänomen wäre ein eigenständiger,
    publikationswürdiger Forschungsbeitrag.

  \item \textbf{Feature-Engineering entlang von Repository-Grenzen.}
    Theorem~\ref{thm:repository-zerlegung} hat gezeigt, dass ein
    Tiefenpfadschema eine eindeutige Zerlegung des Anwendungscodes in
    Module $M_1,\ldots,M_l$ induziert. Eine naheliegende Erweiterung
    wäre die Übertragung dieser Zerlegung auf die
    \emph{Feature-Pipeline} eines ML-Systems: Für jeden Knoten $R_i$
    existiert ein \emph{Feature-Extraktor} $\phi_i: R_i \to
    \mathbb{R}^{d_i}$, und das Gesamtfeature eines
    $\textsc{Auftrag}$-Tupels für ein Vorhersagemodell (z.\,B.\
    „Wahrscheinlichkeit einer Rückgabe“) ergibt sich durch
    Konkatenation $\phi_1(t_1)\oplus\cdots\oplus\phi_l(t_l)$ entlang
    des Pfads von der Wurzel zum betrachteten Tupel. Eine
    ausführliche Behandlung könnte zeigen, dass diese
    Pfad-Konkatenation -- ähnlich wie die in
    Abschnitt~\ref{sec:lazy-loading} diskutierte Pfadtraversierung
    via Lazy Loading -- entweder \emph{eager} (als ein einziger
    materialisierter Feature-View, vgl.\
    Abschnitt~\ref{sec:select}, „Virtuelle Tabellen“) oder
    \emph{lazy} (zur Trainingszeit aus den Rohdaten berechnet)
    realisiert werden kann, mit denselben Kompromissen wie beim
    $N{+}1$-Problem (Bemerkung~\ref{rem:n-plus-1}).

  \item \textbf{Vektorindizes und approximative Nächste-Nachbarn-Suche
    als neue Kantenart.} Moderne PostgreSQL-Erweiterungen
    (\texttt{pgvector}) erlauben die Speicherung von
    Embedding-Vektoren als Spaltentyp und die effiziente
    Nächste-Nachbarn-Suche mittels Indizes (HNSW, IVF-Flat). Dies
    eröffnet die Möglichkeit, $G_{\mathcal{S}}$ um eine neue Art von
    Kante zu erweitern: \emph{semantische Ähnlichkeitskanten}
    zwischen Tupeln derselben Relation $R_i$, die nicht durch
    Fremdschlüssel, sondern durch Ähnlichkeit im
    Embedding-Raum definiert sind. Eine Erweiterung dieser Arbeit
    könnte formal untersuchen, ob und wie sich die in
    Kapitel~\ref{ch:tiefenentwurf} entwickelten Maße $\depth$,
    $\breadth$, $\tau$ auf einen \emph{hybriden} Graphen
    übertragen lassen, der sowohl klassische
    Fremdschlüsselkanten (gemäß Definition~\ref{def:er-graph}) als
    auch $k$-Nächste-Nachbarn-Kanten im Embedding-Raum enthält, und
    welche Konsequenzen dies für die Indexstrategien aus
    Abschnitt~\ref{sec:distinct-on-effizienz} hat -- etwa die Frage,
    ob \texttt{DISTINCT ON} in Kombination mit einer
    \texttt{ORDER BY embedding <-> :query\_vector}-Klausel
    (Nächste-Nachbarn-Operator von \texttt{pgvector}) weiterhin der in
    Theorem~\ref{thm:distinct-on-aequivalenz} bewiesenen Äquivalenz
    zu \texttt{ROW\_NUMBER()} genügt.
\end{itemize}

\subsection{Hybride Graph-relationale Systeme und alternative
Anfragesprachen}
\label{sec:ausblick-hybrid}

Abschnitt~\ref{sec:ausblick-nosql} hat bereits angedeutet, dass
graphbasierte Systeme (Neo4j u.\,a.) den Fremdschlüsselgraphen
$G_{\mathcal{S}}$ als primäre Datenstruktur behandeln. Eine
ausführliche Erweiterung dieser Arbeit könnte einen vollständigen
Vergleich der in diesem Buch entwickelten Theorie mit den
Anfragesprachen \emph{Cypher}, \emph{Gremlin} und dem neuen
ISO-Standard \emph{GQL} (Graph Query Language, ISO/IEC 39075:2024)
durchführen. Von besonderem Interesse wäre dabei:

\begin{itemize}
  \item \textbf{Pfadausdrücke als natives Sprachmittel.} Während in
    SQL ein Tiefenpfad $R_1\to R_2\to\cdots\to R_l$ durch $l-1$
    explizite \texttt{JOIN}-Klauseln (Beispiel~\ref{ex:join-tiefe})
    traversiert werden muss, erlaubt Cypher Pfadausdrücke der Form
    \texttt{(r1)-[:GEHOERT\_ZU]->(r2)-[:GEHOERT\_ZU]->\ldots->(rl)}
    direkt als Teil des Anfragemusters. Eine formale Untersuchung
    könnte zeigen, dass für \emph{reine} Tiefenpfadschemata
    (Definition~\ref{def:tiefenpfad}) eine Cypher-Anfrage mit einem
    Pfadausdruck der Länge $l-1$ \emph{syntaktisch} prägnanter, aber
    \emph{semantisch äquivalent} zu der in
    Abschnitt~\ref{sec:select} eingeführten SQL-Verbundkette ist --
    ähnlich der in Theorem~\ref{thm:distinct-on-aequivalenz}
    geführten Äquivalenzbeweise --, während für Sternschemata
    (Definition~\ref{def:stern}) Cypher-Pfadausdrücke der Form
    \texttt{(h)-[]->(s1)}, \texttt{(h)-[]->(s2)}, \ldots keinen
    Prägnanzgewinn gegenüber SQL bieten, da kein gemeinsamer Pfad
    existiert.

  \item \textbf{Property Graphs als Generalisierung des
    Tiefenpfadschemas.} Im Property-Graph-Modell trägt nicht nur
    jeder Knoten, sondern auch jede \emph{Kante} Attribute. Dies
    erlaubt es, $n{:}m$-Beziehungen (vgl.\
    Abschnitt~\ref{sec:er-zu-rel}, Transformationsregel~3) ohne ein
    explizites Verbindungs-Relationenschema zu modellieren: Die
    Attribute einer $n{:}m$-Beziehung (z.\,B.\
    $\textsc{Auftragsposition}.\texttt{menge}$) werden direkt zu
    Kantenattributen. Eine ausführliche Untersuchung könnte
    formalisieren, unter welchen Bedingungen diese Verschiebung von
    „Verbindungsrelation als Knoten“ zu „Verbindungsrelation als
    Kante“ die in Proposition~\ref{prop:er-rel-graph} bewiesene
    Struktur des Fremdschlüsselgraphen $G_{\mathcal{S}}$ erhält bzw.\
    vereinfacht, und ob sich dadurch die Tiefe
    $\depth(G_{\mathcal{S}})$ um genau die Anzahl der auf diese Weise
    „eliminierten“ Verbindungsrelationen reduziert.

  \item \textbf{SQL/PGQ (SQL:2023).} Der jüngste SQL-Standard
    SQL:2023 führt mit \texttt{SQL/PGQ} (Property Graph Queries) die
    Möglichkeit ein, \emph{innerhalb} einer relationalen Datenbank
    eine Graph-Sicht (\texttt{CREATE PROPERTY GRAPH}) auf
    existierende Tabellen zu definieren und mit einer an Cypher
    angelehnten Syntax (\texttt{GRAPH\_TABLE}) abzufragen. Dies ist
    möglicherweise der bedeutendste Standardisierungsschritt seit
    SQL:1999 (Abschnitt~\ref{sec:sql-standards}) für die in diesem
    Buch entwickelte Theorie, da er erlaubt, denselben
    Fremdschlüsselgraphen $G_{\mathcal{S}}$ \emph{wahlweise} als
    relationale Tabellen (klassisches SQL,
    Kapitel~\ref{ch:sql}) oder als Property Graph
    (\texttt{GRAPH\_TABLE}) anzufragen -- ohne Datenduplikation. Eine
    sehr ausführliche Erweiterung dieses Buches sollte SQL/PGQ zum
    Anlass nehmen, sämtliche in Kapitel~\ref{ch:sql} eingeführten
    Beispielanfragen (insbesondere Beispiel~\ref{ex:join-tiefe} und
    Beispiel~\ref{ex:distinct-on-tiefenpfad}) in \emph{beiden}
    Syntaxformen zu präsentieren und empirisch (vgl.\
    Abschnitt~\ref{sec:ausblick-empirie}) zu vergleichen, welche
    Formulierung für Tiefenpfad- bzw.\ Sternschemata vom
    PostgreSQL-Optimierer (sobald SQL/PGQ dort verfügbar ist) in
    effizientere Ausführungspläne übersetzt wird.
\end{itemize}

\subsection{Ein Forschungsprogramm: Vom Gütekriterium zur Theorie der Schema-Topologie}
\label{sec:ausblick-forschungsprogramm}

Die in den vorangegangenen Abschnitten skizzierten Erweiterungen --
automatisiertes Tooling (Abschnitt~\ref{sec:ausblick-tooling}),
NoSQL/NewSQL (Abschnitt~\ref{sec:ausblick-nosql}), empirische
Validierung (Abschnitt~\ref{sec:ausblick-empirie}),
Optimierer-Integration (Abschnitt~\ref{sec:ausblick-optimizer}),
Lehre (Abschnitt~\ref{sec:ausblick-lehre}), herstellerspezifische
Erweiterungen (Abschnitt~\ref{sec:ausblick-distinct-on}), temporale
Datenbanken (Abschnitt~\ref{sec:ausblick-temporal}), verteilte
Systeme (Abschnitt~\ref{sec:ausblick-verteilt}), Sicherheit
(Abschnitt~\ref{sec:ausblick-sicherheit}), formale Verifikation
(Abschnitt~\ref{sec:ausblick-verifikation}), maschinelles Lernen
(Abschnitt~\ref{sec:ausblick-ml}) und hybride Graph-relationale
Systeme (Abschnitt~\ref{sec:ausblick-hybrid}) -- sind keine
unabhängigen Einzelthemen, sondern lassen sich als Facetten
\emph{eines einzigen}, übergeordneten Forschungsprogramms auffassen,
das hier abschließend skizziert werden soll: einer
\emph{allgemeinen Theorie der Schema-Topologie}.

\textbf{Leitfrage.} Die Leitfrage eines solchen Forschungsprogramms
lautet: \emph{Existiert für jede relevante Eigenschaft $P$ eines
Datenbanksystems (Normalisierungsgrad, Anomaliefreiheit,
Anfrageeffizienz, Wartbarkeit, Sperrverhalten, Verteilungseignung,
Eignung für maschinelles Lernen, Sicherheitsklassifizierbarkeit,
Verifizierbarkeit, \ldots) eine Funktion $f_P: G_{\mathcal{S}} \to
\mathbb{R}$ (oder $\to$ einen geordneten Wertebereich), sodass $P$
\emph{monoton} in $f_P(G_{\mathcal{S}})$ ist, und ist $f_P$ für alle
relevanten $P$ \emph{simultan} durch dasselbe Strukturmerkmal von
$G_{\mathcal{S}}$ -- nämlich seine Tiefen-Breiten-Charakteristik
$\tau(G_{\mathcal{S}})$ aus Definition~\ref{def:masse} -- bestimmt?}

Dieses Buch hat für eine Reihe konkreter Eigenschaften $P$ -- 3NF
(Korollar~\ref{cor:tiefenpfad3NF}), Anomaliefreiheit
(Theorem~\ref{thm:anomaliefrei}), Selektivität
(Theorem~\ref{thm:selektivitaet}), Schreibkomplexität, lokale
Änderbarkeit (Theorem~\ref{thm:schaemaevolution}), Modulzerlegbarkeit
(Theorem~\ref{thm:repository-zerlegung}) -- jeweils \emph{einzeln}
gezeigt, dass $P$ in $\tau(G_{\mathcal{S}})$ monoton ist (höheres
$\tau$, also „mehr Tiefe relativ zu Breite“, impliziert „besseres
$P$“). Die in diesem Ausblick skizzierten Erweiterungen --
insbesondere die temporale Erweiterung von $\tau$ zu
$\tau(\tau_{\text{Zeit}})$ (Abschnitt~\ref{sec:ausblick-temporal}),
die verteilungsbezogene Interpretation von $\tau$
(Abschnitt~\ref{sec:ausblick-verteilt}) und die
GNN-Receptive-Field-Interpretation von $\tau$
(Abschnitt~\ref{sec:ausblick-ml}) -- legen nahe, dass $\tau$ tatsächlich
ein \emph{universelles} Strukturmerkmal sein könnte, das weit über die
in diesem Buch behandelten Einzelresultate hinaus als
\glqq Ein-Zahl-Zusammenfassung\grqq{} der Qualität eines
Datenbankschemas dient -- vergleichbar der Rolle, die die
\emph{Konditionszahl} einer Matrix in der numerischen linearen
Algebra spielt: Eine einzelne Kennzahl, die unmittelbar Auskunft über
das Verhalten eines Systems unter einer Vielzahl unterschiedlicher
Operationen gibt.

Ein vollständiger Beweis dieser \glqq universellen Monotonie-These\grqq{}
-- für \emph{alle} relevanten Eigenschaften $P$ und nicht nur für die
in diesem Buch einzeln behandelten -- wäre ein Forschungsprogramm, das
sich über mehrere Dissertationen erstrecken könnte und das die
Datenbanktheorie mit der Graphentheorie (insbesondere der Theorie der
Baumweite und verwandter struktureller Graphparameter, die in der
algorithmischen Graphentheorie eine ganz ähnliche Rolle spielen wie
$\tau$ in dieser Arbeit) auf eine neue, einheitliche Grundlage stellen
würde. Die vorliegende Arbeit versteht sich als erster, in sich
geschlossener Baustein dieses größeren Programms.

\subsection{Schlussbemerkung}

Die Leitidee dieses Buches -- \emph{Tiefe vor Breite} -- ist sowohl
formal beweisbar als auch praktisch handlungsleitend. Die in den
vorangegangenen Kapiteln entwickelte Theorie liefert nicht nur eine
Erklärung dafür, \emph{warum} normalisierte, tiefenorientierte Schemata
in der Praxis robuster, wartbarer und konsistenter sind, sondern auch
ein \emph{Werkzeug} -- die Maße $\depth$, $\breadth$ und $\tau$ --, mit
dem sich diese Eigenschaft für jedes beliebige Schema quantitativ
bewerten lässt. Die in Abschnitt~\ref{sec:ausblick-tooling} bis
\ref{sec:ausblick-lehre} skizzierten Weiterführungen zeigen, dass dieses
Werkzeug weit über den Rahmen einer einführenden Darstellung hinaus
Anwendung finden kann -- von automatisierten Refactoring-Werkzeugen über
verteilte NewSQL-Architekturen bis hin zur Gestaltung
datenbankzentrierter Curricula.


\appendix
\chapter{Syntaxübersicht und ergänzende Materialien}

\section{Ãœbersicht zentraler SQL-Sprachelemente}

Tabelle~\ref{tab:sql-uebersicht} fasst die in diesem Buch verwendeten
zentralen SQL-Sprachelemente und ihren Bezug zu den entsprechenden
Kapiteln zusammen.

\begin{table}[H]
\centering
\small
\begin{tabular}{@{}llp{6cm}@{}}
\toprule
\textbf{Sprachelement} & \textbf{Kapitel} & \textbf{Bezug zum Tiefenpfadparadigma} \\
\midrule
\texttt{CREATE TABLE \ldots REFERENCES} & \ref{ch:entwurf} &
  Definiert eine Kante im Fremdschlüsselgraphen $G_{\mathcal{S}}$ \\
\texttt{SELECT \ldots JOIN} & \ref{ch:sql} &
  Traversiert eine oder mehrere Kanten von $G_{\mathcal{S}}$ \\
\texttt{CREATE VIEW} & \ref{ch:sql} &
  Materialisiert einen Pfad in $G_{\mathcal{S}}$ als flache Sicht \\
\texttt{ON DELETE CASCADE} & \ref{ch:sql}, \ref{ch:betrieb} &
  Propagation entlang einer Kante von $G_{\mathcal{S}}$ \\
\texttt{GRANT}/\texttt{REVOKE} & \ref{ch:betrieb} &
  Entitätspräzise Rechtevergabe pro Knoten von $G_{\mathcal{S}}$ \\
\texttt{CREATE TRIGGER} & \ref{ch:plsql} &
  An genau einen Knoten von $G_{\mathcal{S}}$ gebundene Logik \\
\texttt{information\_schema} & \ref{ch:betrieb} &
  Quelle zur automatisierten Rekonstruktion von $G_{\mathcal{S}}$ \\
\bottomrule
\end{tabular}
\caption{Zentrale SQL-Sprachelemente und ihr Bezug zum Fremdschlüsselgraphen}
\label{tab:sql-uebersicht}
\end{table}

\section{Glossar der formalen Symbole}

\begin{table}[H]
\centering
\small
\begin{tabular}{@{}lp{10cm}@{}}
\toprule
\textbf{Symbol} & \textbf{Bedeutung} \\
\midrule
$\R$, $R$ & Relation bzw.\ Relationenschema \\
$\att(R)$ & Attributmenge von $R$ \\
$\Sigma_R$ & Menge der Integritätsbedingungen von $R$ \\
$K \to A$ & Funktionale Abhängigkeit: $K$ bestimmt $A$ \\
$X \twoheadrightarrow Y$ & Mehrwertige Abhängigkeit \\
$G_{\mathcal{S}} = (V,E)$ & Fremdschlüsselgraph eines Schemas $\mathcal{S}$ \\
$\depth(G)$ & Tiefe von $G$ (längster Pfad) \\
$\breadth(G)$ & Breite von $G$ (maximaler Verzweigungsgrad) \\
$\tau(G)$ & Tiefen-Breiten-Verhältnis $\depth(G)/\breadth(G)$ \\
$\delta(\mathcal{Q})$ & Join-Tiefe einer Anfrage $\mathcal{Q}$ \\
$\rho$ & Redundanzfaktor \\
$\sigma_P$, $\pi$, $\Join$, $\rho_{A\to B}$ & Selektion, Projektion, Verbund, Umbenennung (relationale Algebra) \\
\bottomrule
\end{tabular}
\caption{Glossar der in diesem Buch verwendeten formalen Symbole}
\end{table}

\section{Hinweise zu Software-Bezugsquellen}

In Anlehnung an die im Vorbild von Kleinschmidt und
Rank~\cite{Kleinschmidt2002} enthaltenen Hinweise auf
Software-Bezugsquellen sei angemerkt, dass die in diesem Buch
verwendeten SQL- und PL/SQL-Konstrukte sich an verbreiteten
Open-Source- und kommerziellen relationalen Datenbanksystemen
orientieren und mit geringfügigen syntaktischen Anpassungen auf diesen
Systemen lauffähig sind. Für die in Kapitel~\ref{ch:hostsprachen}
behandelte DBI-Schnittstelle sind entsprechende Treibermodule für die
gängigen Open-Source-Datenbanksysteme verfügbar.


\bibliographystyle{plainnat}
\begin{thebibliography}{99}

% ============================================================
% Grundlagen: Relationenmodell, Normalformen, Algebra
% ============================================================

\bibitem{Armstrong1974}
Armstrong, W.\,W. (1974).
\newblock Dependency structures of data base relationships.
\newblock \emph{Proceedings of IFIP Congress}, 74, 580--583.
\newblock Grundlegende Arbeit zu den nach Armstrong benannten
Inferenzregeln für funktionale Abhängigkeiten, die der
Hüllenoperation $X^+_\Sigma$ (Definition~\ref{def:attributhuelle})
zugrunde liegen.

\bibitem{Ambler2003}
Ambler, S.\,W. (2003).
\newblock \emph{Agile Database Techniques: Effective Strategies for the Agile
Software Developer}.
\newblock Wiley.
\newblock Praxisorientierte Darstellung von Schema-Refactoring-Techniken,
relevant für den in Abschnitt~\ref{sec:ausblick-tooling} skizzierten
Schema-Linter.

\bibitem{Boyce1974}
Boyce, R.\,F.; Codd, E.\,F. (1974).
\newblock Relational completeness of data base sublanguages.
\newblock In: Rustin, R. (Hrsg.): \emph{Data Base Systems}, Courant Computer
Science Symposia, Vol.\,6, Prentice-Hall, S.\,65--98.
\newblock Ursprungsarbeit zur Boyce-Codd-Normalform (Definition~3.4).

\bibitem{Codd1970}
Codd, E.\,F. (1970).
\newblock A relational model of data for large shared data banks.
\newblock \emph{Communications of the ACM}, 13(6), 377--387.
\newblock Die Gründungsarbeit des relationalen Modells
(Definition~\ref{def:relation}); Ausgangspunkt sämtlicher in diesem
Buch behandelten Theorie.

\bibitem{Codd1974}
Codd, E.\,F. (1974).
\newblock Recent investigations in relational data base systems.
\newblock \emph{Proceedings of IFIP Congress}, 1021--1036.
\newblock Formuliert 1NF, 2NF und 3NF (Abschnitt~\ref{sec:normalformen})
sowie frühe Fassungen der in Abschnitt~\ref{sec:codd-rules}
diskutierten Zusatzforderungen.

\bibitem{Fagin1977}
Fagin, R. (1977).
\newblock Multivalued dependencies and a new normal form for relational
databases.
\newblock \emph{ACM Transactions on Database Systems}, 2(3), 262--278.
\newblock Einführung der vierten Normalform (Definition~3.5) und der
mehrwertigen Abhängigkeiten, zentral für
Proposition~\ref{prop:mvd_stern}.

\bibitem{Fagin1981}
Fagin, R. (1981).
\newblock A normal form for relational databases that is based on domains
and keys.
\newblock \emph{ACM Transactions on Database Systems}, 6(3), 387--415.
\newblock Einführung der Domain-Key-Normalform (DKNF), der
„stärksten“ klassischen Normalform, als Bezugspunkt für
Theorem~\ref{thm:synthese-korrekt}.

\bibitem{Ullman1988}
Ullman, J.\,D. (1988).
\newblock \emph{Principles of Database and Knowledge-Base Systems}, Vol.\,1.
\newblock Computer Science Press.
\newblock Standardquelle für das Lossless-Join-Lemma
(Theorem~\ref{thm:lossless}) und den Synthesealgorithmus
(Algorithmus~\ref{alg:synthese}).

\bibitem{Date2003}
Date, C.\,J. (2003).
\newblock \emph{An Introduction to Database Systems} (8th ed.).
\newblock Addison-Wesley.
\newblock Umfassendes einführendes Lehrbuch; Quelle der
Anomalieklassifikation in Abschnitt~\ref{sec:anomalien}.

\bibitem{Elmasri2015}
Elmasri, R.; Navathe, S.\,B. (2015).
\newblock \emph{Fundamentals of Database Systems} (7th ed.).
\newblock Pearson.
\newblock Breit angelegtes Lehrbuch mit Kapiteln zu ER-Modellierung
(Abschnitt~\ref{sec:ermodell}), Normalisierung und verteilten
Datenbanken (Abschnitt~\ref{sec:ausblick-verteilt}).

\bibitem{Kleinschmidt2002}
Kleinschmidt, P.; Rank, C. (2002).
\newblock \emph{Relationale Datenbanksysteme: Eine praktische Einführung}
(2., überarbeitete und erweiterte Auflage).
\newblock Springer-Verlag, Berlin Heidelberg.
\newblock Strukturelles Vorbild für die Gliederung dieses Buches
(vgl.\ Abschnitt~1.4).

\bibitem{Ramakrishnan2003}
Ramakrishnan, R.; Gehrke, J. (2003).
\newblock \emph{Database Management Systems} (3rd ed.).
\newblock McGraw-Hill.
\newblock Quelle für die in Proposition~\ref{prop:jointiefe} verwendete
Diskussion der Join-Tiefe und Anfrageoptimierung.

\bibitem{Graefe1993}
Graefe, G. (1993).
\newblock Query evaluation techniques for large databases.
\newblock \emph{ACM Computing Surveys}, 25(2), 73--169.
\newblock Ãœbersichtsarbeit zu Hash-Joins und weiteren
Verbundalgorithmen, zitiert in der Bemerkung zu
Proposition~\ref{prop:jointiefe}.

\bibitem{Brewer2000}
Brewer, E.\,A. (2000).
\newblock Towards robust distributed systems (Keynote).
\newblock \emph{Proceedings of ACM PODC}.
\newblock Erste Formulierung des CAP-Theorems, zentral für
Abschnitt~\ref{sec:ausblick-verteilt}.

\bibitem{GilbertLynch2002}
Gilbert, S.; Lynch, N. (2002).
\newblock Brewer's conjecture and the feasibility of consistent, available,
partition-tolerant web services.
\newblock \emph{ACM SIGACT News}, 33(2), 51--59.
\newblock Formaler Beweis des CAP-Theorems, Grundlage der Diskussion
des Konsistenzgradienten in Abschnitt~\ref{sec:ausblick-verteilt}.

\bibitem{Garcia-Molina1987}
Garcia-Molina, H.; Salem, K. (1987).
\newblock Sagas.
\newblock \emph{Proceedings of ACM SIGMOD}, 249--259.
\newblock Ursprüngliche Definition des Saga-Patterns, aufgegriffen in
Abschnitt~\ref{sec:ausblick-verteilt} im Kontext der
Tiefenpfadtraversierung.

\bibitem{Vogels2009}
Vogels, W. (2009).
\newblock Eventually consistent.
\newblock \emph{Communications of the ACM}, 52(1), 40--44.
\newblock Praxisnahe Darstellung eventueller Konsistenz, Grundlage
für den in Abschnitt~\ref{sec:ausblick-verteilt} postulierten
Konsistenzgradienten entlang des Tiefenpfads.

\bibitem{Corbett2013}
Corbett, J.\,C.; Dean, J.; Epstein, M.; Fikes, A.; Frost, C.;
Furman, J.\,J.; Ghemawat, S.; Gubarev, A.; Heiser, C.; Hochschild, P.;
et al. (2013).
\newblock Spanner: Google's globally distributed database.
\newblock \emph{ACM Transactions on Computer Systems}, 31(3), 1--22.
\newblock Beispiel eines NewSQL-Systems mit globaler Verteilung,
relevant für Abschnitt~\ref{sec:ausblick-nosql} und
Abschnitt~\ref{sec:ausblick-verteilt}.

\bibitem{Newman2021}
Newman, S. (2021).
\newblock \emph{Building Microservices} (2nd ed.).
\newblock O'Reilly.
\newblock Standardwerk zu Microservice-Architekturen einschließlich
Database-per-Service-Mustern (Abschnitt~\ref{sec:ausblick-verteilt}).

% ============================================================
% Temporale Datenbanken
% ============================================================

\bibitem{Snodgrass1999}
Snodgrass, R.\,T. (1999).
\newblock \emph{Developing Time-Oriented Database Applications in SQL}.
\newblock Morgan Kaufmann.
\newblock Standardreferenz für bitemporale Datenmodellierung,
Grundlage von Abschnitt~\ref{sec:ausblick-temporal}.

\bibitem{Jensen1994}
Jensen, C.\,S.; Clifford, J.; Elmasri, R.; Gadia, S.\,K.; Hayes, P.;
Jajodia, S.; et al. (1994).
\newblock A consensus glossary of temporal database concepts.
\newblock \emph{ACM SIGMOD Record}, 23(1), 52--64.
\newblock Begriffliche Grundlage für Gültigkeitszeit und
Transaktionszeit, verwendet in Abschnitt~\ref{sec:ausblick-temporal}.

\bibitem{Kimball1996}
Kimball, R. (1996).
\newblock \emph{The Data Warehouse Toolkit}.
\newblock Wiley.
\newblock Quelle für Sternschemata (Definition~\ref{def:stern}) und
Slowly Changing Dimensions, diskutiert in
Abschnitt~\ref{sec:ausblick-temporal}.

\bibitem{KimballRoss2013}
Kimball, R.; Ross, M. (2013).
\newblock \emph{The Data Warehouse Toolkit: The Definitive Guide to
Dimensional Modeling} (3rd ed.).
\newblock Wiley.
\newblock Aktualisierte Darstellung der Slowly-Changing-Dimension-Muster
(Typ~1 bis Typ~7), aufgegriffen in
Abschnitt~\ref{sec:ausblick-temporal}.

\bibitem{Kulkarni2012}
Kulkarni, K.; Michels, J.-E. (2012).
\newblock Temporal features in SQL:2011.
\newblock \emph{ACM SIGMOD Record}, 41(3), 34--43.
\newblock Beschreibung der \texttt{PERIOD FOR}- und
\texttt{SYSTEM\_VERSIONING}-Klauseln aus
Abschnitt~\ref{sec:ausblick-temporal}.

% ============================================================
% Datenschutz, Sicherheit, Mandantenfähigkeit
% ============================================================

\bibitem{Saltzer1975}
Saltzer, J.\,H.; Schroeder, M.\,D. (1975).
\newblock The protection of information in computer systems.
\newblock \emph{Proceedings of the IEEE}, 63(9), 1278--1308.
\newblock Formuliert das Prinzip der minimalen Rechtevergabe,
zentral für Abschnitt~\ref{sec:zugriffsrechte}.

\bibitem{Popa2011}
Popa, R.\,A.; Redfield, C.\,M.\,S.; Zeldovich, N.; Balakrishnan, H. (2011).
\newblock CryptDB: Protecting confidentiality with encrypted query processing.
\newblock \emph{Proceedings of ACM SOSP}, 85--100.
\newblock Grundlagenarbeit zu ordnungserhaltender und deterministischer
Verschlüsselung von Datenbankspalten, zentral für
Abschnitt~\ref{sec:ausblick-sicherheit}.

\bibitem{Curino2011}
Curino, C.; Jones, E.; Popa, R.\,A.; Malviya, N.; Wu, E.; Madden, S.;
Balakrishnan, H.; Zeldovich, N. (2011).
\newblock Relational Cloud: A database-as-a-service for the cloud.
\newblock \emph{Proceedings of CIDR}, 235--240.
\newblock Architekturmuster für Mandantenfähigkeit (Multi-Tenancy),
diskutiert in Abschnitt~\ref{sec:ausblick-sicherheit}.

\bibitem{Voigt2017}
Voigt, P.; von dem Bussche, A. (2017).
\newblock \emph{The EU General Data Protection Regulation (GDPR): A
Practical Guide}.
\newblock Springer.
\newblock Juristische Grundlage des in
Abschnitt~\ref{sec:ausblick-sicherheit} diskutierten
„Rechts auf Löschung“ und der Pseudonymisierungsanforderungen.

% ============================================================
% Formale Verifikation und Beweisassistenten
% ============================================================

\bibitem{Cormen2009}
Cormen, T.\,H.; Leiserson, C.\,E.; Rivest, R.\,L.; Stein, C. (2009).
\newblock \emph{Introduction to Algorithms} (3rd ed.).
\newblock MIT Press.
\newblock Referenz für Beweisstil und Laufzeitanalyse, vgl.\
Theorem~\ref{thm:laufzeit-minimalcover}.

\bibitem{Tarjan1972}
Tarjan, R.\,E. (1972).
\newblock Depth-first search and linear graph algorithms.
\newblock \emph{SIAM Journal on Computing}, 1(2), 146--160.
\newblock Grundlagenarbeit zu DFS und dem DFS-Wald
(Definition~\ref{def:dfswald}).

\bibitem{Bertot2004}
Bertot, Y.; Castéran, P. (2004).
\newblock \emph{Interactive Theorem Proving and Program Development:
Coq'Art: The Calculus of Inductive Constructions}.
\newblock Springer.
\newblock Standardeinführung in den Beweisassistenten Coq, Grundlage
der in Abschnitt~\ref{sec:ausblick-verifikation} skizzierten
Formalisierung.

\bibitem{Nipkow2002}
Nipkow, T.; Paulson, L.\,C.; Wenzel, M. (2002).
\newblock \emph{Isabelle/HOL: A Proof Assistant for Higher-Order Logic}.
\newblock Springer (LNCS 2283).
\newblock Alternative zu Coq für die in
Abschnitt~\ref{sec:ausblick-verifikation} diskutierte
maschinengeprüfte Formalisierung relationaler Algebra.

\bibitem{Lamport2002}
Lamport, L. (2002).
\newblock \emph{Specifying Systems: The TLA+ Language and Tools for
Hardware and Software Engineers}.
\newblock Addison-Wesley.
\newblock Grundlage der in Abschnitt~\ref{sec:ausblick-verifikation}
vorgeschlagenen Modellprüfung von Transaktionsschedules.

\bibitem{Holzmann2003}
Holzmann, G.\,J. (2003).
\newblock \emph{The SPIN Model Checker: Primer and Reference Manual}.
\newblock Addison-Wesley.
\newblock Alternatives Modellprüfungswerkzeug, ergänzend zu TLA+ in
Abschnitt~\ref{sec:ausblick-verifikation}.

% ============================================================
% Maschinelles Lernen, Graphdatenbanken, hybride Systeme
% ============================================================

\bibitem{Hamilton2020}
Hamilton, W.\,L. (2020).
\newblock \emph{Graph Representation Learning}.
\newblock Morgan \& Claypool.
\newblock Einführendes Lehrbuch zu Graph Neural Networks, Grundlage
der Receptive-Field-Diskussion in
Abschnitt~\ref{sec:ausblick-ml}.

\bibitem{Fey2023}
Fey, M.; Hu, W.; Huang, K.; Lenssen, J.\,E.; Ranjan, R.; Robinson, J.;
Leskovec, J.; et al. (2023).
\newblock Relational deep learning: Graph representation learning on
relational databases.
\newblock \emph{arXiv preprint arXiv:2312.04615}.
\newblock Aktuelle Arbeit zu relationalem Deep Learning direkt auf
Fremdschlüsselgraphen $G_{\mathcal{S}}$, zentrale Referenz für
Abschnitt~\ref{sec:ausblick-ml}.

\bibitem{Topping2021}
Topping, J.; Di Giovanni, F.; Chamberlain, B.\,P.; Dong, X.;
Bronstein, M.\,M. (2021).
\newblock Understanding over-squashing and bottlenecks on graphs via
curvature.
\newblock \emph{Proceedings of ICLR}.
\newblock Formale Untersuchung des Over-Squashing-Phänomens,
aufgegriffen in Abschnitt~\ref{sec:ausblick-ml} im Vergleich von
Tiefenpfad- und Sternschemata.

\bibitem{Angles2018}
Angles, R.; Arenas, M.; Barceló, P.; Hogan, A.; Reutter, J.;
Vrgoč, D. (2018).
\newblock Foundations of modern query languages for graph databases.
\newblock \emph{ACM Computing Surveys}, 50(5), 1--40.
\newblock Ãœbersichtsarbeit zu Cypher, Gremlin und SPARQL, Grundlage
für den Vergleich in Abschnitt~\ref{sec:ausblick-hybrid}.

\bibitem{Robinson2015}
Robinson, I.; Webber, J.; Eifrem, E. (2015).
\newblock \emph{Graph Databases} (2nd ed.).
\newblock O'Reilly.
\newblock Praxisnahe Einführung in Property-Graph-Modelle, relevant
für Abschnitt~\ref{sec:ausblick-hybrid}.

\bibitem{Deutsch2022}
Deutsch, A.; Francis, N.; Green, A.; Hare, K.; Li, B.; Libkin, L.;
Lindaaker, T.; Marsault, V.; Martens, W.; Michels, J.; et al. (2022).
\newblock Graph pattern matching in GQL and SQL/PGQ.
\newblock \emph{Proceedings of ACM SIGMOD}, 2246--2258.
\newblock Beschreibung von SQL/PGQ (SQL:2023) und GQL, zentrale
Referenz für die hybride Anfragesprachendiskussion in
Abschnitt~\ref{sec:ausblick-hybrid}.

\bibitem{Pgvector2023}
Kane, A. (2023).
\newblock pgvector: Open-source vector similarity search for Postgres.
\newblock \emph{GitHub-Repository}, \url{https://github.com/pgvector/pgvector}.
\newblock Erweiterung für Vektorindizes in PostgreSQL, diskutiert im
Kontext semantischer Ähnlichkeitskanten in
Abschnitt~\ref{sec:ausblick-ml}.

\bibitem{Roddick1995}
Roddick, J.\,F. (1995).
\newblock A survey of schema versioning issues for database systems.
\newblock \emph{Information and Software Technology}, 37(7), 383--393.
\newblock Ãœbersicht zu Schemaevolution, Grundlage von
Theorem~\ref{thm:schaemaevolution}.

\bibitem{Curino2008}
Curino, C.; Moon, H.\,J.; Zaniolo, C. (2008).
\newblock Graceful database schema evolution: The PRISM workbench.
\newblock \emph{Proceedings of VLDB}, 761--772.
\newblock Werkzeugunterstützung für Schemaevolution, relevant für den
in Abschnitt~\ref{sec:ausblick-tooling} skizzierten Schema-Linter.

\bibitem{Fowler2018}
Fowler, M. (2018).
\newblock \emph{Refactoring: Improving the Design of Existing Code}
(2nd ed.).
\newblock Addison-Wesley.
\newblock Allgemeine Refactoring-Prinzipien, übertragen auf
Datenbankschemata in Abschnitt~\ref{sec:ausblick-tooling}.

\bibitem{Melton2003}
Melton, J.; Simon, A.\,R. (2003).
\newblock \emph{SQL:1999 -- Understanding Relational Language Components}.
\newblock Morgan Kaufmann.
\newblock Detaillierte Darstellung des SQL:1999-Standards
(Abschnitt~\ref{sec:sql-standards}), insbesondere rekursiver Anfragen
und objektrelationaler Erweiterungen.

\bibitem{Gulutzan1999}
Gulutzan, P.; Pelzer, T. (1999).
\newblock \emph{SQL-99 Complete, Really}.
\newblock CMP Books.
\newblock Praxisorientierte Referenz zu SQL:1999-Konstrukten
einschließlich \texttt{CASE}, \texttt{CHECK} und
Trigger-Syntax (Kapitel~\ref{ch:plsql}).

\bibitem{Obe2017}
Obe, R.; Hsu, L.\,S. (2017).
\newblock \emph{PostgreSQL: Up and Running} (3rd ed.).
\newblock O'Reilly.
\newblock Praktische Einführung in PostgreSQL-spezifische
Erweiterungen wie \texttt{DISTINCT ON}
(Abschnitt~\ref{sec:distinct-on}) und \texttt{LATERAL}-Joins.

\bibitem{PostgreSQLDoc2024}
The PostgreSQL Global Development Group (2024).
\newblock \emph{PostgreSQL 16 Documentation, Chapter 7: Queries}.
\newblock \url{https://www.postgresql.org/docs/16/queries.html}.
\newblock Offizielle Dokumentation der \texttt{SELECT
DISTINCT ON}-Syntax (Definition~\ref{def:distinct-on}) und der
Bereichstypen für temporale Anfragen
(Abschnitt~\ref{sec:ausblick-temporal}).

% ============================================================
% Programmiersprachenanbindung
% ============================================================

\bibitem{Descartes2014}
Descartes, A.; Bunce, T. (2014).
\newblock \emph{Programming the Perl DBI}.
\newblock O'Reilly.
\newblock Standardreferenz zur DBI-Schnittstelle
(Abschnitt~\ref{sec:perl-dbi}).

\bibitem{ISO9075-2016}
ISO/IEC 9075-2:2016.
\newblock \emph{Information technology -- Database languages -- SQL --
Part 2: Foundation (SQL/Foundation)}.
\newblock International Organization for Standardization.
\newblock Normativer Standard für Embedded SQL
(Abschnitt~\ref{sec:embedded-sql}) und Cursor-Semantik
(Bemerkung~\ref{rem:cursor-iterator}).

\bibitem{Bauer2015}
Bauer, C.; King, G.; Gregory, G. (2015).
\newblock \emph{Java Persistence with Hibernate} (2nd ed.).
\newblock Manning.
\newblock Standardwerk zu objektrelationalem Mapping,
Grundlage von Abschnitt~\ref{sec:orm} einschließlich Lazy Loading
(Proposition~\ref{prop:lazy-loading}) und des
$N{+}1$-Problems (Bemerkung~\ref{rem:n-plus-1}).

% ============================================================
% Empirische Software-Engineering-Forschung
% ============================================================

\bibitem{Sjoberg2005}
Sjøberg, D.\,I.\,K.; Hannay, J.\,E.; Hansen, O.; Kampenes, V.\,B.;
Karahasanovic, A.; Liborg, N.-K.; Rekdal, A.\,C. (2005).
\newblock A survey of controlled experiments in software engineering.
\newblock \emph{IEEE Transactions on Software Engineering}, 31(9), 733--753.
\newblock Methodische Grundlage für die in
Abschnitt~\ref{sec:ausblick-empirie} vorgeschlagene empirische Studie
zu Schema-Topologie und Wartungsaufwand.

\bibitem{Kitchenham2004}
Kitchenham, B.\,A.; Pfleeger, S.\,L.; Pickard, L.\,M.; Jones, P.\,W.;
Hoaglin, D.\,C.; El Emam, K.; Rosenberg, J. (2002).
\newblock Preliminary guidelines for empirical research in software
engineering.
\newblock \emph{IEEE Transactions on Software Engineering}, 28(8), 721--734.
\newblock Leitlinien für die methodische Durchführung der in
Abschnitt~\ref{sec:ausblick-empirie} skizzierten Studie.

% ============================================================
% Ergänzende Literatur: Geschichte und Grundlagen der Datenbanktechnik
% ============================================================

\end{thebibliography}


\end{document}

